% Generated mechanically from frozen input p12/ja/build/ch115/obsolete_full.tex.
% Do not edit this generated file; edit the bound locale source instead.

% OK, start here.
%

\setcounter{chapter}{114}
\chapter{旧結果}


\phantomsection
\label{obsolete-section-phantom}


% Frozen obsolete.tex lines 19--73.
\section{序論}
\label{obsolete-section-introduction}

\noindent
本章には，現在では「旧結果」となったいくつかの補題を収録する
（\cite{Miller} を参照せよ）．


\section{予備事項}
\label{obsolete-section-preliminaries}

\begin{remark}
\label{obsolete-remark-composition-of-adjoints-isomorphic-to-identity}
この注に以前含まれていた内容は，現在では Categories の Lemma
\ref{categories-lemma-adjoint-fully-faithful} と Lemma
\ref{categories-lemma-left-adjoint-composed-fully-faithful} を
組み合わせた結果に包含されている．
\end{remark}


\section{ホモロジー代数}
\label{obsolete-section-homological-algebra}

\begin{remark}
\label{obsolete-remark-weak-serre-subcategory}
以下の内容は Homology の Lemma
\ref{homology-lemma-biregular-ss-converges} と Lemma
\ref{homology-lemma-first-quadrant-ss} に包含されるため，
現在では旧結果である．
$\mathcal{A}$ をアーベル圏とする．
$\mathcal{C} \subset \mathcal{A}$ を弱 Serre 部分圏とする
（Homology の Definition
\ref{homology-definition-serre-subcategory} を参照せよ）．
$K^{\bullet, \bullet}$ は Homology の Lemma
\ref{homology-lemma-first-quadrant-ss} を適用できる二重複体であり，
ある $r \geq 0$ に対して，すべての対象
${}'E_r^{p, q}$ が $\mathcal{C}$ に属すると仮定する．このとき，
すべてのコホモロジー群 $H^n(sK^\bullet)$ は $\mathcal{C}$ に属する．
実際，仮定から ${}'d_r^{p, q}$ の核と像は $\mathcal{C}$ に属する．
したがって，各 ${}'E_{r + 1}^{p, q}$ が $\mathcal{C}$ に属することが分かる．
帰納法により，各 ${}'E_\infty^{p, q}$ が $\mathcal{C}$ に属する．
ゆえに各 $H^n(sK^\bullet)$ は，その各部分商が $\mathcal{C}$ に属する
有限フィルトレーションをもつ．$\mathcal{C}$ が拡大について閉じていることから，
主張どおり $H^n(sK^\bullet)$ が $\mathcal{C}$ に属すると結論できる．
$K^{\bullet, \bullet}$ に付随する第2スペクトル系列についても同じ結果が成り立つ．
同様に，$(K^\bullet, F)$ が Homology の Lemma
\ref{homology-lemma-biregular-ss-converges} を適用できるフィルター付き複体であり，
ある $r \geq 0$ に対してすべての対象 $E_r^{p, q}$ が
$\mathcal{C}$ に属するならば，各 $H^n(K^\bullet)$ は
$\mathcal{C}$ の対象である．
\end{remark}
% Frozen obsolete.tex lines 74--253.
\section{旧結果となった代数の補題}
\label{obsolete-section-algebra}

\begin{lemma}
\label{obsolete-lemma-finite-presentation-module-independent}
$M$ を有限表示 $R$-加群とする．任意の全射
$\alpha : R^{\oplus n} \to M$ に対し，$\alpha$ の核は有限 $R$-加群である．
\end{lemma}

\begin{proof}
これは Algebra の Lemma \ref{algebra-lemma-extension} の特別な場合である．
\end{proof}

\begin{lemma}
\label{obsolete-lemma-p-ring-map}
$\varphi : R \to S$ を環準同型とする．次を仮定する：
\begin{enumerate}
\item 任意の $x \in S$ に対して，$x^n$ が $\varphi$ の像に属するような
$n > 0$ が存在する；
\item 任意の $x \in \Ker(\varphi)$ に対して，$x^n = 0$ となる
$n > 0$ が存在する．
\end{enumerate}
このとき $\varphi$ はスペクトル上の同相写像を誘導する．さらに，素数 $p$ があり，
\begin{enumerate}
\item[(a)] $S$ は $R$-代数として，ある $n > 0$ に対して
$x^{p^n} \in \varphi(R)$ かつ $p^nx \in \varphi(R)$ を満たす元 $x$ たちで生成され，
\item[(b)] $\varphi$ の核は冪零元たちで生成される
\end{enumerate}
とする．このとき (1) と (2) が成り立つ．また，任意の環準同型
$R \to R'$ に対して，環準同型 $R' \to R' \otimes_R S$ も
(a)，(b)，(1)，(2) を満たし，とくにスペクトル上の同相写像を誘導する．
\end{lemma}

\begin{proof}
これは Algebra の Lemma \ref{algebra-lemma-powers} と Lemma
\ref{algebra-lemma-p-ring-map} を組み合わせたものである．
\end{proof}

\noindent
次の技術的補題は，本質的有限型の環準同型について，ファイバー上の任意の関係式の列を，
適切な意味で全空間へ持ち上げられることを述べる．

\begin{lemma}
\label{obsolete-lemma-lift-elements-ideal}
$R \to S$ を環準同型とする．
$\mathfrak p \subset R$ を素イデアルとする．
$\mathfrak q \subset S$ を $\mathfrak p$ の上にある素イデアルとする．
$S_{\mathfrak q}$ は $R_\mathfrak p$ 上本質的有限型であると仮定する．
次が与えられているとする：
\begin{enumerate}
\item 整数 $n \geq 0$；
\item 素イデアル $\mathfrak a \subset \kappa(\mathfrak p)[x_1, \ldots, x_n]$；
\item 全射な $\kappa(\mathfrak p)$-準同型
$$
\psi : (\kappa(\mathfrak p)[x_1, \ldots, x_n])_{\mathfrak a}
\longrightarrow
S_{\mathfrak q}/\mathfrak p S_{\mathfrak q},
$$
および
\item $\Ker(\psi)$ の元 $\overline{f}_1, \ldots, \overline{f}_e$．
\end{enumerate}
このとき，次が存在する：
\begin{enumerate}
\item 整数 $m \geq 0$；
\item 元 $g \in S$, $g \not\in \mathfrak q$；
\item 準同型
$$
\Psi :
R[x_1, \ldots, x_n, x_{n + 1}, \ldots, x_{n + m}]
\longrightarrow
S_g,
$$
および
\item $\Ker(\Psi)$ の元
$f_1, \ldots, f_e, f_{e + 1}, \ldots, f_{e + m}$．
\end{enumerate}
これらは次を満たす：
\begin{enumerate}
\item 次の図式は可換である：
$$
\xymatrix{
R[x_1, \ldots, x_{n + m}] \ar[d]_\Psi
\ar[rr]_-{x_{n + j} \mapsto 0} & &
(\kappa(\mathfrak p)[x_1, \ldots, x_n])_{\mathfrak a} \ar[d]^\psi \\
S_g \ar[rr] & &
S_{\mathfrak q}/\mathfrak p S_{\mathfrak q}
},
$$
\item 元 $f_i$, $i \leq n$ は，局所環
$$
(\kappa(\mathfrak p)[x_1, \ldots, x_{n + m}])_{
(\mathfrak a, x_{n + 1}, \ldots, x_{n + m})},
$$
において $\overline{f}_i$ の単元倍へ写る；
\item 元 $f_{e + j}$ は同じ局所環において $x_{n + j}$ の単元倍へ写る；
\item 誘導される準同型 $R[x_1, \ldots, x_{n + m}]_{\mathfrak b}
\to S_{\mathfrak q}$ は全射である．ただし
$\mathfrak b = \Psi^{-1}(\mathfrak qS_g)$ とする．
\end{enumerate}
\end{lemma}

\begin{proof}
$R$ と $S$ がそれぞれ極大イデアル $\mathfrak p$ と $\mathfrak q$ をもつ
局所環である場合に補題を証明すれば十分であると主張する．実際，必要な性質を
すべて満たす
$$
\Psi' : R_{\mathfrak p}[x_1, \ldots, x_{n + m}]
\longrightarrow
S_{\mathfrak q}
$$
および $f_1', \ldots, f_{e + m}' \in
R_{\mathfrak p}[x_1, \ldots, x_{n + m}]$ が構成されたとする．このとき，
各 $ff_k'$ がある元 $f_k \in R[x_1, \ldots, x_{n + m}]$ に由来するような
$f \in R$, $f \not \in \mathfrak p$ が存在する．さらに，適当な
$g \in S$, $g \not \in \mathfrak q$ に対して，元 $\Psi'(x_i)$ は
元 $y_i \in S_g$ の像である．$\Psi(x_i) = y_i$ という規則で定まる
$R$-代数準同型を $\Psi$ とする．$\Psi(f_i)$ は局所化
$S_{\mathfrak q}$ で零なので，必要なら $g$ を取り替えることにより
$\Psi(f_i) = 0$ と仮定できる．これで主張が証明された．

\medskip\noindent
したがって，$R$ と $S$ はそれぞれ極大イデアル $\mathfrak p$ と
$\mathfrak q$ をもつ局所環であると仮定してよい．
$y_i \bmod \mathfrak pS = \psi(x_i)$ となる
$y_1, \ldots, y_n \in S$ を取る．
$S$ がその局所化となるような $R$-部分代数を生成する元を
$y_{n + 1}, \ldots, y_{n + m} \in S$ とする．このような元は，
$S$ が $R$ 上本質的有限型であるという仮定から存在する．$\psi$ は全射なので，
ある $h_j \in \kappa(\mathfrak p)[x_1, \ldots, x_n]_{\mathfrak a}$ により
$y_{n + j} \bmod \mathfrak pS = \psi(h_j)$ と書ける．共通分母
$d \in \kappa(\mathfrak p)[x_1, \ldots, x_n]$,
$d \not \in \mathfrak a$ を用いて $h_j = g_j/d$,
$g_j \in \kappa(\mathfrak p)[x_1, \ldots, x_n]$ と書く．
$g_j$ と $d$ の持ち上げ $G_j, D \in R[x_1, \ldots, x_n]$ を選ぶ．
$y_{n + j}' = D(y_1, \ldots, y_n) y_{n + j} - G_j(y_1, \ldots, y_n)$
と置く．構成により $y_{n + j}' \in \mathfrak p S$ である．
$y_1, \ldots, y_n, y_n', \ldots, y_{n + m}'$ が，その局所化が $S$ となる
$S$ の $R$-部分代数を生成することは明らかである．
$$
\Psi : R[x_1, \ldots, x_{n + m}] \to S
$$
を，$i = 1, \ldots, n$ に対して $x_i$ を $y_i$ へ，
$j = 1, \ldots, m$ に対して $x_{n + j}$ を $y'_{n + j}$ へ送る準同型とする．
(1) と (4) は構成から明らかである．さらに，イデアル $\mathfrak b$ は
多項式環 $\kappa(\mathfrak p)[x_1, \ldots, x_{n + m}]$ のイデアル
$(\mathfrak a, x_{n + 1}, \ldots, x_{n + m})$ の上へ写る．

\medskip\noindent
$J = \Ker(\Psi)$ と記す．短完全列
$$
0 \to J_{\mathfrak b}
\to R[x_1, \ldots, x_{n + m}]_{\mathfrak b}
\to S_{\mathfrak q}
\to 0.
$$
がある．全射性は，上での
$y_1, \ldots, y_n, y_n', \ldots, y_{n + m}'$ の選び方から従う．
このことから
$$
J_{\mathfrak b}/ \mathfrak pJ_{\mathfrak b}
\to \kappa(\mathfrak p)[x_1, \ldots, x_{n + m}]_{
(\mathfrak a, x_{n + 1}, \ldots, x_{n + m})}
\to S_{\mathfrak q}/\mathfrak pS_{\mathfrak q}
\to 0
$$
は完全である．構成により，最後の準同型の下で $x_i$ は $\psi(x_i)$ へ写り，
$x_{n + j}$ は零へ写る．したがって，補題の (2) と (3) にあるような
$f_i$ を容易に選べる．
\end{proof}
% Frozen obsolete.tex lines 255--383.
\begin{lemma}
\label{obsolete-lemma-smooth-independent-presentation}
$R \to S$ を有限表示な環準同型とする．$R$ 上の $S$ のある表示
$\alpha$ に対して，素朴な余接複体 $\NL(\alpha)$ が，次数 $0$ に置かれた
有限射影 $S$-加群と擬同型であるならば，これは任意の表示について成り立つ．
\end{lemma}

\begin{proof}
Algebra の Lemma \ref{algebra-lemma-NL-homotopy} から直ちに従う．
\end{proof}

\begin{remark}[射影分解]
\label{obsolete-remark-projective-resolution}
$R$ を環とする．任意の集合 $S$ に対して，$S$ 上の自由 $R$-加群を
$F(S)$ と記す．このとき任意の左 $R$-加群は次の2段階分解をもつ：
$$
F(M \times M) \oplus F(R \times M) \to F(M) \to M \to 0.
$$
最初の準同型は次の規則で与えられる：
$$
[m_1, m_2] \oplus [r, m] \mapsto [m_1 + m_2] - [m_1] - [m_2] + [rm] - r[m].
$$
\end{remark}

\begin{lemma}
\label{obsolete-lemma-spec-localization-first}
$S$ を $A$ の積閉集合とする．標準環準同型 $A \to S^{-1}A$ が誘導する写像
$$
f: \Spec(S^{-1}A)\longrightarrow \Spec(A)
$$
はその像の上への同相写像であり，
$\Im(f) = \{ \mathfrak p \in \Spec(A) : \mathfrak p\cap S = \emptyset \}$
である．
\end{lemma}

\begin{proof}
これは Algebra の Lemma \ref{algebra-lemma-spec-localization} と重複している．
\end{proof}

\begin{lemma}
\label{obsolete-lemma-finite-type-flat-over-integral-algebra}
$A \to B$ を有限型かつ平坦な環準同型とし，$A$ は整域であるとする．
このとき $B$ は有限表示 $A$-代数である．
\end{lemma}

\begin{proof}
More on Flatness の Proposition
\ref{flat-proposition-flat-finite-type-finite-presentation-domain} の特別な場合である．
\end{proof}

\begin{lemma}
\label{obsolete-lemma-helper-finite-type-flat-finite-presentation}
$R$ を分数体 $K$ をもつ整域とする．
$S = R[x_1, \ldots, x_n]$ を $R$ 上の多項式環とする．
$M$ を有限 $S$-加群とし，$M$ は $R$ 上平坦であると仮定する．
すべての部分環 $R \subset R' \subset K$, $R \not = R'$ に対して，
加群 $M \otimes_R R'$ が $S \otimes_R R'$ 上有限表示ならば，
$M$ は $S$ 上有限表示である．
\end{lemma}

\begin{proof}
実は，補題の主張にあるすべての $R'$ に対して $M \otimes_R R'$ が
有限表示であるという仮定がなくても，$M$ は有限表示なので，この補題は正しい．
これは More on Flatness の Proposition
\ref{flat-proposition-flat-finite-type-finite-presentation-domain} から従う．
この補題には当初，誤った証明があり（間隙を見つけた Ofer Gabber に感謝する），
引用した命題の別証明に用いられていた．この補題を復活させるには，
$R$ が局所正規整域の場合に，可換代数の諸章の結果だけを用いる正しい議論が必要である．
そのような議論があれば
\href{mailto:stacks.project@gmail.com}{stacks.project@gmail.com}
まで電子メールを送ってほしい．
\end{proof}

\begin{lemma}
\label{obsolete-lemma-relative-effective-cartier-algebra}
$A \to B$ を環準同型とする．$f \in B$ とし，次を仮定する：
\begin{enumerate}
\item $A \to B$ は平坦である；
\item $f$ は非零因子である；
\item $A \to B/fB$ は平坦である．
\end{enumerate}
このとき，任意のイデアル $I \subset A$ に対して，写像
$f : B/IB \to B/IB$ は単射である．
\end{lemma}

\begin{proof}
$B$ と $B/fB$ が $A$ 上平坦であることにより，
$IB = I \otimes_A B$ および $I(B/fB) = I \otimes_A B/fB$ であることに注意する．
とくに $IB/fIB \cong I \otimes_A B/fB$ は $B/fB$ へ単射的に写る．
したがって，完全な行をもつ次の図式に蛇の補題を適用すれば結果が従う：
$$
\xymatrix{
0 \ar[r] &
I \otimes_A B \ar[r] \ar[d]^f &
B \ar[r] \ar[d]^f &
B/IB \ar[r] \ar[d]^f &
0 \\
0 \ar[r] &
I \otimes_A B \ar[r] &
B \ar[r] &
B/IB \ar[r] &
0
}
$$
\end{proof}

\begin{lemma}
\label{obsolete-lemma-faithfully-flat-injective}
$R \to S$ が忠実平坦な環準同型ならば，任意の $R$-加群 $M$ に対して，
写像 $M \to S \otimes_R M$, $x \mapsto 1 \otimes x$ は単射である．
\end{lemma}

\begin{proof}
この補題は Algebra の Lemma
\ref{algebra-lemma-faithfully-flat-universally-injective} と重複している．
\end{proof}

\begin{remark}
\label{obsolete-remark-section-colimits}
この参照・タグは以前，Smoothing Ring Maps のある Section を指していたが，
その内容はその後 Algebra の Section \ref{algebra-section-colimits-flat} に
包含された．
\end{remark}
% Frozen obsolete.tex lines 385--577.
\begin{lemma}
\label{obsolete-lemma-not-domain}
$(R, \mathfrak m)$ を次元 $1$ の被約 Noether 局所環とし，
$x \in \mathfrak m$ を非零因子とする．
$\mathfrak q_1, \ldots, \mathfrak q_r$ を $R$ の極小素イデアルとする．このとき
$$
\text{length}_R(R/(x)) = \sum\nolimits_i \text{ord}_{R/\mathfrak q_i}(x)
$$
である．
\end{lemma}

\begin{proof}
Chow Homology の Lemma
\ref{chow-lemma-additivity-divisors-restricted} の特別な（非常に容易な）場合である．
\end{proof}

\begin{lemma}
\label{obsolete-lemma-bound-primes}
$A$ を次元 $2$ の Noether 局所正規整域とする．非零な
$f \in \mathfrak m$ に対して，$A$ の穴あきスペクトル上の $f$ に付随する因子を
$\text{div}(f) = \sum n_i (\mathfrak p_i)$ と記す．
$|f| = \sum n_i$ と置く．このとき，すべての $g \in \mathfrak m^N$ に対して
$|f + g| \leq M$ となる整数 $N$ と $M$ が存在する．
\end{lemma}

\begin{proof}
$f, h$ が $A$ の正則列となるような $h \in \mathfrak m$ を取る
（これは Algebra の Lemma \ref{algebra-lemma-criterion-normal} と Lemma
\ref{algebra-lemma-depth-drops-by-one} から従う）．
$M = \text{length}_A(A/(f, h))$ とし，
$\mathfrak m^N \subset (f, h)$ となる任意の整数 $N$ を取れば補題が成り立つことを示す．
$\sqrt{(f, h)} = \mathfrak m$ なので，そのような整数 $N$ は存在する．
$(f, h) = (f + g, h)$ であるから，すべての $g \in \mathfrak m^N$ に対して
$M = \text{length}_A(A/(f + g, h))$ であることに注意する．さらに，これから
$f + g, h$ も $A$ の正則列であることが従う；Algebra の Lemma
\ref{algebra-lemma-reformulate-CM} を参照せよ．ここで
$\text{div}(f + g ) = \sum m_j (\mathfrak q_j)$ と仮定する．写像
$$
c : A/(f + g) \longrightarrow \prod A/\mathfrak q_j^{(m_j)}
$$
を考える．ここで $\mathfrak q_j^{(m_j)}$ は記号冪である；Algebra の Section
\ref{algebra-section-symbolic-power} を参照せよ．$A$ は正規なので，
$A_{\mathfrak q_i}$ は離散付値環であり，したがって
$$
A_{\mathfrak q_i}/(f + g) =
A_{\mathfrak q_i}/\mathfrak q_i^{m_i} A_{\mathfrak q_i} =
(A/\mathfrak q_i^{(m_i)})_{\mathfrak q_i}
$$
である．$V(f + g, h) = \{\mathfrak m\}$ なので，$h$ を可逆にすると $c$ は同型になる
（細部は省略する）．$h$ は $A/(f + g)$ 上の非零因子なので，
$A/(f + g, h)$ の長さは，Chow Homology の Section
\ref{chow-section-periodic-complexes} で定義された Herbrand 商
$e_A(A/(f + g), 0, h)$ に等しい．同様に，$A/(h, \mathfrak q_j^{(m_j)})$ の長さは
$e_A(A/\mathfrak q_j^{(m_j)}, 0, h)$ に等しい．そこで
\begin{align*}
M & = \text{length}_A(A/(f + g, h) \\
& =
e_A(A/(f + g), 0, h) \\
& =
\sum\nolimits_i e_A(A/\mathfrak q_j^{(m_j)}, 0, h) \\
& =
\sum\nolimits_i \sum\nolimits_{m = 0, \ldots, m_j - 1}
e_A(\mathfrak q_j^{(m)}/\mathfrak q_j^{(m + 1)}, 0, h)
\end{align*}
を得る．これらの等式は Chow Homology の Lemma
\ref{chow-lemma-additivity-periodic-length} と Lemma
\ref{chow-lemma-finite-periodic-length} から従う．ここではとくに，上で論じたように
$c$ の余核が有限長であることを用いた．Nakayama の補題により，
$e_A(\mathfrak q^{(m)}/\mathfrak q^{(m + 1)}, 0, h)$ が少なくとも $1$ であることは
直接証明できる．これで補題の証明が完了する．
\end{proof}

\begin{lemma}
\label{obsolete-lemma-flat-over-gorenstein-gorenstein-fibre}
$A \to B$ を Noether 局所環の平坦な局所準同型とする．
$A$ と $B/\mathfrak m_A B$ が Gorenstein ならば，$B$ も Gorenstein である．
\end{lemma}

\begin{proof}
Dualizing Complexes の Lemma
\ref{dualizing-lemma-flat-under-gorenstein} から直ちに従う．
\end{proof}

\begin{lemma}
\label{obsolete-lemma-kill-local}
$(A, \mathfrak m)$ を Noether 局所環とする．$I \subset A$ をイデアルとし，
$M$ を有限 $A$-加群とする．$s$ を整数とし，次を仮定する：
\begin{enumerate}
\item $A$ は双対化複体をもつ；
\item $\mathfrak p \not \in V(I)$ かつ
$V(\mathfrak p) \cap V(I) \not = \{\mathfrak m\}$ ならば，
$\text{depth}_{A_\mathfrak p}(M_\mathfrak p) + \dim(A/\mathfrak p) > s$．
\end{enumerate}
このとき，$V(J) \cap V(I) = \{\mathfrak m\}$ であり，
$i \leq s$ に対して $JI^n$ が $H^i_\mathfrak m(M)$ を零化するような
$n > 0$ とイデアル $J \subset A$ が存在する．
\end{lemma}

\begin{proof}
Local Cohomology の Lemma
\ref{local-cohomology-lemma-sitting-in-degrees} により，
$i \leq s$ に対して有限 $A$-加群
$E^i = \text{Ext}^{-i}_A(M, \omega_A^\bullet)$ についてこれを示せばよい．
$E^0 \oplus \ldots \oplus E^s$ の台 $Z$ は $\Spec(A)$ の閉集合であり，
(2) にあるような素イデアルを一つも含まない．したがって，補題の主張にあるような
ある $J$ とある指数に対して，その台は $V(JI^n)$ に含まれる．
\end{proof}

\begin{lemma}
\label{obsolete-lemma-algebraize-local-cohomology-bis-bis}
$(A, \mathfrak m)$ を Noether 局所環とする．$I \subset A$ をイデアルとし，
$M$ を有限 $A$-加群とする．$s$ と $d$ を整数とし，次を仮定する：
\begin{enumerate}
\item[(a)] $A$ は双対化複体をもつ；
\item[(b)] $\text{cd}(A, I) \leq d$；
\item[(c)] $\mathfrak p \not \in V(I)$ ならば，
$\text{depth}_{A_\mathfrak p}(M_\mathfrak p) > s$ または
$\text{depth}_{A_\mathfrak p}(M_\mathfrak p) + \dim(A/\mathfrak p) > d + s$．
\end{enumerate}
このとき，$A, I, \mathfrak m, M$ は Algebraic and Formal Geometry の Lemma
\ref{algebraization-lemma-algebraize-local-cohomology-bis} の仮定を満たし，
$i \leq s$ に対して
$H^i_\mathfrak m(M) \to \lim H^i_\mathfrak m(M/I^nM)$ は同型であり，
これらの加群は $I$ のある冪によって零化される．
\end{lemma}

\begin{proof}
Algebraic and Formal Geometry の Lemma
\ref{algebraization-lemma-algebraize-local-cohomology-bis} の仮定は，
より一般的な Algebraic and Formal Geometry の Lemma
\ref{algebraization-lemma-bootstrap-bis-bis} により満たされる．
次に Algebraic and Formal Geometry の Lemma
\ref{algebraization-lemma-algebraize-local-cohomology-bis} の結論が第2の主張を与える．
\end{proof}

\begin{lemma}
\label{obsolete-lemma-combine-one}
Algebraic and Formal Geometry の Situation
\ref{algebraization-situation-bootstrap} において，
$H^s_\mathfrak a(M) = \lim H^s_\mathfrak a(M/I^nM)$ である．
\end{lemma}

\begin{proof}
これは Algebraic and Formal Geometry の Theorem
\ref{algebraization-theorem-final-bootstrap} から直ちに従う．
この補題の元の版には追加の仮定があったが，この定理によって置き換えられた．
\end{proof}

\begin{lemma}
\label{obsolete-lemma-fully-faithful-simple-one}
$A$ を Noether 環とする．$f \in \mathfrak a$ を $A$ のあるイデアルの元とする．
$U = \Spec(A) \setminus V(\mathfrak a)$ と置き，次を仮定する：
\begin{enumerate}
\item $A$ は双対化複体をもち，$f$ に関して完備である；
\item $A_f$ は $(S_2)$ であり，すべての極小素イデアル
$\mathfrak p \subset A$ に対して，$f \not \in \mathfrak p$ かつ
$\mathfrak q \in V(\mathfrak p) \cap V(\mathfrak a)$ ならば
$\dim((A/\mathfrak p)_\mathfrak q) \geq 3$ である．
\end{enumerate}
このとき，完備化函手
$$
\textit{Coh}(\mathcal{O}_U)
\longrightarrow
\textit{Coh}(U, I\mathcal{O}_U),
\quad
\mathcal{F} \longmapsto \mathcal{F}^\wedge
$$
は，有限局所自由な対象からなる充満部分圏上で充満忠実である．
\end{lemma}

\begin{proof}
この補題は Algebraic and Formal Geometry の Lemma
\ref{algebraization-lemma-fully-faithful-simple-two} の特別な場合である．
\end{proof}
% Frozen obsolete.tex lines 578--819.
\section{Zariski の主定理に関する補題}
\label{obsolete-section-ZMT}

\noindent
この節の補題は，もともと Zariski の主定理の（代数的な版の）証明，すなわち
Algebra の Theorem \ref{algebra-theorem-main-theorem} で用いられていた．

\begin{lemma}
\label{obsolete-lemma-change-equation-multiply}
$R$ を環とし，$\varphi : R[x] \to S$ を環準同型とする．$t \in S$ とする．
$t$ が $R[x]$ 上整ならば，任意の $a \in R$ に対して
$\varphi(a)^\ell t$ が，$y \mapsto \varphi(ax)$ および
$r\in R$ に対する $r \mapsto \varphi(r)$ により定義される
$\varphi_a : R[y] \to S$ 上整となるような $\ell \geq 0$ が存在する．
\end{lemma}

\begin{proof}
$f_i \in R[x]$ により
$t^d + \sum_{i < d} \varphi(f_i)t^i = 0$ と書く．$\ell$ を，すべての
$f_i$ の $x$ に関する次数の最大値とする．等式に $\varphi(a)^\ell$ を掛けて
$\varphi(a)^\ell t^d + \sum_{i < d} \varphi(a^\ell f_i)t^i = 0$ を得る．
各 $\varphi(a^\ell f_i)$ は $\varphi_a$ の像に属することに注意する．
結果は Algebra の Lemma \ref{algebra-lemma-make-integral-trivial} から従う．
\end{proof}

\begin{lemma}
\label{obsolete-lemma-make-integral-less-trivial}
$\varphi : R \to S$ を環準同型とする．$t \in S$ が関係式
$\varphi(a_0) + \varphi(a_1)t + \ldots + \varphi(a_n) t^n = 0$ を満たすとする．
$u_n = \varphi(a_n)$，$u_{n-1} = u_n t + \varphi(a_{n-1})$ と置き，
同様にして $u_1 = u_2 t + \varphi(a_1)$ まで定める．このとき
$u_n, u_{n-1}, \ldots, u_1$ と
$u_nt, u_{n-1}t, \ldots, u_1t$ はすべて $R$ 上整であり，
$S$ のイデアル $(\varphi(a_0), \ldots, \varphi(a_n))$ と
$(u_n, \ldots, u_1)$ は等しい．
\end{lemma}

\begin{proof}
$n$ に関する帰納法で証明する．$u_n = \varphi(a_n)$ なので，Algebra の Lemma
\ref{algebra-lemma-make-integral-trivial} から $u_nt$ は $R$ 上整である．もちろん
$u_n = \varphi(a_n)$ も $R$ 上整である．したがって
$u_{n - 1} = u_n t  + \varphi(a_{n - 1})$ は $R$ 上整であり
（Algebra の Lemma \ref{algebra-lemma-integral-closure-is-ring} を参照せよ），
$$
\varphi(a_0) + \varphi(a_1)t + \ldots + \varphi(a_{n - 1})t^{n - 1} +
u_{n - 1}t^{n - 1} = 0.
$$
$S'$ を $S$ における $R$ の整閉包とすると，準同型 $S' \to S$ と
上の等式に帰納法の仮定を適用することにより，
$u_{n-1}, \ldots, u_1$ および $u_{n-1}t, \ldots, u_1t$ もすべて
$S'$ に属することが分かる．イデアルに関する主張は，元の形と
$u_1t + \varphi(a_0) = 0$ から直ちに従う．
\end{proof}

\begin{lemma}
\label{obsolete-lemma-make-integral-not-in-ideal}
$\varphi : R \to S$ を環準同型とする．$t \in S$ が関係式
$\varphi(a_0) + \varphi(a_1)t + \ldots + \varphi(a_n) t^n = 0$ を満たすとする．
$J \subset S$ を，少なくとも一つの $i$ に対して
$\varphi(a_i) \not \in J$ となるイデアルとする．このとき，
$u$ と $ut$ がともに $R$ 上整であるような
$u \in S$, $u \not\in J$ が存在する．
\end{lemma}

\begin{proof}
Lemma \ref{obsolete-lemma-make-integral-less-trivial} により，元 $u_i$ の少なくとも一つは
$J$ に属さないので，直ちに従う．
\end{proof}

\noindent
次の二つの補題は，一つの（非退化な）方程式で切り出される
$\mathbf{P}^1_R$ の閉部分スキームを記述する方法を与える．

\begin{lemma}
\label{obsolete-lemma-P1}
$R$ を環とする．$F(X, Y) \in R[X, Y]$ を次数 $d$ の斉次式とする．
$R$ の任意の素イデアル $\mathfrak p$ に対して，$F$ の少なくとも一つの係数が
$\mathfrak p$ に属さないと仮定する．次数付き環として
$S = R[X, Y]/(F)$ と置く．このとき，すべての $n \geq d$ に対して
$R$-加群 $S_n$ は階数 $d$ の有限局所自由加群である．
\end{lemma}

\begin{proof}
$R$-加群 $S_n$ は表示
$$
R[X, Y]_{n-d} \to R[X, Y]_n \to S_n \to 0.
$$
をもつ．したがって Algebra の Lemma \ref{algebra-lemma-cokernel-flat} により，
すべての素イデアル $\mathfrak p$ に対して，$F$ による乗法が単射
$\kappa(\mathfrak p)[X, Y] \to \kappa(\mathfrak p)[X, Y]$ を誘導することを
示せば十分である．これは，$F$ が $\mathfrak p$ を法として零多項式へ写らない
という仮定から明らかである．階数に関する主張もここから明らかである．
\end{proof}

\begin{lemma}
\label{obsolete-lemma-rel-prime-pols}
$k$ を体とする．$F, G \in k[X, Y]$ をそれぞれ次数 $d, e$ の斉次式とする．
$F, G$ は互いに素であると仮定する．このとき，$G$ による乗法は
$S = k[X, Y]/(F)$ 上で単射である．
\end{lemma}

\begin{proof}
これは「互いに素」の一つの定義である．別の定義を用いるならば，
それがこの定義と同値であることを示せる．
\end{proof}

\begin{lemma}
\label{obsolete-lemma-P1-localize}
$R$ を環とする．$F(X, Y) \in R[X, Y]$ を次数 $d$ の斉次式とする．
次数付き環として $S = R[X, Y]/(F)$ と置く．
$\mathfrak p \subset R$ を素イデアルとし，$F$ のある係数が
$\mathfrak p$ に属さないとする．
このとき，$f \in R$, $f \not\in \mathfrak p$，整数 $e$，および
$G \in R[X, Y]_e$ であって，すべての $n \geq d$ に対して $G$ による乗法が
同型 $(S_n)_f \to (S_{n + e})_f$ を誘導するものが存在する．
\end{lemma}

\begin{proof}
証明の途中で（有限回），$f\in R$, $f\not\in \mathfrak p$ に対して
$R$ を $R_f$ に取り替えてよい．まずそのような取り替えを行い，
$F$ のある係数が $R$ で可逆であるとする．とくに Lemma \ref{obsolete-lemma-P1} により，
$n \geq d$ に対して加群 $S_n$ は階数 $d$ の局所自由加群である．
$\kappa(\mathfrak p)[X, Y]$ における $G$ の像が $F(X, Y)$ の像と互いに素になるような
任意の $G \in R[X, Y]_e$ を取る（ある $e$ に対してこれは可能である）．
$G$ による乗法が誘導する写像 $S_d \to S_{d + e}$ に Algebra の Lemma
\ref{algebra-lemma-cokernel-flat} を適用する．$G$ の選び方と Lemma
\ref{obsolete-lemma-rel-prime-pols} により，
$S_d \otimes \kappa(\mathfrak p) \to S_{d + e} \otimes \kappa(\mathfrak p)$
は全単射である．したがって，適当な $f$ に対して $R$ を $R_f$ に取り替えれば，
$G : S_d \to S_{d + e}$ は全単射であると仮定できる．さらにこれは，
$R$ のすべての素イデアル $\mathfrak p'$ に対して，
$\kappa(\mathfrak p')[X, Y]$ における $G$ の像が $F$ の像と互いに素であることを
意味する．そこで Algebra の Lemma \ref{algebra-lemma-cokernel-flat} を再び用いると，
すべての写像 $G : S_d \to S_{d + e}$, $n \geq d$ は同型であることが分かる．
\end{proof}

\begin{remark}
\label{obsolete-remark-algebra}
$R$ を環とする．$F \in R[X, Y]_d$ と $G \in R[X, Y]_e$ があり，
$S = R[X, Y]/(F)$ と置いたとき，(1) すべての $n \geq d$ に対して
$S_n$ は階数 $d$ の有限局所自由加群であり，(2) $G$ による乗法が
すべての $n \geq d$ に対して同型 $S_n \to S_{n + e}$ を定めると仮定する．
この場合，有限局所自由 $R$-代数 $A$ を次のように定義できる：
\begin{enumerate}
\item $R$-加群として $A = S_{ed}$ とする；
\item 乗法 $A \times A \to A$ は，$S_{2ed}$ において
$G^d H_3 = H_1 H_2$ であるとき，かつそのときに限り $H_1 H_2 = H_3$ とする
規則によって与える．
\end{enumerate}
$G^d$ による乗法が全単射 $S_{de} \to S_{2de}$ を誘導するので，
これは意味をもつ．これが環構造を定めることは容易に分かる．
元 $G^d$ が環 $A$ の単位元を定めるという紛らわしい事実に注意する．
\end{remark}

\begin{lemma}
\label{obsolete-lemma-finite-after-localization}
$R$ を環とし，$f \in R$ とする．環 $S$, $S'$ と，次の可換図式をなす実線の矢印が
与えられているとする：
$$
\xymatrix{
& S'' \ar@{-->}[rd] \ar@{-->}[dd] &
\\
R \ar[rr] \ar@{-->}[ru] \ar[d] &  & S \ar[d]
\\
R_f \ar[r] & S' \ar[r] & S_f
}
$$
$R_f \to S'$ は有限であると仮定する．このとき，有限な環準同型
$R \to S''$ と図式中の点線の矢印を，$S' = (S'')_f$ となるように選べる．
\end{lemma}

\begin{proof}
$S'$ が $R_f$ 上 $x_i$, $i = 1, \ldots, w$ で生成されるとする．
$P_i(x_i) = 0$ となるモニック多項式 $P_i(t) \in R_f[t]$ を取る．
$P_i$ の次数を $d_i > 0$ とする．一様な $n$ に対して
$P_i(t) = t^{d_i} + \sum_{j < d_i} (a_{ij}/f^n) t^j$ と書く．
また，$S_f$ における $x_i$ の像を，適当な $g_i \in S$ により
$g_i / f^n$ と書く．このとき $S$ の元
$\xi_i = f^{nd_i} g_i^{d_i} + \sum_{j < d_i} f^{n(d_i - j)} a_{ij} g_i^j$
は $f$ のある冪によって零化される．したがって $n$ を $n'$ へ増やし，
$g_i$ を $f^{n' - n}g_i$ で取り替えることにより，$\xi_i = 0$ と仮定できる．
すると $S'$ は元 $f^n x_i$ たちで生成され，各元は $R$ に係数をもつモニック多項式
$Q_i(t) = t^{d_i} + \sum_{j < d_i} f^{n(d_i - j)} a_{ij} t^j$ の零点である．
また構成により，$S$ において $Q_i(f^ng_i) = 0$ である．したがって，上のような
可換図式に収まる有限 $R$-代数
$S'' = R[z_1, \ldots, z_w]/(Q_1(z_1), \ldots, Q_w(z_w))$ を得る．
準同型 $\alpha : S'' \to S$ は $z_i$ を $f^ng_i$ へ送り，
準同型 $\beta : S'' \to S'$ は $z_i$ を $f^nx_i$ へ送る．
この時点ではまだ，$\beta$ が同型 $(S'')_f \cong S'$ を誘導するとは限らない．
今のところ，この写像が全射であることだけが分かっている．問題は，
$S'$ では零へ写るが自身は零でない元 $h/f^n \in (S'')_f$ があり得ることである．
この場合，$\beta(h)$ は $S$ の元であり，ある $N$ に対して
$f^N \beta(h) = 0$ である．したがって $f^N h$ は $S''$ のイデアル
$J = \{h \in S'' \mid \alpha(h) = 0 \text{ and }
\beta(h) = 0\}$ の元である．すると $S''/J$ が求めるものになることは容易に分かる．
\end{proof}
% Frozen obsolete.tex lines 820--945.
\section{形式的に滑らかな環準同型}
\label{obsolete-section-formally-smooth}

\begin{lemma}
\label{obsolete-lemma-formally-smooth-smooth}
$R$ を環とする．$S$ を $R$-代数とする．$S$ が有限表示であり，
$R$ 上形式的に滑らかならば，$S$ は $R$ 上滑らかである．
\end{lemma}

\begin{proof}
Algebra の Proposition \ref{algebra-proposition-smooth-formally-smooth} を参照せよ．
\end{proof}

\begin{remark}
\label{obsolete-remark-equation-derivatives}
このタグは以前，Algebraization of Formal Spaces の Proposition
\ref{restricted-proposition-approximate} の証明中のある等式を指していたが，
内容の再配置によって使われなくなった．
\end{remark}

\begin{remark}
\label{obsolete-remark-equation-ci}
このタグは以前，Algebraization of Formal Spaces の Proposition
\ref{restricted-proposition-approximate} の証明中のある等式を指していたが，
内容の再配置によって使われなくなった．
\end{remark}

\begin{remark}
\label{obsolete-remark-equation-in-ideal}
このタグは以前，Algebraization of Formal Spaces の Proposition
\ref{restricted-proposition-approximate} の証明中のある等式を指していたが，
内容の再配置によって使われなくなった．
\end{remark}

\begin{remark}
\label{obsolete-remark-equation-derivatives-analogue}
このタグは以前，Algebraization of Formal Spaces の Proposition
\ref{restricted-proposition-approximate} の証明中のある等式を指していたが，
内容の再配置によって使われなくなった．
\end{remark}

\begin{remark}
\label{obsolete-remark-equation-go-down}
このタグは以前，Algebraization of Formal Spaces の Lemma
\ref{restricted-lemma-lift-approximation} の証明中のある等式を指していたが，
内容の再配置によって使われなくなった．
\end{remark}

\begin{lemma}
\label{obsolete-lemma-get-morphism-general}
$A$ を Noether 環とする．$I \subset A$ をイデアルとする．
$t$ を $I$ の生成元の最小数とする．
$C$ を Noether な $I$-進完備 $A$-代数とする．
$I \subset A \to C$ のみに依存し，次の性質をもつ整数 $d \geq 0$ が存在する．
次が与えられているとする：
\begin{enumerate}
\item $c \geq 0$ と，Algebraization of Formal Spaces の Equation
(\ref{restricted-equation-C-prime}) にある $B$ であって，$a \in I^c$ に対して
$a$ による乗法が $D(B)$ において $\NL_{B/A}^\wedge$ 上零であるもの；
\item 整数 $n > 2t\max(c, d)$；
\item $A/I^n$-代数準同型 $\psi_n : B/I^nB \to C/I^nC$．
\end{enumerate}
このとき，$m = \lfloor \frac{n}{t} \rfloor$ と置けば，
$\psi_n \bmod I^{m - c} = \varphi \bmod I^{m - c}$ を満たす
$A$-代数準同型 $\varphi : B \to C$ が存在する．
\end{lemma}

\begin{proof}
この補題は，より強い Algebraization of Formal Spaces の Lemma
\ref{restricted-lemma-get-morphism-general-better} によって旧結果となった．
実際，この補題をそこから導く．

\medskip\noindent
上の主張にあるような $I \subset A \to C$ が与えられたとする．
Local Cohomology の Section \ref{local-cohomology-section-uniform} で得られる整数を
$d(\text{Gr}_I(C))$ および $q(\text{Gr}_I(C))$ と記す．
$t$ は $IC$ の生成元の最小数の上界なので，
$d(\text{Gr}_I(C)) + 1 \leq t$ であることに注意する；Local Cohomology の Section
\ref{local-cohomology-section-uniform} の議論を参照せよ．そうでなければ補題は
何も述べないので，$t \geq 1$ と仮定してよい．
$$
d = q(\text{Gr}_I(C))
$$
とすれば補題が成り立つと主張する．実際，$c$, $B$, $n$, $\psi_n$ が
上の主張にあるようなものとする．このとき
$$
n > 2t\max(c, d) \Rightarrow n \geq 2tc + 1 \Rightarrow
n \geq 2(d(\text{Gr}_I(C)) + 1)c + 1
$$
である．他方，
$$
n > 2t\max(c, d) \Rightarrow n > t(c + d) \Rightarrow
n \geq q(C) + tc \geq q(\text{Gr}_I(C)) + (d(\text{Gr}_I(C)) + 1)c
$$
である．したがって Algebraization of Formal Spaces の Lemma
\ref{restricted-lemma-get-morphism-general-better} の仮定が満たされ，
$I^{n - (d(\text{Gr}_I(C)) + 1)c}C$ を法として $\psi_n$ と合同な
$A$-代数準同型 $\varphi : B \to C$ を得る．
さらに
\begin{align*}
n - (d(\text{Gr}_I(C)) + 1)c
& = \frac{n}{t} + \frac{(t - 1)n}{t} - (d(\text{Gr}_I(C)) + 1)c \\
& \geq \frac{n}{t} + \frac{(d(\text{Gr}_I(C))n}{t} - (d(\text{Gr}_I(C)) + 1)c \\
& > \frac{n}{t} + \frac{d(\text{Gr}_I(C))2tc}{t} - (d(\text{Gr}_I(C)) + 1)c \\
& = \frac{n}{t} + 2d(\text{Gr}_I(C))c - (d(\text{Gr}_I(C)) + 1)c \\
& = \frac{n}{t} + d(\text{Gr}_I(C))c - c \\
& \geq m - c
\end{align*}
なので，望むとおり $I^{m - c}C$ を法として $\varphi$ と $\psi_n$ は合同である．
\end{proof}
% Frozen obsolete.tex lines 946--999.
\section{サイトと層}
\label{obsolete-section-sites}

\begin{remark}[下付き感嘆符から順像への写像は存在しない]
\label{obsolete-remark-from-shriek-to-star}
$U$ をサイト $\mathcal{C}$ の対象とする．$\mathcal{C}/U$ 上の任意のアーベル層
$\mathcal{G}$ に対して，標準写像
$$
c : j_{U!}\mathcal{G} \longrightarrow j_{U*}\mathcal{G}
$$
が存在するかを問いたくなるかもしれない．このようなものを構成することは，写像
$j_U^{-1}j_{U!}\mathcal{G} \to \mathcal{G}$ を構成することと同じである．
制限は層化と可換であることに注意する．したがって Modules on Sites の Lemma
\ref{sites-modules-lemma-extension-by-zero} の前層を用いることができる．ゆえに，
$V/U$ に対して制限と両立する写像
$$
\bigoplus\nolimits_{\varphi \in \Mor_\mathcal{C}(V, U)}
\mathcal{G}(V \xrightarrow{\varphi} U)
\longrightarrow
\mathcal{G}(V/U)
$$
を定義すれば十分である．$\varphi$ が構造射
$\varphi_0 : V \to U$ である直和成分だけでは $1$ を取り，それ以外のすべての
直和成分上では零である写像を取れそうに見える．しかし，これは制限写像と両立しない．
実際，$\alpha : V' \to V$ が $\mathcal{C}$ の射ならば，
$\mathcal{C}/U$ の対象 $V'/U$ の構造射を
$\varphi'_0 = \varphi_0 \circ \alpha$ と記す．次の図式が可換であることを
確かめる必要がある：
$$
\xymatrix{
\bigoplus\nolimits_{\varphi \in \Mor_\mathcal{C}(V, U)}
\mathcal{G}(V \xrightarrow{\varphi} U)
\ar[d] \ar[r] &
\mathcal{G}(V/U) \ar[d] \\
\bigoplus\nolimits_{\varphi' \in \Mor_\mathcal{C}(V', U)}
\mathcal{G}(V' \xrightarrow{\varphi'} U)
\ar[r] &
\mathcal{G}(V'/U)
}
$$
ここで問題となるのは，$\varphi_0$ と異なる射 $\varphi : V \to U$ であって，
$\varphi \circ \alpha = \varphi'_0$ を満たすものが存在し得ることである．すると左の
垂直矢印は $\varphi$ に対応する直和成分を，下の水平矢印が $1$ に等しくなる直和成分へ
送るので，ほぼ確実に図式は可換でない．
\end{remark}
% Frozen obsolete.tex lines 1000--1224.
\section{コホモロジー}
\label{obsolete-section-cohomology}

\noindent
コホモロジーに関する旧結果となった補題を挙げる．

\begin{lemma}
\label{obsolete-lemma-ML-general}
$I$ を環 $A$ のイデアルとする．$X$ を $\Spec(A)$ 上のスキームとする．
$$
\ldots \to \mathcal{F}_3 \to \mathcal{F}_2 \to \mathcal{F}_1
$$
を，$\mathcal{F}_n = \mathcal{F}_{n + 1}/I^n\mathcal{F}_{n + 1}$ を満たす
$\mathcal{O}_X$-加群の逆系とする．
$$
\bigoplus\nolimits_{n \geq 0} H^1(X, I^n\mathcal{F}_{n + 1})
$$
が次数付き $\bigoplus_{n \geq 0} I^n/I^{n + 1}$-加群として昇鎖条件を満たすと
仮定する．このとき，逆系 $M_n = \Gamma(X, \mathcal{F}_n)$ は
Mittag--Leffler 条件を満たす．
\end{lemma}

\begin{proof}
これは，より一般的な Cohomology の Lemma
\ref{cohomology-lemma-ML-general} の特別な場合である．
\end{proof}

\begin{lemma}
\label{obsolete-lemma-ML-general-better}
$I$ を環 $A$ のイデアルとする．$X$ を $\Spec(A)$ 上のスキームとする．
$$
\ldots \to \mathcal{F}_3 \to \mathcal{F}_2 \to \mathcal{F}_1
$$
を，$\mathcal{F}_n = \mathcal{F}_{n + 1}/I^n\mathcal{F}_{n + 1}$ を満たす
$\mathcal{O}_X$-加群の逆系とする．$n$ が与えられたとき，
$$
H^1_n =
\bigcap\nolimits_{m \geq n}
\Im\left(
H^1(X, I^n\mathcal{F}_{m + 1}) \to H^1(X, I^n\mathcal{F}_{n + 1})
\right)
$$
と定める．$\bigoplus H^1_n$ が次数付き
$\bigoplus_{n \geq 0} I^n/I^{n + 1}$-加群として昇鎖条件を満たすならば，逆系
$M_n = \Gamma(X, \mathcal{F}_n)$ は Mittag--Leffler 条件を満たす．
\end{lemma}

\begin{proof}
これは，より一般的な Cohomology の Lemma
\ref{cohomology-lemma-ML-general-better} の特別な場合である．
\end{proof}

\begin{lemma}
\label{obsolete-lemma-topology-I-adic-general}
$I$ を環 $A$ の有限生成イデアルとする．$X$ を $\Spec(A)$ 上のスキームとする．
$$
\ldots \to \mathcal{F}_3 \to \mathcal{F}_2 \to \mathcal{F}_1
$$
を，$\mathcal{F}_n = \mathcal{F}_{n + 1}/I^n\mathcal{F}_{n + 1}$ を満たす
$\mathcal{O}_X$-加群の逆系とする．
$$
\bigoplus\nolimits_{n \geq 0} H^0(X, I^n\mathcal{F}_{n + 1})
$$
が次数付き $\bigoplus_{n \geq 0} I^n/I^{n + 1}$-加群として昇鎖条件を満たすと
仮定する．このとき，$M = \lim \Gamma(X, \mathcal{F}_n)$ 上の極限位相は
$I$-進位相である．
\end{lemma}

\begin{proof}
これは，より一般的な Cohomology の Lemma
\ref{cohomology-lemma-topology-I-adic-general} の特別な場合である．
\end{proof}

\begin{lemma}
\label{obsolete-lemma-pullback-K-flat}
$(\Sh(\mathcal{C}), \mathcal{O}_\mathcal{C})$ を環付きトポスとする．
任意の $\mathcal{O}_\mathcal{C}$-加群の複体 $\mathcal{G}^\bullet$ に対して，
擬同型 $\mathcal{K}^\bullet \to \mathcal{G}^\bullet$ が存在し，任意の環付きトポスの射
$f : (\Sh(\mathcal{D}), \mathcal{O}_\mathcal{D}) \to
(\Sh(\mathcal{C}), \mathcal{O}_\mathcal{C})$ に対して，
$f^*\mathcal{K}^\bullet$ は $\mathcal{O}_\mathcal{D}$-加群の K-flat 複体である．
\end{lemma}

\begin{proof}
これは Cohomology on Sites の Lemma
\ref{sites-cohomology-lemma-K-flat-resolution} と Lemma
\ref{sites-cohomology-lemma-pullback-K-flat} から従う．
\end{proof}

\begin{remark}
\label{obsolete-remark-pullback-K-flat}
この注では以前，K-flat 複体の引き戻しが K-flat であるか否かについて知られている
ことを論じていたが，現在では Cohomology on Sites の Lemma
\ref{sites-cohomology-lemma-pullback-K-flat} によって旧結果となっている．
\end{remark}

\noindent
次の補題は，特別な場合に導来極限のコホモロジー層を計算する．

\begin{lemma}
\label{obsolete-lemma-Rlim-of-system}
$(\mathcal{C}, \mathcal{O})$ を環付きサイトとする．$(K_n)$ を
$D(\mathcal{O})$ の対象の逆系とする．
$\mathcal{B} \subset \Ob(\mathcal{C})$ を部分集合とし，
$d \in \mathbf{N}$ とする．次を仮定する：
\begin{enumerate}
\item すべての $n$ に対して $K_n$ は $D^+(\mathcal{O})$ の対象である；
\item $q \in \mathbf{Z}$ に対して $n(q)$ が存在し，$n \geq n(q)$ ならば
$H^q(K_{n + 1}) \to H^q(K_n)$ は同型である；
\item $\mathcal{C}$ のすべての対象は，その各要素が $\mathcal{B}$ に属する被覆をもつ；
\item すべての $U \in \mathcal{B}$ に対し，$p > d$ およびすべての $q$ について
$H^p(U, H^q(K_n)) = 0$ である．
\end{enumerate}
このとき，すべての $m \in \mathbf{Z}$ に対して
$H^m(R\lim K_n) = \lim H^m(K_n)$ である．
\end{lemma}

\begin{proof}
$K = R\lim K_n$ と置く．$U \in \mathcal{B}$ とする．各 $n$ に対してスペクトル系列
$$
H^p(U, H^q(K_n)) \Rightarrow H^{p + q}(U, K_n)
$$
がある．$K_n$ は下に有界なので，これは収束する；Derived Categories の Lemma
\ref{derived-lemma-two-ss-complex-functor} を参照せよ．
$m \in \mathbf{Z}$ を固定すると，仮定 (4) により，$H^m(U, K_n)$ に寄与するのは
$0 \leq p \leq d$ かつ $m - d \leq q \leq m$ である
$H^p(U, H^q(K_n))$ だけである．仮定 (2) により，これは
$n \geq \max{n(m), \ldots, n(m - d)}$ となるや否や
$H^m(U, K_{n + 1}) \to H^m(U, K_n)$ が同型になることを意味する．函手
$R\Gamma(U, -)$ は Injectives の Lemma
\ref{injectives-lemma-RF-commutes-with-Rlim} により導来極限と可換する．したがって
$$
H^m(U, K) = H^m(R\lim R\Gamma(U, K_n))
$$
である．他方，上で複体 $R\Gamma(U, K_n)$ のコホモロジー群は最終的に一定であると
分かった．したがって More on Algebra の Remark
\ref{more-algebra-remark-compare-derived-limit} により，$U \in \mathcal{B}$ に依存しない
ある上界について，すべての $n \gg 0$ に対し $H^m(U, K)$ は
$H^m(U, K_n)$ に等しい．そのような $n$ を取る．最後に，$H^m(K)$ は前層
$U \mapsto H^m(U, K)$ の層化であり，$H^m(K_n)$ は前層
$U \mapsto H^m(U, K_n)$ の層化であることを思い出す．これらの前層は
$\mathcal{B}$ の要素上で同じ値をもつ．したがって仮定 (3) は，その層化も同じである
ことを保証する．補題が従う．
\end{proof}

\begin{lemma}
\label{obsolete-lemma-trivialities-cohomological-descent-abelian}
Simplicial Spaces の Situation
\ref{spaces-simplicial-situation-simplicial-site} において，
Simplicial Spaces の Remark \ref{spaces-simplicial-remark-augmentation-site} にあるような
サイト $\mathcal{D}$ への増大写像を $a_0$ とする．厳密充満な弱 Serre 部分圏
$$
\mathcal{A} \subset \textit{Ab}(\mathcal{D}),\quad
\mathcal{A}_n \subset \textit{Ab}(\mathcal{C}_n)
$$
が与えられているとする．このとき
\begin{enumerate}
\item[(1)] すべての $n$ に対して $\mathcal{C}_n$ への制限が
$\mathcal{A}_n$ に属するような $\mathcal{C}_{total}$ 上のアーベル層
$\mathcal{F}$ の集まりは，厳密充満な弱 Serre 部分圏
$\mathcal{A}_{total} \subset \textit{Ab}(\mathcal{C}_{total})$ である．
\end{enumerate}
すべての $n$ に対して $a_n^{-1}$ が $\mathcal{A}$ を $\mathcal{A}_n$ へ送るならば，
\begin{enumerate}
\item[(2)] $a^{-1}$ は $\mathcal{A}$ を $\mathcal{A}_{total}$ へ送り，
\item[(3)] $a^{-1}$ は $D_\mathcal{A}(\mathcal{D})$ を
$D_{\mathcal{A}_{total}}(\mathcal{C}_{total})$ へ送る．
\end{enumerate}
すべての $n, q$ に対して $R^qa_{n, *}$ が $\mathcal{A}_n$ を
$\mathcal{A}$ へ送るならば，
\begin{enumerate}
\item[(4)] すべての $q$ に対して $R^qa_*$ は $\mathcal{A}_{total}$ を
$\mathcal{A}$ へ送り，
\item[(5)] $Ra_*$ は $D_{\mathcal{A}_{total}}^+(\mathcal{C}_{total})$ を
$D_\mathcal{A}^+(\mathcal{D})$ へ送る．
\end{enumerate}
\end{lemma}

\begin{proof}
興味深い主張は (4) と (5) だけである．(4) は Simplicial Spaces の Lemma
\ref{spaces-simplicial-lemma-augmentation-spectral-sequence} にあるスペクトル系列と，
Homology の Lemma \ref{homology-lemma-biregular-ss-converges} から従う．
次に (5) は，切断によって与えられる
$D_{\mathcal{A}_{total}}^+(\mathcal{C}_{total})$ の対象 $K$ 上の標準フィルトレーションに
付随するスペクトル系列を考えることから従う．詳細は省略する．
\end{proof}

\begin{remark}
\label{obsolete-remark-cohomology-topics}
このタグは以前，扱うべき話題を列挙したコホモロジーの章のある節を指していた．
\end{remark}

\begin{remark}
\label{obsolete-remark-sites-cohomology-topics}
このタグは以前，扱うべき話題を列挙したコホモロジーの章のある節を指していた．
\end{remark}

\begin{remark}
\label{obsolete-remark-V-implies-C}
このタグは以前，Cohomology on Sites の Lemma
\ref{sites-cohomology-lemma-V-implies-C-general} の特別な場合であって，
Cohomology on Sites の Lemma
\ref{sites-cohomology-lemma-compare-qc-zar} に記述された状況に関するものを指していた．
\end{remark}
% Frozen obsolete.tex lines 1225--1439.
\begin{remark}
\label{obsolete-remark-V-implies-cohomology}
このタグは以前，Cohomology on Sites の Lemma
\ref{sites-cohomology-lemma-V-implies-cohomology-general} の特別な場合であって，
Cohomology on Sites の Lemma
\ref{sites-cohomology-lemma-compare-qc-zar} に記述された状況に関するものを指していた．
\end{remark}

\begin{remark}
\label{obsolete-remark-induction-step-V-C}
このタグは以前，Cohomology on Sites の Lemma
\ref{sites-cohomology-lemma-induction-step-V-C-general} の特別な場合であって，
Cohomology on Sites の Lemma
\ref{sites-cohomology-lemma-compare-qc-zar} に記述された状況に関するものを指していた．
\end{remark}

\begin{remark}
\label{obsolete-remark-V-implies-C-etale-fppf}
このタグは以前，Cohomology on Sites の Lemma
\ref{sites-cohomology-lemma-V-implies-C-general} の特別な場合であって，
\'Etale Cohomology の Lemma
\ref{etale-cohomology-lemma-compare-fppf-etale} に記述された状況に関するものを
指していた．
\end{remark}

\begin{remark}
\label{obsolete-remark-V-implies-cohomology-etale-fppf}
このタグは以前，Cohomology on Sites の Lemma
\ref{sites-cohomology-lemma-V-implies-cohomology-general} の特別な場合であって，
\'Etale Cohomology の Lemma
\ref{etale-cohomology-lemma-compare-fppf-etale} に記述された状況に関するものを
指していた．
\end{remark}

\begin{remark}
\label{obsolete-remark-induction-step-V-C-etale-fppf}
このタグは以前，Cohomology on Sites の Lemma
\ref{sites-cohomology-lemma-induction-step-V-C-general} の特別な場合であって，
\'Etale Cohomology の Lemma
\ref{etale-cohomology-lemma-compare-fppf-etale} に記述された状況に関するものを
指していた．
\end{remark}

\begin{remark}
\label{obsolete-remark-V-implies-C-etale-ph}
このタグは以前，Cohomology on Sites の Lemma
\ref{sites-cohomology-lemma-V-implies-C-general} の特別な場合であって，
\'Etale Cohomology の Lemma
\ref{etale-cohomology-lemma-compare-ph-etale} に記述された状況に関するものを
指していた．
\end{remark}

\begin{remark}
\label{obsolete-remark-V-implies-cohomology-etale-ph}
このタグは以前，Cohomology on Sites の Lemma
\ref{sites-cohomology-lemma-V-implies-cohomology-general} の特別な場合であって，
\'Etale Cohomology の Lemma
\ref{etale-cohomology-lemma-compare-ph-etale} に記述された状況に関するものを
指していた．
\end{remark}

\begin{remark}
\label{obsolete-remark-V-implies-cohomology-etale-ph-extra}
このタグは以前，Cohomology on Sites の Lemma
\ref{sites-cohomology-lemma-V-implies-cohomology-extra-general} の特別な場合であって，
\'Etale Cohomology の Lemma
\ref{etale-cohomology-lemma-compare-ph-etale} に記述された状況に関するものを
指していた．
\end{remark}

\begin{remark}
\label{obsolete-remark-make-class-zero}
このタグは以前，Cohomology on Sites の Lemma
\ref{sites-cohomology-lemma-make-class-zero-general} の特別な場合であって，
\'Etale Cohomology の Lemma
\ref{etale-cohomology-lemma-compare-ph-etale} に記述された状況に関するものを
指していた．
\end{remark}

\begin{remark}
\label{obsolete-remark-induction-step-V-C-etale-ph}
このタグは以前，Cohomology on Sites の Lemma
\ref{sites-cohomology-lemma-induction-step-V-C-general} の特別な場合であって，
\'Etale Cohomology の Lemma
\ref{etale-cohomology-lemma-compare-ph-etale} に記述された状況に関するものを
指していた．
\end{remark}

\begin{remark}
\label{obsolete-remark-V-implies-C-etale-h}
このタグは以前，Cohomology on Sites の Lemma
\ref{sites-cohomology-lemma-V-implies-C-general} の特別な場合であって，
\'Etale Cohomology の Lemma
\ref{etale-cohomology-lemma-compare-h-etale} に記述された状況に関するものを
指していた．
\end{remark}

\begin{remark}
\label{obsolete-remark-V-implies-cohomology-etale-h}
このタグは以前，Cohomology on Sites の Lemma
\ref{sites-cohomology-lemma-V-implies-cohomology-general} の特別な場合であって，
\'Etale Cohomology の Lemma
\ref{etale-cohomology-lemma-compare-h-etale} に記述された状況に関するものを
指していた．
\end{remark}

\begin{remark}
\label{obsolete-remark-V-implies-cohomology-etale-h-extra}
このタグは以前，Cohomology on Sites の Lemma
\ref{sites-cohomology-lemma-V-implies-cohomology-extra-general} の特別な場合であって，
\'Etale Cohomology の Lemma
\ref{etale-cohomology-lemma-compare-h-etale} に記述された状況に関するものを
指していた．
\end{remark}

\begin{remark}
\label{obsolete-remark-induction-step-V-C-etale-h}
このタグは以前，Cohomology on Sites の Lemma
\ref{sites-cohomology-lemma-induction-step-V-C-general} の特別な場合であって，
\'Etale Cohomology の Lemma
\ref{etale-cohomology-lemma-compare-h-etale} に記述された状況に関するものを
指していた．
\end{remark}

\begin{remark}
\label{obsolete-remark-how-used}
このタグは以前，\'etale cohomology の章にあったが，再編成のため現在では
そこに置くのは適切でない．タグの内容は次のとおりであった：
\'Etale Cohomology の Lemma \ref{etale-cohomology-lemma-when-ctf} を用いると，
$f : X \to Y$ がスキームの分離的有限型射であり，$Y$ が Noether ならば，
$Rf_!$ が函手
$D_{ctf}(X_\etale, \Lambda) \to D_{ctf}(Y_\etale, \Lambda)$ を誘導することを
証明できる．$Y$ が体のスペクトルであり $X$ が曲線である場合のこの議論の例は，
The Trace Formula の Proposition
\ref{trace-proposition-projective-curve-constructible-cohomology} にある．
\end{remark}

\begin{lemma}
\label{obsolete-lemma-glue-f-upper-shriek}
$f : X \to Y$ をスキームの局所準有限射とする．次を満たす一意な函手
$f^! : \textit{Ab}(Y_\etale) \to \textit{Ab}(X_\etale)$ が存在する：
\begin{enumerate}
\item $f \circ j$ が分離的となる任意の開射 $j : U \to X$ に対して，標準同型
$j^! \circ f^! = (f \circ j)^!$ が存在する；
\item $U \subset U' \subset X$ に対するこれらの同型は，More \'Etale Cohomology の
Lemma \ref{more-etale-lemma-upper-shriek-restriction} にある同型と両立する．
\end{enumerate}
\end{lemma}

\begin{proof}
More \'Etale Cohomology の Lemma
\ref{more-etale-lemma-lqf-f-upper-shriek} と Lemma
\ref{more-etale-lemma-upper-shriek-restriction} から直ちに従う．
\end{proof}

\begin{proposition}
\label{obsolete-proposition-lqf-shriek}
$f : X \to Y$ を局所準有限射とする．随伴函手
$f_! : \textit{Ab}(X_\etale) \to \textit{Ab}(Y_\etale)$ および
$f^! : \textit{Ab}(Y_\etale) \to \textit{Ab}(X_\etale)$ で，次の性質をもつものが
存在する：
\begin{enumerate}
\item 函手 $f^!$ は More \'Etale Cohomology の Lemma
\ref{more-etale-lemma-lqf-f-upper-shriek} で構成されたものである；
\item $f \circ j$ が分離的となる任意の開射 $j : U \to X$ に対して，標準同型
$f_! \circ j_! = (f \circ j)_!$ が存在する；
\item $U \subset U' \subset X$ に対するこれらの同型は，More \'Etale Cohomology の
Lemma \ref{more-etale-lemma-f-shriek-composition} にある同型と両立する．
\end{enumerate}
\end{proposition}

\begin{proof}
More \'Etale Cohomology の Section
\ref{more-etale-section-finite-support} と Section
\ref{more-etale-section-duality-locally-quasi-finite} を参照せよ．
\end{proof}

\begin{lemma}
\label{obsolete-lemma-lqf-f-upper-shriek-stalk}
$f : X \to Y$ をスキームの局所準有限射とする．アーベル群 $A$ と幾何点
$\overline{y} : \Spec(k) \to Y$ に対して，
$f^!(\overline{y}_*A) = \prod\nolimits_{f(\overline{x}) = \overline{y}}
\overline{x}_*A$ である．
\end{lemma}

\begin{proof}
More \'Etale Cohomology の Lemma
\ref{more-etale-lemma-lqf-f-upper-shriek} にある対応する主張から従う．
\end{proof}

\begin{lemma}
\label{obsolete-lemma-lqf-f-shriek-composition}
$f : X \to Y$ と $g : Y \to Z$ を合成可能なスキームの局所準有限射とする．このとき
$g_! \circ f_! = (g \circ f)_!$ および
$f^! \circ g^! = (g \circ f)^!$ である．
\end{lemma}

\begin{proof}
More \'Etale Cohomology の Lemma
\ref{more-etale-lemma-lqf-shriek-composition} と Lemma
\ref{more-etale-lemma-upper-shriek-restriction} を組み合わせればよい．
\end{proof}
% Frozen obsolete.tex lines 1440--1586.
\section{微分次数付き代数}
\label{obsolete-section-dga}


\begin{lemma}
\label{obsolete-lemma-P-not-preserved-base-change}
$(A, \text{d})$ と $(B, \text{d})$ を微分次数付き代数とする．
$N$ を性質 (P) をもつ微分次数付き $(A, B)$-双加群とする．
$M$ を性質 (P) をもつ微分次数付き $A$-加群とする．このとき
$Q = M \otimes_A N$ は，$D(B)$ において $M \otimes_A^\mathbf{L} N$ を表す
微分次数付き $B$-加群であり，さらに微分次数付き部分加群によるフィルトレーション
$$
0 = F_{-1}Q \subset F_0Q \subset F_1Q \subset \ldots \subset Q
$$
をもつ．ここで $Q = \bigcup F_pQ$ であり，包含
$F_iQ \to F_{i + 1}Q$ は許容単射，商 $F_{i + 1}Q/F_iQ$ は微分次数付き
$B$-加群として $(A \otimes_R B)[k]$ の直和と同型である．
\end{lemma}

\begin{proof}
$M$ と $N$ 上のフィルトレーション $F_\bullet$ を選ぶ．
$Q = M \otimes_A N$ 上のフィルトレーションを
$$
F_n(Q) = \sum\nolimits_{i + j = n}
F_i(M) \otimes_A F_j(N)
$$
で定める．これは明らかに微分次数付き $B$-部分加群である．
$$
F_n(Q)/F_{n - 1}(Q) =
\bigoplus\nolimits_{i + j = n}
F_i(M)/F_{i - 1}(M) \otimes_A F_j(N)/F_{j - 1}(N)
$$
である．たとえば，これは $M$ のフィルトレーションが次数付き $A$-加群の圏で
分裂することから従う．仮定により右辺の商は $A$ と $A \otimes_R B$ のシフトの
直和と同型であり，また
$A \otimes_A (A \otimes_R B) = A \otimes_R B$ である．したがって左辺は
微分次数付き $B$-加群として $A \otimes_R B$ のシフトの直和である．
（注意：$Q$ は $(A, B)$-双加群の構造をもたない．）
これで補題の第1の主張が証明された．第2の主張は Differential Graded Algebra の
Lemma \ref{dga-lemma-derived-bc} にある函手の定義から直ちに従う．
\end{proof}










\section{単体的方法}
\label{obsolete-section-simplicial}


\begin{lemma}
\label{obsolete-lemma-equiv}
仮定と記法は Simplicial の Lemma \ref{simplicial-lemma-section} と同じとする．
射 $f$ の切断 $g : U \to V$ が存在し，合成 $g \circ f$ は $V$ 上の恒等写像と
ホモトピー同値である．とくに，射 $f$ はホモトピー同値である．
\end{lemma}

\begin{proof}
Simplicial の Lemma \ref{simplicial-lemma-section} と Lemma
\ref{simplicial-lemma-trivial-kan-homotopy} から直ちに従う．
\end{proof}

\begin{lemma}
\label{obsolete-lemma-cosk-hom-deltak}
$\mathcal{C}$ を有限余積と有限極限をもつ圏とする．$X$ を $\mathcal{C}$ の対象とする．
$k \geq 0$ とする．標準写像
$$
\Hom(\Delta[k], X)
\longrightarrow
\text{cosk}_1 \text{sk}_1 \Hom(\Delta[k], X)
$$
は同型である．
\end{lemma}

\begin{proof}
任意の単体的対象 $V$ に対して，
\begin{eqnarray*}
\Mor(V, \text{cosk}_1 \text{sk}_1 \Hom(\Delta[k], X))
& = &
\Mor(\text{sk}_1 V, \text{sk}_1 \Hom(\Delta[k], X)) \\
& = &
\Mor(i_{1!} \text{sk}_1 V, \Hom(\Delta[k], X)) \\
& = &
\Mor(i_{1!} \text{sk}_1 V \times \Delta[k], X)
\end{eqnarray*}
である．第1の等式は $\text{sk}$ と $\text{cosk}$ の随伴性による．
第2の等式は $i_{1!}$ と $\text{sk}_1$ の随伴性による．第1の等式は Simplicial の
Definition \ref{simplicial-definition-hom-from-simplicial-set} による．ここで最後の
$X$ は値 $X$ をもつ定単体的対象を表す．Simplicial の Lemma
\ref{simplicial-lemma-augmentation-howto} により，この集合の元は
$i_{1!} \text{sk}_1 V \times \Delta[k]$ の次数 $0$ と $1$ の項だけに依存する．
これらは $V \times \Delta[k]$ の次数 $0$ と $1$ の項に一致する；Simplicial の Lemma
\ref{simplicial-lemma-recovering-U-for-real} を参照せよ．したがって上の集合は
$\Mor(V \times \Delta[k], X) = \Mor(V, \Hom(\Delta[k], X))$ に等しい．
\end{proof}

\begin{lemma}
\label{obsolete-lemma-cosk0-hom-deltak}
$\mathcal{C}$ を圏とする．$X$ を，自己積 $X \times \ldots \times X$ が存在する
$\mathcal{C}$ の対象とする．$k \geq 0$ とし，$C[k]$ は Simplicial の Example
\ref{simplicial-example-simplex-cosimplicial-set} にあるものとする．
Simplicial の Lemma \ref{simplicial-lemma-morphism-into-product} の記法の下で，標準写像
$$
\Hom(C[k], X)_1
\longrightarrow
(\text{cosk}_0 \text{sk}_0 \Hom(C[k], X))_1
$$
は，写像
$$
\prod\nolimits_{\alpha : [k] \to [1]} X
\longrightarrow
X \times X
$$
と同一視される．これは $\alpha$ が定値写像である因子への射影である．
\end{lemma}

\begin{proof}
これは Hypercoverings の Lemma
\ref{hypercovering-lemma-covering} の証明中で示されている．
\end{proof}
% Frozen obsolete.tex lines 1587--1755.
\section{スキームに関する結果}
\label{obsolete-section-devissage}

\noindent
余分と思われる補題を挙げる．

\begin{lemma}
\label{obsolete-lemma-stein-projective}
$(R, \mathfrak m, \kappa)$ を局所環とする．
$X \subset \mathbf{P}^n_R$ を閉部分スキームとする．
$R = \Gamma(X, \mathcal{O}_X)$ と仮定する．このとき特殊ファイバー
$X_k$ は幾何学的に連結である．
\end{lemma}

\begin{proof}
これは More on Morphisms の Theorem
\ref{more-morphisms-theorem-stein-factorization-general} の特別な場合である．
\end{proof}

\begin{lemma}
\label{obsolete-lemma-property-irreducible-higher-rank}
$X$ を Noether スキームとする．$Z_0 \subset X$ を生成点 $\xi$ をもつ既約閉集合とする．
$\mathcal{P}$ を $X$ 上の連接層の性質で，次を満たすものとする：
\begin{enumerate}
\item 連接層の任意の短完全列において，三つのうち二つが性質
$\mathcal{P}$ をもてば，残る一つもこの性質をもつ．
\item 連接層の直和が $\mathcal{P}$ を満たすならば，双方がこれを満たす．
\item 任意の整閉部分スキーム $Z \subset Z_0 \subset X$,
$Z \not = Z_0$ と，任意の準連接イデアル層
$\mathcal{I} \subset \mathcal{O}_Z$ に対して，
$(Z \to X)_*\mathcal{I}$ は $\mathcal{P}$ を満たす．
\item $X$ 上のある連接層 $\mathcal{G}$ が存在し，次を満たす：
\begin{enumerate}
\item $\text{Supp}(\mathcal{G}) = Z_0$；
\item $\mathcal{G}_\xi$ は $\mathfrak m_\xi$ によって零化される；
\item $\mathcal{G}$ は性質 $\mathcal{P}$ を満たす．
\end{enumerate}
\end{enumerate}
このとき，その台が $Z_0$ に含まれる $X$ 上の任意の連接層
$\mathcal{F}$ は性質 $\mathcal{P}$ を満たす．
\end{lemma}

\begin{proof}
証明は Cohomology of Schemes の Lemma
\ref{coherent-lemma-property-irreducible} の証明の変形である．その証明とまったく同様に，
台が $Z_0$ に真に含まれる任意の連接層は性質 $\mathcal{P}$ をもつことが分かる．

\medskip\noindent
(3) にあるような連接層 $\mathcal{G}$ を考える．Cohomology of Schemes の Lemma
\ref{coherent-lemma-prepare-filter-irreducible} により，$Z_0$ 上のイデアル層
$\mathcal{I}$ と短完全列
$$
0 \to
\left((Z_0 \to X)_*\mathcal{I}\right)^{\oplus r} \to
\mathcal{G} \to
\mathcal{Q} \to 0
$$
が存在する．ここで $\mathcal{Q}$ の台は $Z_0$ に真に含まれる．とくに
$\mathcal{G}$ の台は $Z$ に等しいので，$r > 0$ であり，$\mathcal{I}$ は零でない．
$\mathcal{Q}$ は性質 $\mathcal{P}$ をもつから，
$\left((Z_0 \to X)_*\mathcal{I}\right)^{\oplus r}$ も性質 $\mathcal{P}$ をもつ．
(2) により，$(Z_0 \to X)_*\mathcal{I}$ が性質 $\mathcal{P}$ をもつと分かる．
これを Cohomology of Schemes の Lemma
\ref{coherent-lemma-property-irreducible} の証明の適切な箇所に組み込めば補題を得る．
いくつかの細部は省略する．
\end{proof}

\begin{lemma}
\label{obsolete-lemma-property-higher-rank}
$X$ を Noether スキームとする．$\mathcal{P}$ を $X$ 上の連接層の性質で，
次を満たすものとする：
\begin{enumerate}
\item 連接層の任意の短完全列において，三つのうち二つが性質
$\mathcal{P}$ をもてば，残る一つもこの性質をもつ．
\item 連接層の直和が $\mathcal{P}$ を満たすならば，双方がこれを満たす．
\item 生成点 $\xi$ をもつ任意の整閉部分スキーム $Z \subset X$ に対して，
ある連接層 $\mathcal{G}$ が存在し，次を満たす：
\begin{enumerate}
\item $\text{Supp}(\mathcal{G}) = Z$；
\item $\mathcal{G}_\xi$ は $\mathfrak m_\xi$ によって零化される；
\item $\mathcal{G}$ は性質 $\mathcal{P}$ を満たす．
\end{enumerate}
\end{enumerate}
このとき，$X$ 上の任意の連接層は性質 $\mathcal{P}$ を満たす．
\end{lemma}

\begin{proof}
これは Lemma \ref{obsolete-lemma-property-irreducible-higher-rank} から，
Cohomology of Schemes の Lemma \ref{coherent-lemma-property} が
Cohomology of Schemes の Lemma
\ref{coherent-lemma-property-irreducible} から従うのとまったく同じ仕方で従う．
\end{proof}

\begin{lemma}
\label{obsolete-lemma-section-maps-back-into}
$X$ をスキームとする．$\mathcal{L}$ を可逆 $\mathcal{O}_X$-加群とする．
$s \in \Gamma(X, \mathcal{L})$ を切断とする．
$\mathcal{F}' \subset \mathcal{F}$ を準連接 $\mathcal{O}_X$-加群とする．次を仮定する：
\begin{enumerate}
\item $X$ は準コンパクトである；
\item $\mathcal{F}$ は有限型である；
\item $\mathcal{F}'|_{X_s} = \mathcal{F}|_{X_s}$．
\end{enumerate}
このとき，$\mathcal{F}$ 上の $s^n$ による乗法が $\mathcal{F}'$ を経由するような
$n \geq 0$ が存在する．
\end{lemma}

\begin{proof}
言い換えると，ある $n \geq 0$ に対して
$s^n\mathcal{F} \subset
\mathcal{F}' \otimes_{\mathcal{O}_X} \mathcal{L}^{\otimes n}$ であると主張している．
さらに言い換えると，商写像 $\mathcal{F} \to \mathcal{F}/\mathcal{F}'$ は
$s$ のある冪を掛けると零になると主張している．これは Properties の Lemma
\ref{properties-lemma-section-maps-backwards} から従う．
\end{proof}

\begin{lemma}
\label{obsolete-lemma-bound-degree-in-nbhd-generic-point}
$f : X \to Y$ をスキームの射とする．次を仮定する：
\begin{enumerate}
\item $X$ と $Y$ は整スキームである；
\item $f$ は局所有限型かつ支配的である；
\item $f$ は準コンパクトまたは分離的である；
\item $f$ は生成的に有限である，すなわち Morphisms の Lemma
\ref{morphisms-lemma-finite-degree} の (1)--(5) のいずれかが成り立つ．
\end{enumerate}
このとき，$f^{-1}(V) \to V$ が次数 $\deg(X/Y)$ の有限局所自由射となるような
空でない開集合 $V \subset Y$ が存在する．とくに，$f^{-1}(V) \to V$ のファイバーの
次数は $\deg(X/Y)$ で上から抑えられる．
\end{lemma}

\begin{proof}
$f^{-1}(V) \to V$ が有限となるように $V$ を選べる．次に $V$ を縮小し，
生成的平坦性（Morphisms の Proposition
\ref{morphisms-proposition-generic-flatness}）により，$f^{-1}(V) \to V$ が
平坦かつ有限表示であると仮定できる．すると Morphisms の Lemma
\ref{morphisms-lemma-finite-flat} により，この射は有限局所自由である．
$V$ は既約なので，この射は一定の次数をもつ．最後の主張はこれと Morphisms の Lemma
\ref{morphisms-lemma-finite-locally-free-universally-bounded} から従う．
\end{proof}
% Frozen obsolete.tex lines 1756--1995.
\section{多様体の導来圏}
\label{obsolete-section-equiv}

\noindent
もともと多様体の導来圏の章の一部であったが，現在では不要になった補題を挙げる．

\begin{lemma}
\label{obsolete-lemma-preserves-Coh}
$k$ を体とする．$X$ を $k$ 上有限型で分離的かつ正則なスキームとする．
$F : D_{perf}(\mathcal{O}_X) \to D_{perf}(\mathcal{O}_X)$ を
$k$-線形完全函手とする．$\dim(\text{Supp}(\mathcal{F})) = 0$ を満たす任意の連接
$\mathcal{O}_X$-加群 $\mathcal{F}$ に対して，$D_{perf}(\mathcal{O}_X)$ の
$M$ に関して函手的な $k$-ベクトル空間の同型
$$
\Hom_X(\mathcal{F}, M) = \Hom_X(\mathcal{F}, F(M))
$$
があると仮定する．このとき，台となる位相空間上で恒等写像を誘導する
$k$ 上の自己同型 $f : X \to X$
\footnote{これはしばしば $f$ が恒等写像であることを強制する；Varieties の Lemma
\ref{varieties-lemma-automorphism} を参照せよ．} と，可逆 $\mathcal{O}_X$-加群
$\mathcal{L}$ が存在し，$F$ と
$F'(M) = f^*M \otimes_{\mathcal{O}_X}^\mathbf{L} \mathcal{L}$ は sibling である．
\end{lemma}

\begin{proof}
Derived Categories of Varieties の Lemma \ref{equiv-lemma-duality-at-point} により，
台が閉点である任意の連接 $\mathcal{O}_X$-加群 $\mathcal{F}$ に対して，
$M$ に関して函手的な同型
$$
H^0(X, M \otimes^\mathbf{L}_{\mathcal{O}_X} \mathcal{F}) =
H^0(X, F(M) \otimes^\mathbf{L}_{\mathcal{O}_X} \mathcal{F})
$$
がある．

\medskip\noindent
$x \in X$ を閉点とし，上の主張を，$x$ で値 $\kappa(x)$ をとる摩天楼層
$\mathcal{F} = \mathcal{O}_x$ に適用する．すべての $p \in \mathbf{Z}$ に対して
$$
\dim_{\kappa(x)}
\text{Tor}^{\mathcal{O}_{X, x}}_p(M_x, \kappa(x)) =
\dim_{\kappa(x)}
\text{Tor}^{\mathcal{O}_{X, x}}_p(F(M)_x, \kappa(x))
$$
を得る．とくに，$i > 0$ に対して $H^i(M) = 0$ ならば，Derived Categories of
Varieties の Lemma \ref{equiv-lemma-orthogonal-point-sheaf} により，
$i > 0$ に対して $H^i(F(M)) = 0$ である．

\medskip\noindent
$\mathcal{E}$ が階数 $r$ の局所自由加群ならば，$F(\mathcal{E})$ も階数 $r$ の
局所自由加群である．実際，$\mathcal{O}_{X, x}$ 上の完全複体 $K$ が
$$
\dim_{\kappa(x)} \text{Tor}^{\mathcal{O}_{X, x}}_i(K, \kappa(x)) =
\left\{
\begin{matrix}
r & \text{if} & i = 0 \\
0 & \text{if} & i \not = 0
\end{matrix}
\right.
$$
を満たすならば，それは次数 $0$ に置かれた階数 $r$ の自由加群に等しい．
たとえば More on Algebra の Lemma
\ref{more-algebra-lemma-lift-perfect-from-residue-field} を参照せよ．

\medskip\noindent
$M$ の台が閉部分スキーム $Z \subset X$ に含まれるならば，$F(M)$ の台も
$Z$ に含まれる．実際，$x \not \in Z$ に対して
$M \otimes_{\mathcal{O}_X}^\mathbf{L} \mathcal{O}_x = 0$ であり，したがって
$F(M)$ についても同じことが成り立つ．ゆえに Derived Categories of Varieties の
Lemma \ref{equiv-lemma-orthogonal-point-sheaf} から結論を得る．

\medskip\noindent
とくに $F(\mathcal{O}_x)$ の台は $\{x\}$ に含まれる．
$H^i(\mathcal{O}_x) \not = 0$ となる最小の整数を $i \in \mathbf{Z}$ とする．
$i \leq 0$ であることが分かっている．もし $i < 0$ ならば，射
$\mathcal{O}_x[-i] \to F(\mathcal{O}_x)$ が存在し，これはすべての射
$\mathcal{O}_x[-i] \to \mathcal{O}_x$ が零であることに矛盾する．したがって
$F(\mathcal{O}_x) = \mathcal{H}[0]$ である．ここで $\mathcal{H}$ は
$x$ における摩天楼層である．

\medskip\noindent
$\mathcal{G}$ を $\dim(\text{Supp}(\mathcal{G})) = 0$ を満たす連接
$\mathcal{O}_X$-加群とする．このときフィルトレーション
$$
0 = \mathcal{G}_0 \subset \mathcal{G}_1 \subset \ldots \subset
\mathcal{G}_n = \mathcal{G}
$$
であって，$n \geq i \geq 1$ に対して商 $\mathcal{G}_i/\mathcal{G}_{i - 1}$ が
ある閉点 $x_i \in X$ に対する $\mathcal{O}_{x_i}$ と同型であるものが存在する．
すると識別三角形
$$
F(\mathcal{G}_{i - 1}) \to F(\mathcal{G}_i) \to F(\mathcal{O}_{x_i})
$$
を得る．帰納法により，$F(\mathcal{G}_i)$ は次数 $0$ に置かれた連接層である．

\medskip\noindent
$\mathcal{G}$ を連接 $\mathcal{O}_X$-加群とする．$i > 0$ に対して
$H^i(F(\mathcal{G})) = 0$ であることが分かっている．矛盾を導くため，
ある $i < 0$ に対して $H^i(F(\mathcal{G}))$ が零でないと仮定する．
この性質をもつ最小の $i$ を選ぶと，$D_{perf}(\mathcal{O}_X)$ に射
$H^i(F(\mathcal{G}))[-i] \to F(\mathcal{G})$ がある．
$H^i(F(\mathcal{G}))$ の台にある閉点 $x \in X$ を選ぶ．More on Algebra の Lemma
\ref{more-algebra-lemma-kollar-kovacs-pseudo-coherent} により，写像
$$
H^i(F(\mathcal{G}))_x \otimes_{\mathcal{O}_{X, x}}
\mathcal{O}_{X, x}/\mathfrak m_x^n
\longrightarrow
\text{Tor}^{\mathcal{O}_{X, x}}_{-i}(F(\mathcal{G})_x,
\mathcal{O}_{X, x}/\mathfrak m_x^n)
$$
が零でないような $n > 0$ が存在する．次に $m \geq 1$ を取り，短完全列
$$
0 \to \mathfrak m_x^m \mathcal{G} \to \mathcal{G} \to
\mathcal{G}/\mathfrak m_x^m\mathcal{G} \to 0
$$
を考える．上で見たように，$F(\mathcal{G}/\mathfrak m_x^m\mathcal{G})$ は
次数 $0$ に置かれた層である．したがって
$H^i(F(\mathfrak m_x^m \mathcal{G})) \to H^i(F(\mathcal{G}))$ は同型である．
可換図式
$$
\xymatrix{
H^i(F(\mathfrak m_x^m\mathcal{G}))_x \otimes_{\mathcal{O}_{X, x}}
\mathcal{O}_{X, x}/\mathfrak m_x^n \ar[r] \ar[d] &
\text{Tor}^{\mathcal{O}_{X, x}}_{-i}(F(\mathfrak m_x^m\mathcal{G})_x,
\mathcal{O}_{X, x}/\mathfrak m_x^n) \ar[d] \\
H^i(F(\mathcal{G}))_x \otimes_{\mathcal{O}_{X, x}}
\mathcal{O}_{X, x}/\mathfrak m_x^n \ar[r] &
\text{Tor}^{\mathcal{O}_{X, x}}_{-i}(F(\mathcal{G})_x,
\mathcal{O}_{X, x}/\mathfrak m_x^n)
}
$$
を考える．左の垂直矢印は同型で，下の矢印は零でないので，すべての
$m \geq 1$ に対して右の垂直矢印は零でないと結論できる．他方，証明の第1段落により，
この矢印は写像
$$
\text{Tor}^{\mathcal{O}_{X, x}}_{-i}(\mathfrak m_x^m\mathcal{G}_x,
\mathcal{O}_{X, x}/\mathfrak m_x^n)
\longrightarrow
\text{Tor}^{\mathcal{O}_{X, x}}_{-i}(\mathcal{G}_x,
\mathcal{O}_{X, x}/\mathfrak m_x^n)
$$
と同型であることが分かっている．しかし More on Algebra の Lemma
\ref{more-algebra-lemma-tor-annihilated} により，$m \gg n$ に対してこの矢印は零である．
これは求めていた矛盾である．

\medskip\noindent
以上から $F$ は連接加群を保つ．Derived Categories of Varieties の Lemma
\ref{equiv-lemma-exact-functor-preserving-Coh} により，$F$ は Fourier--Mukai 函手
$F'$ と sibling である．この函手は，$\text{pr}_1$ を介して $X$ 上平坦であり，
$\text{pr}_2$ を介して $X$ 上有限である連接 $\mathcal{O}_{X \times X}$-加群
$\mathcal{K}$ によって与えられる．$F(\mathcal{O}_X)$ は次数 $0$ に置かれた
可逆 $\mathcal{O}_X$-加群 $\mathcal{L}$ なので，
$$
\mathcal{L} \cong F(\mathcal{O}_X) \cong F'(\mathcal{O}_X) \cong
\text{pr}_{2, *}\mathcal{K}
$$
である．したがって Functors and Morphisms の Lemma
\ref{functors-lemma-pushforward-invertible-pre} により，
$\text{pr}_2 \circ s = \text{id}_X$ かつ $\mathcal{K} = s_*\mathcal{L}$ を満たす射
$s : X \to X \times X$ が存在する．$f = \text{pr}_1 \circ s$ と置く．このとき
\begin{align*}
F'(M)
& =
R\text{pr}_{2, *}(L\text{pr}_1^*K \otimes \mathcal{K}) \\
& =
R\text{pr}_{2, *}(L\text{pr}_1^*M \otimes s_*\mathcal{L}) \\
& =
R\text{pr}_{2, *}(Rs_*(Lf^*M \otimes \mathcal{L})) \\
& =
Lf^*M \otimes \mathcal{L}
\end{align*}
である．第3段階では Derived Categories of Schemes の Lemma
\ref{perfect-lemma-cohomology-base-change} を用いた．すべての閉点 $x \in X$ に対して
加群 $F(\mathcal{O}_x)$ の台は $x$ に含まれるので，$f$ は $X$ の台となる
位相空間上の恒等写像を誘導する．$f$ が同型であることはまだ示す必要があり，
次の段落でこれを行う．

\medskip\noindent
$x \in X$ を閉点とする．$n \geq 1$ に対して，値
$\mathcal{O}_{X, x}/\mathfrak m_x^n$ をもつ $x$ における摩天楼層を
$\mathcal{O}_{x, n}$ と記す．
$$
\Hom_X(\mathcal{O}_{x, m}, \mathcal{O}_{x, n}) \cong
\Hom_X(\mathcal{O}_{x, m}, F(\mathcal{O}_{x, n})) \cong
\Hom_X(\mathcal{O}_{x, m}, f^*\mathcal{O}_{x, n} \otimes \mathcal{L})
$$
という同型があり，$\mathcal{O}_{x, n}$ どうしの $\mathcal{O}_X$-加群準同型に関して
函手的である．（第1の同型は仮定により存在し，第2の同型は $F$ と $F'$ が
sibling であることによる．）$m \geq n$ に対して，$\mathcal{O}_{x, n}$ への作用を介して
$\mathcal{O}_{X, x}/\mathfrak m^n =
\Hom_X(\mathcal{O}_{x, m}, \mathcal{O}_{x, n})$ であるから，すべての $n$ に対して
$f^\sharp : \mathcal{O}_{X, x}/\mathfrak m_x^n \to
\mathcal{O}_{X, x}/\mathfrak m_x^n$ は全単射であると結論できる．したがって $f$ は
閉点における完備局所環上で同型を誘導し，ゆえに \'etale である
（\'Etale Morphisms の Lemma \ref{etale-lemma-characterize-etale-completions}）．
閉点を見ると，$f$ は \'etale なので開埋め込みである
$\Delta_f : X \to X \times_{f, X, f} X$ は全単射であり，したがって同型である．
ゆえに $f$ は単射射である．最後に，$f$ は同型であると結論する．Descent の Lemma
\ref{descent-lemma-flat-surjective-quasi-compact-monomorphism-isomorphism} により，
それは開埋め込みだからである．
\end{proof}
% Frozen obsolete.tex lines 1996--2129.
\section{正則固有の場合の表現可能性}
\label{obsolete-section-regular-proper}

\noindent
Derived Categories of Varieties の Theorem
\ref{equiv-theorem-bondal-van-den-bergh} を，体上のすべての固有スキームに適用できる
ように改良したので，この節は旧結果である（以前は体上の射影スキームについてのみ
証明していた）．

\begin{lemma}
\label{obsolete-lemma-trace-map}
\begin{reference}
ここで与える証明は \cite[Remark 3.4]{MS} の議論に従う．
\end{reference}
$f : X' \to X$ を整 Noether スキームの固有双有理射とし，$X$ は正則であるとする．
写像 $\mathcal{O}_X \to Rf_*\mathcal{O}_{X'}$ は
$D(\mathcal{O}_X)$ において標準的に分裂する．
\end{lemma}

\begin{proof}
$D(\mathcal{O}_X)$ において $E = Rf_*\mathcal{O}_{X'}$ と置く．
Derived Categories of Schemes の Lemma
\ref{perfect-lemma-direct-image-coherent} により，
$E$ は $D^b_{\textit{Coh}}(\mathcal{O}_X)$ に属することに注意する．
Derived Categories of Schemes の Lemma
\ref{perfect-lemma-perfect-on-regular} により，$E$ は
$D(\mathcal{O}_X)$ の完全対象である．$\mathcal{O}_{X'}$ は代数の層なので，
Cohomology の Remark \ref{cohomology-remark-cup-product} により，相対カップ積
$\mu : E \otimes_{\mathcal{O}_X}^\mathbf{L} E \to E$ がある．
$\sigma : E \otimes E^\vee \to E^\vee \otimes E$ を，対称モノイド圏
$D(\mathcal{O}_X)$ 上の可換制約とする（Cohomology の Lemma
\ref{cohomology-lemma-symmetric-monoidal-derived}）．
$\eta : \mathcal{O}_X \to E \otimes E^\vee$ および
$\epsilon : E^\vee \otimes E \to \mathcal{O}_X$ を Cohomology の Example
\ref{cohomology-example-dual-derived} で構成された写像とする．このとき写像
$$
E \xrightarrow{\eta \otimes 1} E \otimes E^\vee \otimes E
\xrightarrow{\sigma \otimes 1} E^\vee \otimes E \otimes E
\xrightarrow{1 \otimes \mu} E^\vee \otimes E
\xrightarrow{\epsilon} \mathcal{O}_X
$$
を考えられる．この写像は補題の主張にある写像の片側逆であると主張する．これを見るには，
合成 $\mathcal{O}_X \to \mathcal{O}_X$ が恒等写像であることを示せば十分である．
これは $X$ の生成点で（あるいは $f$ が同型となる $X$ の開部分スキーム上で）
確かめてよい．この場合 $E = \mathcal{O}_X$ であり，$\mu$ は通常の乗法写像なので，
結果は明らかである．
\end{proof}

\begin{lemma}
\label{obsolete-lemma-characterize-dbcoh-proper-regular}
$X$ を体 $k$ 上の正則な固有スキームとする．
$K \in \Ob(D_\QCoh(\mathcal{O}_X))$ とする．次は同値である：
\begin{enumerate}
\item $K \in D^b_{\textit{Coh}}(\mathcal{O}_X) = D_{perf}(\mathcal{O}_X)$；
\item $D(\mathcal{O}_X)$ のすべての完全対象 $E$ に対して
$\sum_{i \in \mathbf{Z}} \dim_k \Ext^i_X(E, K) < \infty$．
\end{enumerate}
\end{lemma}

\begin{proof}
(1) の等式は Derived Categories of Schemes の Lemma
\ref{perfect-lemma-perfect-on-regular} により成り立つ．含意 (1) $\Rightarrow$ (2) は
Derived Categories of Varieties の Lemma \ref{equiv-lemma-finiteness} から従う．
含意 (2) $\Rightarrow$ (1) は More on Morphisms の Lemma
\ref{more-morphisms-lemma-characterize-relatively-perfect} から従う．
\end{proof}

\begin{lemma}
\label{obsolete-lemma-bondal-van-den-bergh}
$X$ を体 $k$ 上の正則な固有スキームとする．
\begin{enumerate}
\item $F : D_{perf}(\mathcal{O}_X)^{opp} \to \text{Vect}_k$ を
$k$-線形コホモロジー函手であって，すべての $E \in D_{perf}(\mathcal{O}_X)$ に対し
$$
\sum\nolimits_{n \in \mathbf{Z}} \dim_k F(E[n]) < \infty
$$
を満たすものとする．このとき $F$ は，ある
$K \in D_{perf}(\mathcal{O}_X)$ に対する $E \mapsto \Hom_X(E, K)$ という形の函手と
同型である．
\item $G : D_{perf}(\mathcal{O}_X) \to \text{Vect}_k$ を
$k$-線形ホモロジー函手であって，すべての $E \in D_{perf}(\mathcal{O}_X)$ に対し
$$
\sum\nolimits_{n \in \mathbf{Z}} \dim_k G(E[n]) < \infty
$$
を満たすものとする．このとき $G$ は，ある
$K \in D_{perf}(\mathcal{O}_X)$ に対する $E \mapsto \Hom_X(K, E)$ という形の函手と
同型である．
\end{enumerate}
\end{lemma}

\begin{proof}
これは Derived Categories of Varieties の Theorem
\ref{equiv-theorem-bondal-van-den-bergh} と Lemma
\ref{equiv-lemma-homological-representable} から従う．以下では別の証明も与える．

\medskip\noindent
(1) の証明．導来圏 $D_\QCoh(\mathcal{O}_X)$ は直和をもち，コンパクト生成され，
$D_{perf}(\mathcal{O}_X)$ はコンパクト対象の充満部分圏である；
Derived Categories of Schemes の Lemma
\ref{perfect-lemma-quasi-coherence-direct-sums}，Theorem
\ref{perfect-theorem-bondal-van-den-Bergh}，Proposition
\ref{perfect-proposition-compact-is-perfect} を参照せよ．
Derived Categories of Varieties の Lemma \ref{equiv-lemma-van-den-bergh} により，
ある $K \in \Ob(D_\QCoh(\mathcal{O}_X))$ に対して
$F(E) = \Hom_X(E, K)$ と仮定してよい．すると Lemma
\ref{obsolete-lemma-characterize-dbcoh-proper-regular} により，
$K$ は $D^b_{\textit{Coh}}(\mathcal{O}_X)$ に属する．

\medskip\noindent
(2) の証明．$D_{perf}(\mathcal{O}_X)$ 上の反変函手
$E \mapsto E^\vee$ を考える；Cohomology の Lemma
\ref{cohomology-lemma-dual-perfect-complex} を参照せよ．この函手は
$D_{perf}(\mathcal{O}_X)$ の完全な反自己同値である．したがって函手
$F(E) = G(E^\vee)$ に (1) を適用して，
$G(E^\vee) = \Hom_X(E, K)$ となる $K \in D_{perf}(\mathcal{O}_X)$ を得る．
すると $G(E) = \Hom_X(E^\vee, K) = \Hom_X(K^\vee, E)$ であるから，
$K^\vee$ を取ればよい．
\end{proof}
% Frozen obsolete.tex lines 2130--2232.
\section{商の函手}
\label{obsolete-section-quotients}

\begin{lemma}
\label{obsolete-lemma-factors-through-quotient}
$S = \Spec(R)$ をアフィンスキームとする．$X$ を $S$ 上の代数空間とする．
$q_i : \mathcal{F} \to \mathcal{Q}_i$, $i = 1, 2$ を準連接
$\mathcal{O}_X$-加群の全射とする．$\mathcal{Q}_1$ は $S$ 上平坦であると仮定する．
$T \to S$ をスキームの準コンパクト射とし，次の分解が存在するとする：
$$
\xymatrix{
& \mathcal{F}_T \ar[rd]^{q_{2, T}} \ar[ld]_{q_{1, T}} \\
\mathcal{Q}_{1, T} & & \mathcal{Q}_{2, T} \ar@{..>}[ll]
}
$$
このとき閉部分スキーム $Z \subset S$ が存在して，(a) $T \to S$ は $Z$ を経由し，
(b) $q_{1, Z}$ は $q_{2, Z}$ を経由する．$\Ker(q_2)$ が有限型
$\mathcal{O}_X$-加群であり，$X$ が準コンパクトならば，$Z \to S$ を
有限表示に取ることができる．
\end{lemma}

\begin{proof}
写像 $\Ker(q_2) \to \mathcal{Q}_1$ に Flatness on Spaces の Lemma
\ref{spaces-flat-lemma-F-zero-somewhat-closed} を適用する．
\end{proof}







\section{空間と fpqc 被覆}
\label{obsolete-section-fpqc}

\noindent
代数空間が fpqc 被覆に関する層条件を満たすことを示す Gabber の議論により，
ここの内容は旧結果となった．Properties of Spaces の Section
\ref{spaces-properties-section-fpqc} を参照せよ．

\begin{lemma}
\label{obsolete-lemma-separated-fpqc}
$S$ をスキームとする．$X$ を $S$ 上の代数空間とする．
$\{f_i : T_i \to T\}_{i \in I}$ を $S$ 上のスキームの fpqc 被覆とする．このとき写像
$$
\Mor_S(T, X)
\longrightarrow
\prod\nolimits_{i \in I} \Mor_S(T_i, X)
$$
は単射である．
\end{lemma}

\begin{proof}
Properties of Spaces の Proposition
\ref{spaces-properties-proposition-sheaf-fpqc} から直ちに従う．
\end{proof}

\begin{lemma}
\label{obsolete-lemma-sheaf-fpqc-open-covering}
$S$ をスキームとする．$X$ を $S$ 上の代数空間とする．
$X = \bigcup_{j \in J} X_j$ を Zariski 被覆とする；Spaces の Definition
\ref{spaces-definition-Zariski-open-covering} を参照せよ．各 $X_j$ が fpqc 位相に関する
層性を満たすならば，$X$ も fpqc 位相に関する層性を満たす．
\end{lemma}

\begin{proof}
これは，すべての代数空間が fpqc 位相に関する層性を満たすことから正しい；
Properties of Spaces の Proposition
\ref{spaces-properties-proposition-sheaf-fpqc} を参照せよ．
\end{proof}

\begin{lemma}
\label{obsolete-lemma-sheaf-fpqc-quasi-separated}
$S$ をスキームとする．$X$ を $S$ 上の代数空間とする．
$X$ が $S$ 上 Zariski 局所的に準分離的ならば，$X$ は fpqc 位相に関する層条件を
満たす．
\end{lemma}

\begin{proof}
一般の Properties of Spaces の Proposition
\ref{spaces-properties-proposition-sheaf-fpqc} から直ちに従う．
\end{proof}

\begin{remark}
\label{obsolete-remark-proof-works-when}
この注では以前，Lemma \ref{obsolete-lemma-sheaf-fpqc-quasi-separated} の元の証明
（2009年12月18日）がどの程度一般化できるかを論じていた．
\end{remark}
% Frozen obsolete.tex lines 2233--2362.
\section{非常に reasonable な代数空間}
\label{obsolete-section-very-reasonable}

\noindent
やや旧結果となった内容を収める．

\begin{lemma}
\label{obsolete-lemma-reasonable-kolmogorov}
$S$ をスキームとする．$X$ を $S$ 上の reasonable な代数空間とする．このとき
$|X|$ は Kolmogorov である（Topology の Definition
\ref{topology-definition-generic-point} を参照せよ）．
\end{lemma}

\begin{proof}
定義と Decent Spaces の Lemma
\ref{decent-spaces-lemma-kolmogorov} から従う．
\end{proof}

\noindent
この節の残りでは，非常に reasonable な代数空間についていくつか注意を述べる．
スキーム $U$ と全射な \'etale 準コンパクト射 $U \to X$ が存在するならば，
$X$ は非常に reasonable である；Decent Spaces の Lemma
\ref{decent-spaces-lemma-characterize-very-reasonable} を参照せよ．

\begin{lemma}
\label{obsolete-lemma-scheme-very-reasonable}
スキームは非常に reasonable である．
\end{lemma}

\begin{proof}
恒等写像が準コンパクトかつ全射な \'etale 射なので，これは正しい．
\end{proof}

\begin{lemma}
\label{obsolete-lemma-very-reasonable-Zariski-local}
$S$ をスキームとする．$X$ を $S$ 上の代数空間とする．各 $X_i$ が非常に
reasonable となる Zariski 開被覆 $X = \bigcup X_i$ が存在するならば，
$X$ は非常に reasonable である．
\end{lemma}

\begin{proof}
これは Decent Spaces の Lemma
\ref{decent-spaces-lemma-properties-local} の場合 $(\epsilon)$ である．
\end{proof}

\begin{lemma}
\label{obsolete-lemma-quasi-separated-very-reasonable}
Zariski 局所的に準分離的な代数空間は非常に reasonable である．
とくに，任意の準分離的代数空間は非常に reasonable である．
\end{lemma}

\begin{proof}
これは Decent Spaces の Lemma
\ref{decent-spaces-lemma-bounded-fibres} の含意の一つである．
\end{proof}

\begin{lemma}
\label{obsolete-lemma-representable-very-reasonable}
$S$ をスキームとする．$X$, $Y$ を $S$ 上の代数空間とする．
$Y \to X$ を表現可能な射とする．$X$ が非常に reasonable ならば，
$Y$ も非常に reasonable である．
\end{lemma}

\begin{proof}
これは Decent Spaces の Lemma
\ref{decent-spaces-lemma-representable-properties} の場合 $(\epsilon)$ である．
\end{proof}

\begin{remark}
\label{obsolete-remark-very-reasonable-Zariski-locally-quasi-separated}
非常に reasonable な代数空間の集まりは，Zariski 局所的に準分離的な代数空間の
集まりより真に大きい．$G$ が非準分離的スキーム $U$ に不動点なしに作用する有限群で
あるような，$X = [U/G]$ という形の代数空間を考える（Spaces の Definition
\ref{spaces-definition-quotient} を参照せよ）．実際，この場合
$U \times_X U = U \times G$ であり，$U$ への二つの射影は明らかに準コンパクトなので，
$X$ は非常に reasonable である．他方，対角射
$U \times_X U \to U \times U$ は準コンパクトでなく，したがってこの代数空間は
準分離的でない．ここで，標数が $\not = 2$ の体 $k$ 上の無限次元アフィン空間で，
原点を二重化したものを $U$ とする；Schemes の Example
\ref{schemes-example-not-quasi-separated} を参照せよ．$U$ の二つの原点を
$0_1, 0_2$ とする．$G = \{+1, -1\}$ とし，$-1$ はすべての座標上では
$-1$ として作用し，$0_1$ と $0_2$ を入れ替えるものとする．すると
$[U/G]$ は非常に reasonable であるが，Zariski 局所的に準分離的でない
（詳細は省略する）．
\end{remark}

\noindent
注意：次の補題は慎重に用いるべきである．そこに現れるスキーム $U_i$ は，
分離的であるとも，さらには準分離的であるとも限らない．

\begin{lemma}
\label{obsolete-lemma-very-reasonable-quasi-compact-pieces}
$S$ をスキームとする．$X$ を $S$ 上の非常に reasonable な代数空間とする．
次を満たすスキーム $U_i$ の集合と射 $U_i \to X$ が存在する：
\begin{enumerate}
\item 各 $U_i$ は準コンパクトスキームである；
\item 各 $U_i \to X$ は \'etale である；
\item 二つの射影 $U_i \times_X U_i \to U_i$ はともに準コンパクトである；
\item 射 $\coprod U_i \to X$ は全射（かつ \'etale）である．
\end{enumerate}
\end{lemma}

\begin{proof}
Decent Spaces の Definition
\ref{decent-spaces-definition-very-reasonable} により，(2)，(3)，(4) を満たす
$U_i \to X$ が存在する．$i$ を固定し，$R_i = U_i \times_X U_i$ と置き，射影を
$s, t : R_i \to U_i$ と記す．任意のアフィン開集合 $W \subset U_i$ に対して，
開集合 $W' = t(s^{-1}(W)) \subset U_i$ は準コンパクトかつ $R_i$-不変である
（Groupoids の Lemma \ref{groupoids-lemma-constructing-invariant-opens} を参照せよ）．
したがって $W'$ は準コンパクトスキームで，$W' \to X$ は \'etale であり，
$W' \times_X W' = s^{-1}(W') = t^{-1}(W')$ なので二つの射影
$W' \times_X W' \to W'$ はともに準コンパクトである．ゆえに，$W \subset U_i$ が
$U_i$ のアフィン開被覆の要素を走るときの $W' \to X$ の族が求めるものを与える．
\end{proof}
% Frozen obsolete.tex lines 2363--2471.
\section{代数空間に関する旧結果となった補題}
\label{obsolete-section-obsolete-on-spaces}

\noindent
余分と思われるか，本文でもはや使われていない補題を挙げる．

\begin{lemma}
\label{obsolete-lemma-vanishing-surjective}
Cohomology of Spaces の Situation
\ref{spaces-cohomology-situation-vanishing} において，射
$p : X \to \Spec(A)$ は全射である．
\end{lemma}

\begin{proof}
この補題はもともと Cohomology of Spaces の Proposition
\ref{spaces-cohomology-proposition-vanishing-affine} の証明で用いられていたが，
現在ではその命題の帰結である．
\end{proof}

\begin{lemma}
\label{obsolete-lemma-vanishing-universally-closed}
Cohomology of Spaces の Situation
\ref{spaces-cohomology-situation-vanishing} において，射
$p : X \to \Spec(A)$ は普遍閉である．
\end{lemma}

\begin{proof}
この補題はもともと Cohomology of Spaces の Proposition
\ref{spaces-cohomology-proposition-vanishing-affine} の証明で用いられていたが，
現在ではその命題の帰結である．
\end{proof}

\begin{remark}
\label{obsolete-remark-equation-first}
このタグは以前，Formal Spaces の Lemma
\ref{formal-spaces-lemma-iff-adic-star} の証明中のある等式を指していた．
\end{remark}

\begin{remark}
\label{obsolete-remark-equation-second}
このタグは以前，Formal Spaces の Lemma
\ref{formal-spaces-lemma-iff-adic-star} の証明中のある等式を指していた．
\end{remark}






\section{代数スタックに関する旧結果となった補題}
\label{obsolete-section-obsolete-on-stacks}

\noindent
余分と思われるか，本文でもはや使われていない補題を挙げる．

\begin{lemma}
\label{obsolete-lemma-infinite-sequence}
$S$ を局所 Noether スキームとする．$\mathcal{X}$ を
$(\Sch/S)_{fppf}$ 上の亜群でファイバー付けられた圏で，(RS*) をもつものとする．
$x$ を，$S$ 上有限型のアフィンスキーム $U$ 上の $\mathcal{X}$ の対象とする．
$u_n \in U$, $n \geq 1$ を相異なる有限型点とし，すべての $n$ に対して
$x$ は $u_n$ で versal でないとする．$u_n$ を部分列で取り替えると，
$S$ 上の射
$$
x \to x_1 \to x_2 \to \ldots
\quad\text{in }\mathcal{X}\text{ lying over }\quad
U \to U_1 \to U_2 \to \ldots
$$
であって，次を満たすものが存在する：
\begin{enumerate}
\item 各 $n$ に対して射 $U \to U_n$ は1次の肥厚である；
\item 各 $n$ に対して短完全列
$$
0 \to \kappa(u_n) \to \mathcal{O}_{U_n} \to \mathcal{O}_{U_{n - 1}} \to 0
$$
があり，$n = 1$ に対して $U_0 = U$ とする；
\item 各 $n$ に対して，$u_n$ の開近傍 $W \subset U_n$ と射
$\alpha : x_n|_W \to x$ からなり，合成
$$
x|_{U \cap W} \xrightarrow{\text{restriction of }x \to x_n}
x_n|_W \xrightarrow{\alpha} x
$$
が標準射 $x|_{U \cap W} \to x$ となるような対 $(W, \alpha)$ は
{\bf 存在しない}．
\end{enumerate}
\end{lemma}

\begin{proof}
この補題はもともと versality の開性に関する判定法
（Artin's Axioms の Lemma \ref{artin-lemma-SGE-implies-openness-versality}）の証明で
用いられていたが，Artin's Axioms の Lemma
\ref{artin-lemma-infinite-sequence-pre} に置き換えられ，そこから容易に従う．
実際，$u_n$, $n \geq 1$ を部分列で取り替え，これらの点の間に特殊化がないと
仮定してよい；Properties の Lemma
\ref{properties-lemma-thin-infinite-sequence} を参照せよ．すると Artin's Axioms の
Lemma \ref{artin-lemma-infinite-sequence-pre} を適用して証明を終えられる．
\end{proof}
% Frozen obsolete.tex lines 2472--2663.
\section{スキームに対する余接複体の変種}
\label{obsolete-section-cotangent-schemes-variant}

\noindent
この節ではスキームの射の余接複体の別構成を与える．応用には Cotangent の Section
\ref{cotangent-section-cotangent-schemes-variant} で論じた，より容易な版で足りるので，
この節は現在，旧結果の章に置かれている．

\medskip\noindent
$f : X \to Y$ をスキームの射とする．$\mathcal{C}_{X/Y}$ を，スキームの可換図式
\begin{equation}
\label{obsolete-equation-object}
\vcenter{
\xymatrix{
X \ar[d] & U \ar[l] \ar[d] \ar[r]_i & A \ar[ld] \\
Y & V \ar[l]
}
}
\end{equation}
を対象とする圏とする．ここで
\begin{enumerate}
\item $U$ は $X$ の開部分スキームである；
\item $V$ は $Y$ の開部分スキームである；
\item $P$ を $\mathbf{Z}$ 上の（ある変数集合に関する）多項式代数として，
$V$ 上の同型 $A = V \times \Spec(P)$ が存在する．
\end{enumerate}
言い換えると，$A$ は $V$ 上の（無限次元）アフィン空間である．射は可換図式で与える．

\medskip\noindent
{\bf 記法．} $\mathcal{C}_{X/Y}$ の対象，すなわち図式
(\ref{obsolete-equation-object}) は，しばしば $U \to A$ と記す．このとき，(a) $U$ は
$X$ の開部分スキーム，(b) $U \to A$ は $Y$ 上の射，(c) 構造射
$A \to Y$ の像は開集合 $V \subset Y$，(d) $A \to V$ はアフィン空間であることを
了解する．$V \subset Y$ が $A \to Y$ の像であることを示すときは
$U \to A/V$ と書く．$X_{Zar}$ が $X$ の小 Zariski サイトを表すことを思い出す．
忘却函手
$$
\mathcal{C}_{X/Y} \to X_{Zar},\ (U \to A) \mapsto U
\quad\text{and}\quad
\mathcal{C}_{X/Y} \mapsto Y_{Zar},\ (U \to A/V) \mapsto V.
$$
がある．

\begin{lemma}
\label{obsolete-lemma-category-fibred}
$X \to Y$ をスキームの射とする．
\begin{enumerate}
\item 圏 $\mathcal{C}_{X/Y}$ は $X_{Zar}$ 上ファイバー付けられている．
\item 圏 $\mathcal{C}_{X/Y}$ は $Y_{Zar}$ 上ファイバー付けられている．
\item 圏 $\mathcal{C}_{X/Y}$ は，$U \subset X$, $V \subset Y$ が開で
$f(U) \subset V$ となる対 $(U, V)$ の圏上ファイバー付けられている．
\end{enumerate}
\end{lemma}

\begin{proof}
(1) について．$\mathcal{C}_{X/Y}$ の対象 $U \to A$ と $X_{Zar}$ の射
$U' \to U$ が与えられたとする．$i'$ を $U' \to U$ と $i$ の合成として，
$\mathcal{C}_{X/Y}$ の対象 $i' : U' \to A$ を考える．
$\mathcal{C}_{X/Y}$ の射 $(U' \to A) \to (U \to A)$ は，$X_{Zar}$ 上強 Cartesian
である．

\medskip\noindent
(2) について．対象 $U \to A/V$ と $V' \to V$ が与えられたとする．
$U' = U \cap f^{-1}(V')$ および $A' = V' \times_V A$ と置けば，
$V' \to V$ 上の強 Cartesian 射 $(U' \to A') \to (U \to A)$ を得る．

\medskip\noindent
(3) について．(3) の圏を $(X/Y)_{Zar}$ と記す．$U \to A/V$ と
$(X/Y)_{Zar}$ の射 $(U', V') \to (U, V)$ が与えられたとする．
$A' = V' \times_V A$ を考える．すると射
$(U' \to A'/V') \to (U \to A/V)$ は $(X/Y)_{Zar}$ 上の
$\mathcal{C}_{X/Y}$ において強 Cartesian である．
\end{proof}

\noindent
$X_{Zar}$ から継承した位相を用いて，$\mathcal{C}_{X/Y}$ 上の位相
$\tau_X$ を得る（Stacks の Section \ref{stacks-section-topology} を参照せよ）．
特に断らない限り，$\mathcal{C}_{X/Y}$ 上ではこの位相を考える．正確には，射の族
$\{(U_i \to A_i) \to (U \to A)\}$ が $\mathcal{C}_{X/Y}$ の被覆であることと，
次は同値である：
\begin{enumerate}
\item $U = \bigcup U_i$；
\item すべての $i$ に対して $A_i = A$．
\end{enumerate}
(2) で $A_i \cong A$ を許しても，同じ層の集まりを得る．函手 $u$ はトポスの射
$\pi : \Sh(\mathcal{C}_{X/Y}) \to \Sh(X_{Zar})$ を定める．

\medskip\noindent
サイト $\mathcal{C}_{X/Y}$ にはいくつかの環の層が備わっている．
\begin{enumerate}
\item 規則 $(U \to A) \mapsto \mathcal{O}(A)$ で与えられる層 $\mathcal{O}$．
\item 規則 $(U \to A) \mapsto \mathcal{O}(U)$ で与えられる層
$\underline{\mathcal{O}}_X = \pi^{-1}\mathcal{O}_X$．
\item 規則 $(U \to A/V) \mapsto \mathcal{O}(V)$ で与えられる層
$\underline{\mathcal{O}}_Y$．
\end{enumerate}
環付きトポスの射
\begin{equation}
\label{obsolete-equation-pi-schemes}
\vcenter{
\xymatrix{
(\Sh(\mathcal{C}_{X/Y}), \underline{\mathcal{O}}_X) \ar[r]_i \ar[d]_\pi &
(\Sh(\mathcal{C}_{X/Y}), \mathcal{O}) \\
(\Sh(X_{Zar}), \mathcal{O}_X)
}
}
\end{equation}
を得る．射 $i$ は台となるトポス上の恒等写像で，
$i^\sharp : \mathcal{O} \to \underline{\mathcal{O}}_X$ は明らかな写像である．
写像 $\pi$ は Cohomology on Sites の Situation
\ref{sites-cohomology-situation-fibred-category} の特別な場合である．以下では導来函手
$
Li^* : D(\mathcal{O}) \longrightarrow D(\underline{\mathcal{O}}_X)
$
が重要な役割を果たす．これは
$Ri_* = i_* : D(\underline{\mathcal{O}}_X) \to D(\mathcal{O})$ の左随伴である．また
$
L\pi_! : D(\underline{\mathcal{O}}_X) \longrightarrow D(\mathcal{O}_X)
$
も重要であり，これは
$\pi^* = \pi^{-1} : D(\mathcal{O}_X) \to D(\underline{\mathcal{O}}_X)$ の左随伴である．

\begin{remark}
\label{obsolete-remark-different-topologies}
$Y_{Zar}$ から継承した位相を取ることにより，$\mathcal{C}_{X/Y}$ 上の第2の位相
$\tau_Y$ を得る．第3の位相 $\tau_{X \to Y}$ もある．ここで射の族
$\{(U_i \to A_i) \to (U \to A)\}$ が被覆であることと，
$U = \bigcup U_i$, $V = \bigcup V_i$ かつ $A_i \cong V_i \times_V A$ であることは
同値である．これは，Lemma \ref{obsolete-lemma-category-fibred} の (3) にある対
$(U, V)$ の圏を台となる圏とするサイト $(X/Y)_{Zar}$ 上の位相から継承した位相である．
$(X/Y)_{Zar}$ の被覆は，$U = \bigcup U_i$ かつ $V = \bigcup V_i$ を満たす族
$\{(U_i, V_i) \to (U, V)\}$ である．トポスの射
$$
\xymatrix{
\Sh(\mathcal{C}_{X/Y})
= \Sh(\mathcal{C}_{X/Y}, \tau_X) &
\Sh(\mathcal{C}_{X/Y}, \tau_{X \to Y}) \ar[l] \ar[r] &
\Sh(\mathcal{C}_{X/Y}, \tau_Y)
}
$$
がある（$\tau_X$ が「既定」の位相であることを思い出す）．これらの矢印に対する
引き戻し函手は層化であり，順像は台となる前層上の恒等写像である．トポスの図式
$$
\xymatrix{
\Sh(X_{Zar}) \ar[d]^f & \Sh(\mathcal{C}_{X/Y}) \ar[l]^\pi &
\Sh(\mathcal{C}_{X/Y}, \tau_{X \to Y}) \ar[l] \ar[d] \\
\Sh(Y_{Zar}) & & \Sh(\mathcal{C}_{X/Y}, \tau_Y) \ar[ll]
}
$$
は {\bf 可換でない}．実際，$Y$ 上の零でないアーベル層の引き戻しは
$(\mathcal{C}_{X/Y}, \tau_{X \to Y})$ 上の零でないアーベル層であるが，
そのような層が $X$ 上では零へ引き戻される例を確かに見つけられる．
$Y_{Zar}$ 上の任意の前層 $\mathcal{F}$ は，$(U \to A/V)$ に集合
$\mathcal{F}(V)$ を割り当てる規則により，$\mathcal{C}_{Y/X}$ 上の層
$\underline{\mathcal{F}}$ を与えることに注意する．$\mathcal{F}$ が層であっても，
一般には $\underline{\mathcal{F}} = \pi^{-1}f^{-1}\mathcal{F}$ ではない．これは上の
図式の非可換性と関係する．実際，$\underline{\mathcal{F}}$ は，$\mathcal{F}$ を
下の水平矢印と右の垂直矢印を介して
$\Sh(\mathcal{C}_{X/Y}, \tau_{X \to Y})$ へ引き戻し，その後に順像を取ったものとして
記述できる．一例は層 $\underline{\mathcal{O}}_Y$ である．しかし正しいのは，写像
$\underline{\mathcal{F}} \to \pi^{-1}f^{-1}\mathcal{F}$ が存在し，それが
（後で見るように）$\pi_!$ を適用すると同型へ移されることである．
\end{remark}
% Frozen obsolete.tex lines 2664--2834.
\section{平坦加群の変形と障害}
\label{obsolete-section-modules}

\noindent
この節では，環 $\Lambda$ 上の任意の代数空間 $X$ に対する連接層のスタックの
変形理論の構成を概説する．Quot の Section \ref{quot-section-not-flat} における
議論が改良されたため，この内容は旧結果となっている．

\medskip\noindent
設定は次のとおりである．以下が与えられていると仮定する：
\begin{enumerate}
\item 環 $\Lambda$；
\item $\Lambda$ 上の代数空間 $X$；
\item $\Lambda$-代数 $A$；
$X_A = X \times_{\Spec(\Lambda)} \Spec(A)$ と置く；
\item $A$ 上平坦な有限表示 $\mathcal{O}_{X_A}$-加群 $\mathcal{F}$．
\end{enumerate}
この状況で，全射
$$
0 \to I \to A' \to A \to 0
$$
で，$A'$ が $\Lambda$-代数であり，その核 $I$ が $A'$ の平方零イデアルであるものを
すべて考える．$A'$ が与えられると，$\Spec(\Lambda)$ 上の代数空間の一次の厚化
$X_A \to X_{A'}$ を得る．その各々について，$\mathcal{F}$ を $X_{A'}$ 上の
$A'$ 上平坦な有限表示加群 $\mathcal{F}'$ へ持ち上げる問題を考える．この設定で
Deformation Theory, Lemma \ref{defos-lemma-flat-ringed-topoi} の結果を再現したい．

\medskip\noindent
より正確には，$\textit{Lift}(\mathcal{F}, A')$ を対 $(\mathcal{F}', \alpha)$ の
圏とする．ここで $\mathcal{F}'$ は $X_{A'}$ 上の $A'$ 上平坦な有限表示加群であり，
$\alpha : \mathcal{F}'|_{X_A} \to \mathcal{F}$ は同型である．射
$(\mathcal{F}'_1, \alpha_1) \to (\mathcal{F}'_2, \alpha_2)$ は，
$\alpha_1$ および $\alpha_2$ と両立する同型
$\mathcal{F}'_1 \to \mathcal{F}'_2$ である．
$\textit{Lift}(\mathcal{F}, A')$ の同型類の集合を
$\text{Lift}(\mathcal{F}, A')$ と記す．

\medskip\noindent
$\mathcal{G}$ を，$X_\etale$ 上の $A$ 上平坦な
$\mathcal{O}_X \otimes_\Lambda A$-加群の層とする．圏
$\textit{Lift}(\mathcal{G}, A')$ を導入する．その対象は対
$(\mathcal{G}', \beta)$ であり，$\mathcal{G}'$ は
$A'$ 上平坦な $\mathcal{O}_X \otimes_\Lambda A'$-加群の層，$\beta$ は同型
$\mathcal{G}' \otimes_{A'} A \to \mathcal{G}$ である．

\begin{lemma}
\label{obsolete-lemma-equivalence}
記法と仮定は上のとおりとする．$p : X_A \to X$ を射影とする．$A'$ が
与えられたとき，$p' : X_{A'} \to X$ を射影と記す．函手 $p'_*$ は，
次の圏の間の圏同値を誘導する：
\begin{enumerate}
\item 圏 $\textit{Lift}(\mathcal{F}, A')$；
\item 圏 $\textit{Lift}(p_*\mathcal{F}, A')$．
\end{enumerate}
\end{lemma}

\begin{proof}
FIXME.
\end{proof}

\noindent
$\mathcal{H}$ を，$\mathcal{C}_{X/\Lambda}$ 上の $A$ 上平坦な
$\mathcal{O} \otimes_\Lambda A$-加群の層とする．圏
$\textit{Lift}_\mathcal{O}(\mathcal{H}, A')$ を導入する．その対象は対
$(\mathcal{H}', \gamma)$ であり，$\mathcal{H}'$ は
$A'$ 上平坦な $\mathcal{O} \otimes_\Lambda A'$-加群の層，
$\gamma : \mathcal{H}' \otimes_A A' \to \mathcal{H}$ は
$\mathcal{O} \otimes_\Lambda A$-加群の同型である．

\medskip\noindent
$\mathcal{G}$ を，$X_\etale$ 上の $A$ 上平坦な
$\mathcal{O}_X \otimes_\Lambda A$-加群の層とする．
Cotangent, Equation (\ref{cotangent-equation-pi-spaces}) の射 $i$ と $\pi$ を考える．
$\underline{\mathcal{G}} = \pi^{-1}(\mathcal{G})$ と記す．これは単に規則
$(U \to \mathbf{A}) \mapsto \mathcal{G}(U)$ で与えられるので，
$\underline{\mathcal{O}}_X \otimes_\Lambda A$-加群の層である．
$i_*\underline{\mathcal{G}}$ を，同じ層を
$\mathcal{O} \otimes_\Lambda A$-加群の層とみなしたものとする．

\begin{lemma}
\label{obsolete-lemma-second-equivalence}
記法と仮定は上のとおりとする．函手 $\pi_!$ は，次の圏の間の圏同値を誘導する：
\begin{enumerate}
\item 圏 $\textit{Lift}_\mathcal{O}(i_*\underline{\mathcal{G}}, A')$；
\item 圏 $\textit{Lift}(\mathcal{G}, A')$．
\end{enumerate}
\end{lemma}

\begin{proof}
FIXME.
\end{proof}

\begin{lemma}
\label{obsolete-lemma-second-equivalence-obs}
記法と仮定は Lemma \ref{obsolete-lemma-second-equivalence} のとおりとする．
$D(\mathcal{O}_X \otimes_\Lambda A)$ の対象
$$
L = L(\Lambda, X, A, \mathcal{G}) = L\pi_!(Li^*(i_*(\underline{\mathcal{G}})))
$$
を考える．平方零の核 $I$ をもつ $\Lambda$-代数の全射 $A' \to A$ が与えられると，
次が成り立つ：
\begin{enumerate}
\item 圏 $\textit{Lift}(\mathcal{G}, A')$ が空でないための必要十分条件は，ある類
$\xi \in \Ext^2_{\mathcal{O}_X \otimes A}(L, \mathcal{G} \otimes_A I)$
が零であることである．
\item $\textit{Lift}(\mathcal{G}, A')$ が空でなければ，
$\text{Lift}(\mathcal{G}, A')$ は
$\Ext^1_{\mathcal{O}_X \otimes A}(L, \mathcal{G} \otimes_A I)$
の作用に関して主同次空間をなす．
\item 持ち上げ $\mathcal{G}'$ が与えられたとき，$\text{id}_\mathcal{G}$ へ
引き戻される $\mathcal{G}'$ の自己同型の集合は，標準的に
$\Ext^0_{\mathcal{O}_X \otimes A}(L, \mathcal{G} \otimes_A I)$
と同型である．
\end{enumerate}
\end{lemma}

\begin{proof}
FIXME.
\end{proof}

\noindent
最後に，以上を次のようにまとめる．

\begin{proposition}
\label{obsolete-proposition-conclusion}
$\Lambda$, $X$, $A$, $\mathcal{F}$ は上のとおりとする．$D(X_A)$ の標準的な対象
$L = L(\Lambda, X, A, \mathcal{F})$ が存在し，平方零の核 $I$ をもつ
$\Lambda$-代数の全射 $A' \to A$ が与えられると，次が成り立つ：
\begin{enumerate}
\item 圏 $\textit{Lift}(\mathcal{F}, A')$ が空でないための必要十分条件は，ある類
$\xi \in \Ext^2_{X_A}(L,
\mathcal{F} \otimes_A I)$ が零であることである．
\item $\textit{Lift}(\mathcal{F}, A')$ が空でなければ，
$\text{Lift}(\mathcal{F}, A')$ は
$\Ext^1_{X_A}(L, \mathcal{F} \otimes_A I)$ の作用に関して主同次空間をなす．
\item 持ち上げ $\mathcal{F}'$ が与えられたとき，$\text{id}_\mathcal{F}$ へ
引き戻される $\mathcal{F}'$ の自己同型の集合は，標準的に
$\Ext^0_{X_A}(L, \mathcal{F} \otimes_A I)$ と同型である．
\end{enumerate}
\end{proposition}

\begin{proof}
FIXME.
\end{proof}

\begin{lemma}
\label{obsolete-lemma-pseudo-coherent}
Proposition \ref{obsolete-proposition-conclusion} の状況で，$X \to \Spec(\Lambda)$ が
局所有限型で $\Lambda$ が Noether 環ならば，$L$ は擬連接である．
\end{lemma}

\begin{proof}
FIXME.
\end{proof}
% Frozen obsolete.tex lines 2835--3010.
\section{非平坦な場合の連接層のスタック}
\label{obsolete-section-not-flat}

\noindent
Quot, Theorem \ref{quot-theorem-coherent-algebraic} では，
$f : X \to B$ が平坦であるという仮定は必要ない．この節では，導来代数幾何の
考えに基づいて証明法を修正し，平坦性の仮定を回避する．Quot, Section
\ref{quot-section-not-flat} では，まったく異なる方法により正確に同じ結果を得ている．
このため，本節の方法は旧結果となっている．

\medskip\noindent
Quot, Theorem \ref{quot-theorem-coherent-algebraic} の証明で平坦性を用いる唯一の
段階は，Quot, Lemma \ref{quot-lemma-coherent-defo-thy} の適用である．この補題は，
Artin's Axioms, Section \ref{artin-section-dual} の意味での障害理論を構成するために
用いられる．補題の証明は，Deformation Theory, Lemmas
\ref{defos-lemma-flat-ringed-topoi} および
\ref{defos-lemma-verify-iv-ringed-topoi} に依存する．これらは Deformation Theory,
Section \ref{defos-section-flat-ringed-topoi} にある．これが $f$ の平坦性という
仮定が現れる理由である．先へ進む前に，Deformation Theory, Lemmas
\ref{defos-lemma-flat-ringed-topoi} の結果 (2) と (3) は，$f$ の平坦性を仮定しなくても
成り立つことに注意する．実際，これらが依存する Deformation Theory, Lemmas
\ref{defos-lemma-inf-ext-rel-ringed-topoi} および
\ref{defos-lemma-inf-map-rel-ringed-topoi} には平坦性の仮定がない．

\medskip\noindent
詳細を述べる前に，以下を明確にすると思われるので，導来代数幾何から得られる構成の
動機を説明する．$A$ を局所 Noether 基底 $S$ 上の有限型代数とする．$X$ の $A$ への
「導来基底変換」を $X \otimes^\mathbf{L} A$ と記し，標準的な包含射を
$i : X_A \to X \otimes^\mathbf{L} A$ と記す．対象
$X \otimes^\mathbf{L} A$ は Stacks project では（まだ）定義されていない．これは，
代数空間 $X_A$ に環の単体的層 $\mathcal{O}_{X \otimes^\mathbf{L} A}$ を備えたものと
考えられ，そのホモロジー層は
$$
H_i(\mathcal{O}_{X \otimes^\mathbf{L} A}) =
\text{Tor}^{\mathcal{O}_S}_i(\mathcal{O}_X, A).
$$
である．射 $X \otimes^\mathbf{L} A \to \Spec(A)$ は平坦である（環の単体的層の
各項が $A$-平坦だからである）．したがって，平坦加群の変形に関する通常の内容を
これに適用できる．ゆえに，$i = 0, 1, 2$ に対し，それぞれ無限小自己同型，
無限小変形，障害を与える群
$$
\Ext^i_{X \otimes^\mathbf{L} A}(i_*\mathcal{F},
i_*\mathcal{F} \otimes_A M)
$$
を用いて障害理論を得ることが分かる．$i_*\mathcal{F}$ の
$X \otimes^\mathbf{L} A'$ への平坦変形は，自動的に $i'_*\mathcal{F}'$ の形である．
ここで $\mathcal{F}'$ は $\mathcal{F}$ の平坦変形である．函手 $Li^*$ と
$i_* = Ri_*$ の随伴により，これらの Ext 群は
$$
\Ext^i_{X_A}(Li^*(i_*\mathcal{F}), \mathcal{F} \otimes_A M)
$$
に等しい．したがって Quot, Lemma \ref{quot-lemma-coherent-defo-thy} の証明と
まったく同じ形の障害群が得られる．唯一の変更は，$\mathcal{F}$ の最初の出現を
複体 $Li^*(i_*\mathcal{F})$ で置き換えることである．

\medskip\noindent
以下では，$E(\mathcal{F}) = Li^*(i_*\mathcal{F})$ を「直接」構成し，
$\mathcal{F}$ の変形理論との関係を直接証明することにより，この補題の非平坦版を
証明する．実際，関心があるのは $i = 0, 1, 2$ に対する Ext 群
$\Ext^i_{X_A}(Li^*(i_*\mathcal{F}), \mathcal{F} \otimes_A M)$
だけなので，$\tau_{\geq -2}E(\mathcal{F})$ を構成すれば十分である．さらに，
コホモロジー層を
$$
H^i(E(\mathcal{F})) =
\left\{
\begin{matrix}
0 & \text{if }i > 0 \\
\mathcal{F} & \text{if } i = 0 \\
0 & \text{if } i = -1 \\
\text{Tor}_1^{\mathcal{O}_S}(\mathcal{O}_X, A)
\otimes_{\mathcal{O}_X} \mathcal{F} &
\text{if } i = -2
\end{matrix}
\right.
$$
と同定することさえできる．この観察が，以下の注意における
$E(\mathcal{F})$ の構成を導く．

\begin{remark}[直接構成]
\label{obsolete-remark-construction-E}
$S$ をスキームとする．$f : X \to B$ を $S$ 上の代数空間の射とする．$U$ を
$B$ 上の別の代数空間とする．第2射影を $q : X \times_B U \to U$ と記す．
Cotangent, Section \ref{cotangent-section-fibre-product} の特別三角形
$$
Lq^*L_{U/B} \to L_{X \times_B U/B} \to E \to Lq^*L_{U/B}[1]
$$
を考える．任意の $\mathcal{O}_{X \times_B U}$-加群の層 $\mathcal{F}$ に対し，
Atiyah 類
$$
\mathcal{F} \to
L_{X \times_B U/B}
\otimes_{\mathcal{O}_{X \times_B U}}^\mathbf{L} \mathcal{F}[1]
$$
がある．Cotangent, Section \ref{cotangent-section-atiyah-general} を参照せよ．
これを $E$ への写像と合成し，$D(\mathcal{O}_{X \times_B U})$ における特別三角形
$$
E(\mathcal{F}) \to \mathcal{F} \to
\mathcal{F} \otimes_{\mathcal{O}_{X \times_B U}}^\mathbf{L} E[1] \to
E(\mathcal{F})[1]
$$
を選べる．構成により，Atiyah 類は写像
$$
e_\mathcal{F} :
E(\mathcal{F})
\longrightarrow
Lq^*L_{U/B} \otimes_{\mathcal{O}_{X \times_B U}}^\mathbf{L} \mathcal{F}[1]
$$
へ持ち上がり，これは特別三角形の射
$$
\xymatrix{
\mathcal{F} \otimes^\mathbf{L} Lq^*L_{U/B}[1] \ar[r] &
\mathcal{F} \otimes^\mathbf{L} L_{X \times_B U/B}[1] \ar[r] &
\mathcal{F} \otimes^\mathbf{L} E[1] \\
E(\mathcal{F}) \ar[r] \ar[u]^{e_\mathcal{F}} &
\mathcal{F} \ar[r] \ar[u]^{Atiyah} &
\mathcal{F} \otimes^\mathbf{L} E[1] \ar[u]^{=}
}
$$
に組み込まれる．$S, B, X, f, U, \mathcal{F}$ が与えられるたびに，
$E(\mathcal{F})$ と $e_\mathcal{F}$ の選択を一つ固定する．
\end{remark}

\begin{remark}[障害類の構成]
\label{obsolete-remark-construction-ob}
記法を Remark \ref{obsolete-remark-construction-E} のとおりとし，$i : U \to U'$ を
$B$ 上の $U$ の一次の厚化とする．$\mathcal{I} \subset \mathcal{O}_{U'}$ を，
$B'$ における $B$ を切り出す準連接イデアル層とする．基本三角形
$$
Li^*L_{U'/B} \to L_{U/B} \to L_{U/U'} \to Li^*L_{U'/B}[1]
$$
と写像 $L_{U/U'} \to \mathcal{I}[1]$ は，写像
$e_{U'} : L_{U/B} \to \mathcal{I}[1]$ を定める．これと前の注意の写像
$e_\mathcal{F}$ を組み合わせると，
$$
(\text{id}_\mathcal{F} \otimes Lq^*e_{U'}) \cup e_\mathcal{F} :
E(\mathcal{F})
\longrightarrow
\mathcal{F} \otimes_{\mathcal{O}_{X \times_B U}} q^*\mathcal{I}[2]
$$
を得る（導来テンソル積から通常のテンソル積への写像とも合成した）．言い換えると，
元
$$
\xi_{U'} \in
\Ext^2_{\mathcal{O}_{X \times_B U}}(
E(\mathcal{F}),
\mathcal{F} \otimes_{\mathcal{O}_{X \times_B U}} q^*\mathcal{I})
$$
を得る．
\end{remark}

\begin{lemma}
\label{obsolete-lemma-ob-is-obstruction}
Remark \ref{obsolete-remark-construction-ob} の状況で，$\mathcal{F}$ が $U$ 上平坦であると
仮定する．類 $\xi_{U'}$ が消えることは，$U'$ 上平坦で
$i^*\mathcal{F}' \cong \mathcal{F}$ を満たす
$\mathcal{O}_{X \times_B U'}$-加群 $\mathcal{F}'$ が存在するための必要十分条件である．
\end{lemma}
% Frozen obsolete.tex lines 3011--3193.

\begin{proof}[証明の概略]
Deformation Theory, Lemma \ref{defos-lemma-inf-obs-ext-rel-ringed-topoi} の判定法を
用いる．$\mathcal{O} = \mathcal{O}_{X \times_B U}$ および
$\mathcal{O}' = \mathcal{O}_{X \times_B U'}$ と略記する．短完全列
$$
0 \to \mathcal{I} \to \mathcal{O}_{U'} \to \mathcal{O}_U \to 0.
$$
を考える．$\mathcal{J} \subset \mathcal{O}'$ を，$X \times_B U$ を切り出す
準連接イデアル層とする．上のことから完全列
$$
\text{Tor}_1^{\mathcal{O}_B}(\mathcal{O}_X, \mathcal{O}_U) \to
q^*\mathcal{I} \to \mathcal{J} \to 0
$$
を得る．ここで $\text{Tor}_1^{\mathcal{O}_B}(\mathcal{O}_X, \mathcal{O}_U)$ は
$$
\text{Tor}_1^{h^{-1}\mathcal{O}_B}(p^{-1}\mathcal{O}_X, q^{-1}\mathcal{O}_U)
\otimes_{(p^{-1}\mathcal{O}_X\otimes_{h^{-1}\mathcal{O}_B}q^{-1}\mathcal{O}_U)}
\mathcal{O}.
$$
の略記である．$\mathcal{F}$ とテンソルを取ると，完全列
$$
\mathcal{F} \otimes_\mathcal{O}
\text{Tor}_1^{\mathcal{O}_B}(\mathcal{O}_X, \mathcal{O}_U) \to
\mathcal{F} \otimes_\mathcal{O}
q^*\mathcal{I} \to
\mathcal{F} \otimes_\mathcal{O} \mathcal{J} \to 0
$$
を得る（文字 $\mathcal{I}$ と $\mathcal{J}$ の役割は，Deformation Theory,
Lemma \ref{defos-lemma-inf-obs-ext-rel-ringed-topoi} の記法とは逆であることに
注意せよ）．この補題の条件 (1) は，上の最後の写像が同型であること，すなわち
最初の写像が零であることである．この写像が消えるかどうかは，幾何学的点
$\overline{z} = (\overline{x}, \overline{u}) : \Spec(k) \to X \times_B U$ の
茎で確認できる．$R = \mathcal{O}_{B, \overline{b}}$,
$A = \mathcal{O}_{X, \overline{x}}$, $B = \mathcal{O}_{U, \overline{u}}$，および
$C = \mathcal{O}_{\overline{z}}$ と置く．Cotangent, Lemma
\ref{cotangent-lemma-fibre-product} と $E(\mathcal{F})$ の定義三角形から，
$$
H^{-2}(E(\mathcal{F}))_{\overline{z}} =
\mathcal{F}_{\overline{z}} \otimes \text{Tor}_1^R(A, B)
$$
が分かる．したがって写像 $\xi_{U'}$ は写像
$$
\mathcal{F}_{\overline{z}} \otimes \text{Tor}_1^R(A, B)
\longrightarrow
\mathcal{F}_{\overline{z}} \otimes_B \mathcal{I}_{\overline{u}}
$$
を誘導する．この写像は上で述べた写像の茎と同じであると主張する
（証明は省略する．これは純粋に環論的な主張である）．したがって，Deformation
Theory, Lemma \ref{defos-lemma-inf-obs-ext-rel-ringed-topoi} の条件 (1) は，
写像
$H^{-2}(\xi_{U'}) :
H^{-2}(E(\mathcal{F})) \to \mathcal{F} \otimes \mathcal{I}$
が消えることと同値である．

\medskip\noindent
証明を終えるため，条件 (1) が満たされると仮定して，条件 (2) が
$\xi_{U'}$ の消滅と同値であることを示す．証明の残りでは，
$\mathcal{F} \otimes \mathcal{I}$ を
$\mathcal{F} \otimes_\mathcal{O} q^*\mathcal{I} =
\mathcal{F} \otimes_\mathcal{O} \mathcal{J}$ の意味で用いる．スペクトル系列
$$
\Ext^i(H^{-j}(E(\mathcal{F})), \mathcal{F} \otimes \mathcal{I})
\Rightarrow
\Ext^{i + j}(E(\mathcal{F}), \mathcal{F} \otimes \mathcal{I})
$$
を考え，$H^0(E(\mathcal{F})) = \mathcal{F}$ および
$H^{-1}(E(\mathcal{F})) = 0$ を用いると，完全列
$$
0 \to
\Ext^2(\mathcal{F}, \mathcal{F} \otimes \mathcal{I}) \to
\Ext^2(E(\mathcal{F}), \mathcal{F} \otimes \mathcal{I}) \to
\Hom(H^{-2}(E(\mathcal{F})), \mathcal{F} \otimes \mathcal{I})
$$
があることが分かる．したがって，元 $\xi_{U'}$ は
$\Ext^2(\mathcal{F}, \mathcal{F} \otimes \mathcal{I})$ の元である．
この元が Deformation Theory, Lemma
\ref{defos-lemma-inf-obs-ext-rel-ringed-topoi} の元と一致することを示せば証明は
完了するが，その確認は省略する．
\end{proof}

\begin{lemma}
\label{obsolete-lemma-coherent-defo-thy-general}
Quot, Situation \ref{quot-situation-coherent} において，$S$ は局所 Noether
スキームで $S = B$ であると仮定する．
$\mathcal{X} = \textit{Coh}_{X/B}$ と置く．このとき $\mathcal{X}$ に対して
versal 性の開性が成り立つ（Artin's Axioms, Definition
\ref{artin-definition-openness-versality} を参照せよ）．
\end{lemma}

\begin{proof}[証明の概略]
$U \to S$ をスキームの有限型射，$x$ を $U$ 上の $\mathcal{X}$ の対象，
$u_0 \in U$ を有限型の点で，$x$ が $u_0$ で versal であるものとする．$U$ を
縮小した後，$u_0$ は閉点であると仮定でき（Morphisms, Lemma
\ref{morphisms-lemma-point-finite-type}），また $U = \Spec(A)$ で，$U \to S$ の
像は $S$ のアフィン開集合 $\Spec(\Lambda)$ に含まれると仮定できる．
Artin's Axioms, Lemma \ref{artin-lemma-dual-openness} を用いて補題を証明する．
与えられた対象 $x$ に対応する，$A$ 上平坦な $X_A = \Spec(A) \times_S X$ 上の
連接加群を $\mathcal{F}$ とする．

\medskip\noindent
Remark \ref{obsolete-remark-construction-E} のように $E(\mathcal{F})$ と
$e_\mathcal{F}$ を選ぶ．$E(\mathcal{F})$ のコホモロジー層の記述から，任意の
$A$-加群 $M$ に対して
$$
\Ext^1(E(\mathcal{F}), \mathcal{F} \otimes_A M) =
\Ext^1(\mathcal{F}, \mathcal{F} \otimes_A M)
$$
が分かる．これと Deformation Theory, Lemma
\ref{defos-lemma-inf-ext-rel-ringed-topoi} を用いると，函手の同型
$$
T_x(M) = \Ext^1_{X_A}(E(\mathcal{F}), \mathcal{F} \otimes_A M)
$$
を得る．Lemma \ref{obsolete-lemma-ob-is-obstruction} により，平方零の核 $I$ をもつ
$\Lambda$-代数の任意の全射 $A' \to A$ に対して，障害類
$$
\xi_{A'} \in \Ext^2_{X_A}(E(\mathcal{F}), \mathcal{F} \otimes_A I)
$$
がある．$m = 2$ とし，$i \leq m$ に対する Ext 群
$\Ext^i_{X_A}(E(\mathcal{F}), \mathcal{F} \otimes_A M)$ の計算に
Derived Categories of Spaces, Lemma \ref{spaces-perfect-lemma-compute-ext} を
適用する．$E(\mathcal{F})$ が $D^-_{\textit{Coh}}$ に属することの確認は省略する．
ヒント：Cotangent, Lemma \ref{cotangent-lemma-cotangent-finite} を用いよ．
完全対象 $K \in D(A)$ と，境界写像と両立する函手的同型
$$
H^i(K \otimes_A^\mathbf{L} M)
\longrightarrow
\Ext^i_{X_A}(E(\mathcal{F}), \mathcal{F} \otimes_A M)
$$
が $i \leq m$ に対して得られる．この対象 $K$ と上の同定により，Artin's Axioms,
Situation \ref{artin-situation-dual} にあるデータが得られる．最後に，Artin's Axioms,
Lemma \ref{artin-lemma-dual-obstruction} の条件 (iv) は，Deformation Theory,
Lemma \ref{defos-lemma-verify-iv-ringed-topoi} の変種から従う．その定式化と証明は
省略する．したがって Artin's Axioms, Lemma
\ref{artin-lemma-dual-openness} を適用でき，補題が証明される．
\end{proof}

\begin{theorem}
\label{obsolete-theorem-coherent-algebraic-general}
$S$ をスキームとする．$f : X \to B$ を $S$ 上の代数空間の射とする．$f$ は
有限表示かつ分離的であると仮定する．このとき
$\textit{Coh}_{X/B}$ は $S$ 上の代数スタックである．
\end{theorem}

\begin{proof}
この定理は Quot, Theorem \ref{quot-theorem-coherent-algebraic-general} の複製である．
ここに複製を置く理由は，本節の内容から第2の証明が得られるためである
（本節の冒頭で論じたとおりである）．すなわち，Quot, Theorem
\ref{quot-theorem-coherent-algebraic} の証明とまったく同じように論じるが，
Lemma \ref{obsolete-lemma-coherent-defo-thy-general} を，そこで用いられた Quot, Lemma
\ref{quot-lemma-coherent-defo-thy} の代わりに用いる．
\end{proof}
% Frozen obsolete.tex lines 3194--3270.
\section{修正}
\label{obsolete-section-modifications}

\noindent
ここに，Algebraization of Formal Spaces, Equation
(\ref{restricted-equation-modification}) の圏に関する旧結果を記す．
現在の内容については Algebraization of Formal Spaces, Section
\ref{restricted-section-modifications} を参照されたい．

\begin{lemma}
\label{obsolete-lemma-henselian}
$(A, \mathfrak m, \kappa)$ を Noether 局所環とする．$A$ に対する
Algebraization of Formal Spaces, Equation
(\ref{restricted-equation-modification}) の圏は，$A$ の Hensel 化 $A^h$ に対する
Algebraization of Formal Spaces, Equation
(\ref{restricted-equation-modification}) の圏と同値である．
\end{lemma}

\begin{proof}
これは Algebraization of Formal Spaces, Lemma
\ref{restricted-lemma-equivalence-to-completion} の特別な場合である．
\end{proof}

\noindent
有理特異点に関する次の補題は，曲面特異点の解消を扱う章ではもはや必要ない．

\begin{lemma}
\label{obsolete-lemma-double-dual-rational}
Resolution of Surfaces, Situation \ref{resolve-situation-rational} の状況とする．
$M$ を有限反射的 $A$-加群とする．付随する $\mathcal{O}_S$-加群の引き戻しを
$M \otimes_A \mathcal{O}_X$ と記す．このとき
$M \otimes_A \mathcal{O}_X$ からその二重双対への写像は全射である．
\end{lemma}

\begin{proof}
二重双対を $\mathcal{F} = (M \otimes_A \mathcal{O}_X)^{**}$ とし，評価写像
$M \otimes_A \mathcal{O}_X \to \mathcal{F}$ の像を
$\mathcal{F}' \subset \mathcal{F}$ とする．すると短完全列
$$
0 \to \mathcal{F}' \to \mathcal{F} \to \mathcal{Q} \to 0
$$
がある．$X$ は正規なので，余次元 $1$ の点における局所環
$\mathcal{O}_{X, x}$ は離散付値環である（Properties, Lemma
\ref{properties-lemma-criterion-normal} を参照せよ）．したがって，そのような点では
More on Algebra, Lemma \ref{more-algebra-lemma-cokernel-map-double-dual-dvr} により
$\mathcal{Q}_x = 0$ である．ゆえに $\mathcal{Q}$ は有限個の閉点に台をもち，
Cohomology of Schemes, Lemma
\ref{coherent-lemma-coherent-support-dimension-0} により大域切断で生成される．
$\mathcal{F}'$ は大域切断で生成されるので（Resolution of Surfaces, Lemma
\ref{resolve-lemma-globally-generated}），完全列
$$
0 \to H^0(X, \mathcal{F}') \to H^0(X, \mathcal{F}) \to H^0(X, \mathcal{Q}) \to 0
$$
を得る．$X \to \Spec(A)$ は閉点の補集合上で同型であり，$M$ は反射的なので，
写像
$$
M \to H^0(X, \mathcal{F}') \to H^0(X, \mathcal{F})
$$
は $A$ の任意の非極大素イデアルで局所化した後に同型を誘導する．したがって，
More on Algebra, Lemma
\ref{more-algebra-lemma-check-isomorphism-via-depth-and-ass}，および正規環上の反射的
加群が性質 $(S_2)$ をもつという事実（More on Algebra, Lemma
\ref{more-algebra-lemma-reflexive-over-normal}）により，これらの写像は同型である．
よって望みどおり $\mathcal{Q} = 0$ と結論できる．
\end{proof}
% Frozen obsolete.tex lines 3271--3363.
\section{交叉理論}
\label{obsolete-section-intersection-theory}

\begin{lemma}
\label{obsolete-lemma-good-blowing-up}
$b : X' \to X$ を，体 $k$ 上の滑らかな射影スキームを滑らかな閉部分スキーム
$Z \subset X$ に沿って吹き上げたものとする．図式は
$$
\xymatrix{
E \ar[r]_j \ar[d]_\pi & X' \ar[d]^b \\
Z \ar[r]^i & X
}
$$
である．$K_0(X)$ の元で，$Z$ への制限が $K_0(Z)$ における
$\mathcal{C}_{Z/X}$ の類に等しいものが存在すると仮定する．このとき，ある
$\alpha'' \in K_0(X')$ に対して，$K_0(X')$ において
$[Lb^*\mathcal{O}_Z] = [\mathcal{O}_E] \cdot \alpha''$ が成り立つ．
\end{lemma}

\begin{proof}
スキーム $X$, $X'$, $E$, $Z$ は $k$ 上滑らかかつ射影的であり，したがって
$K'_0(X) = K_0(X) = K_0(\textit{Vect}(X)) =
K_0(D^b_{\textit{Coh}}(X)))$ である．他の $3$ つについても同様である．
Derived Categories of Schemes, Lemmas \ref{perfect-lemma-Noetherian-Kprime},
\ref{perfect-lemma-Kprime-K}, および \ref{perfect-lemma-K-is-old-K} を参照せよ．
この証明では，これらの版を自由に行き来する．$\mathcal{F}$ を定める，有限局所自由
$\mathcal{O}_E$-加群の短完全列
$$
0 \to \mathcal{F} \to \pi^*\mathcal{C}_{Z/X} \to \mathcal{C}_{E/X'} \to 0
$$
を考える．$\mathcal{C}_{E/X'} = \mathcal{O}_{X'}(-E)|_E$ は，可逆
$\mathcal{O}_X$-加群 $\mathcal{O}_{X'}(-E)$ の制限であることに注意する．
$K_0(Z)$ において $i^*\alpha = [\mathcal{C}_{Z/X}]$ を満たす元
$\alpha \in K_0(X)$ を取る．
$\alpha' = b^*\alpha - [\mathcal{O}_{X'}(-E)]$ と置く．このとき
$j^*\alpha' = [\mathcal{F}]$ である．Weil Cohomology Theories, Lemma
\ref{weil-lemma-lambda-operations} により，
$j^*\lambda^i(\alpha') = [\wedge^i(\mathcal{F})]$ と分かる．これは $K_0(X)$ において
$[\mathcal{O}_E] \cdot \alpha' = [\wedge^i\mathcal{F}]$ であることを意味する．
Derived Categories of Schemes, Lemma
\ref{perfect-lemma-projection-formula} を参照せよ．$X$ における $Z$ の既約成分の
余次元の最大値を $r$ とする．省略する計算により，$i \geq 0, 1, \ldots, r - 1$
に対して $H^{-i}(Lb^*\mathcal{O}_Z) = \wedge^i\mathcal{F}$ であり，その他の次数では
零である．したがって $K_0(X)$ において
\begin{align*}
[Lb^*\mathcal{O}_Z] & =
\sum\nolimits_{i = 0, \ldots, r - 1} (-1)^i[\wedge^i\mathcal{F}] \\
& =
\sum\nolimits_{i = 0, \ldots, r - 1} (-1)^i[\mathcal{O}_E] \lambda^i(\alpha') \\
& =
[\mathcal{O}_E] \left(\sum\nolimits_{i = 0, \ldots, r - 1}
(-1)^i \lambda^i(\alpha')\right)
\end{align*}
が成り立つ．これにより，
$\alpha'' = \sum_{i = 0, \ldots, r - 1} (-1)^i \lambda^i(\alpha')$ として補題が
証明される．
\end{proof}



\begin{lemma}
\label{obsolete-lemma-gysin-factors-principal}
$(S, \delta)$ は Chow Homology, Situation \ref{chow-situation-setup} のとおりとする．
$X$ は $S$ 上局所有限型とする．$X$ は整で $n = \dim_\delta(X)$ とする．
$a \in \Gamma(X, \mathcal{O}_X)$ を零でない関数とする．
$i : D = Z(a) \to X$ を $a$ の零点スキームの閉埋入とする．
$f \in R(X)^*$ とする．この場合，$A_{n - 2}(D)$ において
$i^*\text{div}_X(f) = 0$ である．
\end{lemma}

\begin{proof}
Chow Homology, Lemma \ref{chow-lemma-gysin-factors-general} の特別な場合である．
\end{proof}

\begin{remark}
\label{obsolete-remark-not-true-not-quasi-compact}
以前この注意では，Chow Homology, Lemma
\ref{chow-lemma-cycles-k-group} の矢印が一般に同型であるかどうかは不明だと述べていた．
しかし現在では，この事実の証明が見つかっている．
\end{remark}
% Frozen obsolete.tex lines 3364--3452.
\subsection{吹き上げに関する補題}
\label{obsolete-section-blowing-up-lemmas}

\noindent
この節では，吹き上げ上の適当な有効 Cartier 因子によって Cartier 因子を表すことに
関するいくつかの補題を証明する．これらの補題は \cite[Section 2.4]{F} にある．
有限型でない設定でも成り立つように定式化を修正した．Lemma
\ref{obsolete-lemma-blowing-up-intersections} の射 $b$ は無限個の吹き上げの合成となることが
あるが，任意に与えた準コンパクト開集合 $W \subset X$ 上では有限個の吹き上げだけで
足りる（そしてこれが同所の結果である）．

\begin{lemma}
\label{obsolete-lemma-push-pull-effective-Cartier}
$(S, \delta)$ は Chow Homology, Situation \ref{chow-situation-setup} のとおりとする．
$X$, $Y$ は $S$ 上局所有限型とする．$f : X \to Y$ を固有射とする．
$D \subset Y$ を有効 Cartier 因子とする．$X$, $Y$ は整で，
$n = \dim_\delta(X) = \dim_\delta(Y)$ とし，$f$ は支配的であると仮定する．このとき
$$
f_*[f^{-1}(D)]_{n - 1} = [R(X) : R(Y)] [D]_{n - 1}.
$$
特に，$f$ が双有理ならば
$f_*[f^{-1}(D)]_{n - 1} = [D]_{n - 1}$ である．
\end{lemma}

\begin{proof}
Chow Homology, Lemma \ref{chow-lemma-equal-c1-as-cycles}，および $D$ が
$\mathcal{O}_X(D)$ の標準切断 $1_D$ の零点スキームであることから直ちに従う．
\end{proof}

\begin{lemma}
\label{obsolete-lemma-blowing-up-denominators}
$(S, \delta)$ は Chow Homology, Situation \ref{chow-situation-setup} のとおりとする．
$X$ は $S$ 上局所有限型とする．$X$ は整で
$\dim_\delta(X) = n$ と仮定する．$\mathcal{L}$ を可逆
$\mathcal{O}_X$-加群とする．$s$ を $\mathcal{L}$ の零でない有理型切断とする．
$s$ が $U$ 上の $\mathcal{L}$ の切断に対応するような最大の開部分スキームを
$U \subset X$ とする．射影射
$$
\pi : X' \longrightarrow X
$$
で，次を満たすものが存在する：
\begin{enumerate}
\item $X'$ は整である；
\item $\pi|_{\pi^{-1}(U)} : \pi^{-1}(U) \to U$ は同型である；
\item 有効 Cartier 因子 $D, E \subset X'$ で
$$
\pi^*\mathcal{L} = \mathcal{O}_{X'}(D - E),
$$
を満たすものが存在する；
\item 上の同型を介して，有理型切断 $s$ は有理型切断
$1_D \otimes (1_E)^{-1}$ に対応する（Divisors, Definition
\ref{divisors-definition-invertible-sheaf-effective-Cartier-divisor} を参照せよ）；
\item $Z_{n - 1}(X)$ において
$$
\pi_*([D]_{n - 1} - [E]_{n - 1}) = \text{div}_\mathcal{L}(s)
$$
が成り立つ．
\end{enumerate}
\end{lemma}

\begin{proof}
$\mathcal{I} \subset \mathcal{O}_X$ を $s$ の分母の準連接イデアル層とする．
Divisors, Definition
\ref{divisors-definition-regular-meromorphic-ideal-denominators} を参照せよ．
Divisors, Lemma \ref{divisors-lemma-blowing-up-denominators} により (2), (3), (4) を得る．
Divisors, Lemma \ref{divisors-lemma-blow-up-integral-scheme} により (1) を得る．
Divisors, Lemma \ref{divisors-lemma-blowing-up-projective} により射 $\pi$ は射影的である．
残るのは (5) の証明である．Chow Homology, Lemma
\ref{chow-lemma-equal-c1-as-cycles} により
$$
\pi_*(\text{div}_{\mathcal{L}'}(s')) = \text{div}_\mathcal{L}(s).
$$
したがって
$\text{div}_{\mathcal{L}'}(s') = [D]_{n - 1} - [E]_{n - 1}$
を示せば十分である．これは等式 $s' = 1_D \otimes 1_E^{-1}$ と加法性から従う．
Divisors, Lemma \ref{divisors-lemma-c1-additive} を参照せよ．
\end{proof}
% Frozen obsolete.tex lines 3453--3574.
\begin{definition}
\label{obsolete-definition-epsilon}
$(S, \delta)$ は Chow Homology, Situation \ref{chow-situation-setup} のとおりとする．
$X$ は $S$ 上局所有限型とする．$X$ は整で
$\dim_\delta(X) = n$ と仮定する．$D_1, D_2$ を $X$ の二つの有効 Cartier 因子とする．
$Z \subset X$ を $\dim_\delta(Z) = n - 1$ を満たす整閉部分スキームとする．
この状況の {it $\epsilon$-不変量} は
$$
\epsilon_Z(D_1, D_2) = n_Z \cdot m_Z
$$
である．ここで $n_Z$（それぞれ $m_Z$）は，$(n - 1)$-サイクル
$[D_1]_{n - 1}$（それぞれ $[D_2]_{n - 1}$）における $Z$ の係数である．
\end{definition}

\begin{lemma}
\label{obsolete-lemma-two-divisors}
$(S, \delta)$ は Chow Homology, Situation \ref{chow-situation-setup} のとおりとする．
$X$ は $S$ 上局所有限型とする．$X$ は整で
$\dim_\delta(X) = n$ と仮定する．$D_1, D_2$ を $X$ の二つの有効 Cartier 因子とする．
$Z$ をスキーム $D_1 \cap D_2$ の開かつ閉な部分スキームとする．
$\dim_\delta(D_1 \cap D_2 \setminus Z) \leq n - 2$ と仮定する．
このとき射 $b : X' \to X$ と $X'$ 上の Cartier 因子 $D_1', D_2', E$ で，
次の性質をもつものが存在する：
\begin{enumerate}
\item $X'$ は整である；
\item $b$ は射影的である；
\item $b$ は閉部分スキーム $Z$ に沿った $X$ の吹き上げである；
\item $E = b^{-1}(Z)$；
\item $b^{-1}(D_1) = D'_1 + E$，かつ $b^{-1}D_2 = D_2' + E$；
\item $\dim_\delta(D'_1 \cap D'_2) \leq n - 2$，また $Z = D_1 \cap D_2$ ならば
$D'_1 \cap D'_2 = \emptyset$；
\item $\dim_\delta(W') = n - 1$ を満たす任意の整閉部分スキーム $W'$ に対して，
\begin{enumerate}
\item $\epsilon_{W'}(D'_1, E) > 0$ ならば，$W = b(W')$ と置くと
$\dim_\delta(W) = n - 1$ であり，
$$
\epsilon_{W'}(D'_1, E) < \epsilon_W(D_1, D_2),
$$
\item $\epsilon_{W'}(D'_2, E) > 0$ ならば，$W = b(W')$ と置くと
$\dim_\delta(W) = n - 1$ であり，
$$
\epsilon_{W'}(D'_2, E) < \epsilon_W(D_1, D_2),
$$
\end{enumerate}
\end{enumerate}
\end{lemma}

\begin{proof}
準連接イデアル層
$\mathcal{I} = \mathcal{I}_{D_1} + \mathcal{I}_{D_2}$ は，スキーム論的交叉
$D_1 \cap D_2 \subset X$ を定めることに注意する．$Z$ は $D_1 \cap D_2$ の連結成分の
合併なので，任意の $z \in Z$ に対して
$\mathcal{O}_{X, z} \to \mathcal{O}_{Z, z}$ の核は $\mathcal{I}_z$ に等しい．
$b : X' \to X$ を $Z$ に沿った $X$ の吹き上げとする（したがって $Z$ の周りで
Zariski 局所的には，これは $\mathcal{I}$ に沿った $X$ の吹き上げである）．
対応する有効 Cartier 因子を $E = b^{-1}(Z)$ と記す．Divisors, Lemma
\ref{divisors-lemma-blowing-up-gives-effective-Cartier-divisor} を参照せよ．
$Z \subset D_1$ なので $E \subset f^{-1}(D_1)$ であり，したがってある有効 Cartier
因子 $D'_1 \subset X'$ に対して $D_1 = D_1' + E$ である．Divisors, Lemma
\ref{divisors-lemma-difference-effective-Cartier-divisors} を参照せよ．同様に
$D_2 = D_2' + E$ である．これで主張 (1)--(5) が得られた．

\medskip\noindent
$W'$ が (7) (a) または (7) (b) のとおりならば，$W'$ の像 $W$ は
$D_1 \cap D_2$ に含まれることに注意する．$W$ が $Z$ に含まれないならば，$b$ は
$W$ の生成点で同型であり，
$\dim_\delta(W) = \dim_\delta(W') = n - 1$ と分かる．これは
$\dim_\delta(D_1 \cap D_2 \setminus Z) \leq n - 2$ という仮定に反する．したがって
$W \subset Z$ である．ゆえに (6) と (7) を証明するには，$X$ 上 $Z$ の周りで
局所的に作業してよい．

\medskip\noindent
そこで，$X = \Spec(A)$，$A$ は Noether 整域，
$D_1 = \Spec(A/a)$, $D_2 = \Spec(A/b)$, $Z = D_1 \cap D_2$ と仮定してよい．
$I = (a, b)$ と置く．$A$ は整域で $a, b \not = 0$ なので，吹き上げは二つの
チャート $U = \Spec(A[s]/(as - b))$ と
$V = \Spec(A[t]/(bt -a))$ で被覆される．これらのチャートは，$s$ を $t^{-1}$ へ
写す同型
$A[s, s^{-1}]/(as - b) \cong A[t, t^{-1}]/(bt - a)$
を用いて貼り合わせる．有効 Cartier 因子 $E$ は
$\Spec(A[s]/(as - b, a)) \subset U$ および
$\Spec(A[t]/(bt - a, b)) \subset V$ で記述される．閉部分スキーム $D'_1$ は
$\Spec(A[t]/(bt - a, t)) \subset U$ に対応する．閉部分スキーム $D'_2$ は
$\Spec(A[s]/(as -b, s)) \subset V$ に対応する．「$ts = 1$」なので
$D'_1 \cap D'_2 = \emptyset$ である．

\medskip\noindent
高さ 1 の素イデアル $\mathfrak q \subset A[s]/(as - b)$ で
$s, a \in \mathfrak q$ を満たすものがあるとする．対応する $A$ の素イデアルを
$\mathfrak p \subset A$ とする．$a, b \in \mathfrak p$ であることに注意する．
次元公式により $\dim(A_{\mathfrak p}) = 1$ でもある．最後に示すべき主張は
$$
\text{ord}_{A_{\mathfrak p}}(a)
\text{ord}_{A_{\mathfrak p}}(b)
>
\text{ord}_{B_{\mathfrak q}}(a)
\text{ord}_{B_{\mathfrak q}}(s)
$$
である．ここで $B = A[s]/(as - b)$ である．Algebra, Lemma
\ref{algebra-lemma-quasi-finite-extension-dim-1} により，$x = a, b$ に対して
$\text{ord}_{A_{\mathfrak p}}(x) \geq \text{ord}_{B_{\mathfrak q}}(x)$ である．
$\text{ord}_{B_{\mathfrak q}}(s) > 0$ なので，$\text{ord}$ 函数の加法性と
$as = b$ という事実から結論を得る．
\end{proof}
% Frozen obsolete.tex lines 3575--3617.
\begin{definition}
\label{obsolete-definition-locally-finite-sum-effective-Cartier-divisors}
$X$ をスキームとする．$\{D_i\}_{i \in I}$ を $X$ 上の有効 Cartier 因子の
局所有限な族とする．函数 $I \to \mathbf{Z}_{\geq 0}$, $i \mapsto n_i$ が
与えられているとする．{it 有効 Cartier 因子の和}
$D = \sum n_i D_i$ とは，任意の準コンパクト開集合 $U \subset X$ 上で
$D|_U = \sum_{D_i \cap U \not = \emptyset} n_iD_i|_U$ が Divisors, Definition
\ref{divisors-definition-sum-effective-Cartier-divisors} の意味での和となるような，一意な
有効 Cartier 因子 $D \subset X$ である．
\end{definition}

\begin{lemma}
\label{obsolete-lemma-sum-divisors-associated-Weil}
$(S, \delta)$ は Chow Homology, Situation \ref{chow-situation-setup} のとおりとする．
$X$ は $S$ 上局所有限型とする．$X$ は整で
$\dim_\delta(X) = n$ と仮定する．$\{D_i\}_{i \in I}$ を $X$ 上の有効 Cartier
因子の局所有限な族とする．$i \in I$ に対して $n_i \geq 0$ が与えられているとする．
このとき $Z_{n - 1}(X)$ において
$$
[D]_{n - 1} = \sum\nolimits_i n_i[D_i]_{n - 1}
$$
が成り立つ．
\end{lemma}

\begin{proof}
サイクルの等式を証明するので，$X$ 上で局所的に作業してよい．したがって有限和に，
さらに帰納法によって二つの有効 Cartier 因子の和 $D = D_1 + D_2$ に帰着する．
Chow Homology, Lemma \ref{chow-lemma-compute-c1} により，
$D_1 = \text{div}_{\mathcal{O}_X(D_1)}(1_{D_1})$ である．ここで $1_{D_1}$ は
$\mathcal{O}_X(D_1)$ の標準切断を表す．もちろん $D_2$ と $D$ に対しても同じ主張が
成り立つ．同定
$\mathcal{O}_X(D) = \mathcal{O}_X(D_1) \otimes \mathcal{O}_X(D_2)$
を介して $1_D = 1_{D_1} \otimes 1_{D_2}$ なので，Divisors, Lemma
\ref{divisors-lemma-c1-additive} から結論を得る．
\end{proof}
% Frozen obsolete.tex lines 3618--3734.
\begin{lemma}
\label{obsolete-lemma-blowing-up-intersections}
$(S, \delta)$ は Chow Homology, Situation \ref{chow-situation-setup} のとおりとする．
$X$ は $S$ 上局所有限型とする．$X$ は整で
$\dim_\delta(X) = d$ と仮定する．$\{D_i\}_{i \in I}$ を $X$ 上の有効 Cartier
因子の局所有限な族とする．$\{i, j, k\} \subset I$ かつ
$\#\{i, j, k\} = 3$ であるすべての組に対して
$D_i \cap D_j \cap D_k = \emptyset$ と仮定する．このとき次が存在する：
\begin{enumerate}
\item $\dim_\delta(X \setminus U) \leq d - 3$ を満たす開部分スキーム
$U \subset X$；
\item 射 $b : U' \to U$；
\item $U'$ 上の有効 Cartier 因子 $\{D'_j\}_{j \in J}$．
\end{enumerate}
これらは次の性質をもつ：
\begin{enumerate}
\item $b$ は固有射 $b : U' \to U$ である；
\item $U'$ は整である；
\item $b$ は，$D_i|_U$ たちの二つずつの交叉の合併の補集合上で同型である；
\item $\{D'_j\}_{j \in J}$ は $U'$ 上の有効 Cartier 因子の局所有限な族である；
\item $j \not = j'$ ならば
$\dim_\delta(D'_j \cap D'_{j'}) \leq d - 2$；
\item ある $n_{ij} \geq 0$ に対して
$b^{-1}(D_i|_U) = \sum n_{ij} D'_j$．
\end{enumerate}
さらに $X$ が準コンパクトならば，上で $U = X$ と仮定できる．
\end{lemma}

\begin{proof}
まず準コンパクトな場合を証明する．これはおそらく最も興味深い場合でもある．
この場合，吹き上げの列
$$
X = X_0 \xleftarrow{b_0} X_1 \xleftarrow{b_1} X_2 \leftarrow \ldots
$$
と有効 Cartier 因子の有限集合 $\{D_{n, i}\}_{i \in I_n}$ を帰納的に構成する．
各段階で，任意の三重交叉
$D_{n, i} \cap D_{n, j} \cap D_{n, k}$ は空である．さらに各 $n \geq 0$ に対し，
$I_{n + 1} = I_n \amalg P(I_n)$ とする．ここで $P(I_n)$ は $I_n$ の元の対の集合を
表す．最後に
$$
b_n^{-1}(D_{n, i}) = D_{n + 1, i} +
\sum\nolimits_{i' \in I_n, i' \not = i} D_{n + 1, \{i, i'\}}
$$
とする．その結果，各 $n \geq 0$ に対して
$(b_0 \circ \ldots \circ b_n)^{-1}(D_i)$ は因子 $D_{n + 1, j}$,
$j \in I_{n + 1}$ の非負整数係数の線形結合となる．

\medskip\noindent
帰納法の初期値として $X_0 = X$, $I_0 = I$, $D_{0, i} = D_i$ と置く．

\medskip\noindent
$(X_n, \{D_{n, i}\}_{i \in I_n})$ が与えられたとき，$X_{n + 1}$ を閉部分スキーム
$Z_n = \bigcup_{\{i, i'\} \in P(I_n)} D_{n, i} \cap D_{n, i'}$
に沿った $X_n$ の吹き上げとする．三重交叉に関する仮定により，閉部分スキーム
$D_{n, i} \cap D_{n, i'}$ たちは二つずつ互いに素であることに注意する．言い換えると
$Z_n = \coprod_{\{i, i'\} \in P(I_n)} D_{n, i} \cap D_{n, i'}$
と書ける．さらに $D_{n, i} \cap D_{n, i'}$ の Zariski 近傍では，射 $b_n$ は閉部分
スキーム $D_{n, i} \cap D_{n, i'}$ に沿ったスキーム $X_n$ の吹き上げに等しく，
Lemma \ref{obsolete-lemma-two-divisors} の結果を適用できる．したがって
$D_{n + 1, \{i, i'\}} = b_n^{-1}(D_i \cap D_{i'})$ と置くと，有効 Cartier 因子を
得る．Cartier 因子 $D_{n + 1, \{i, i'\}}$ たちは二つずつ互いに素である．
任意の $i' \in I_n$, $i' \not = i$ に対して，明らかに
$b_n^{-1}(D_{n, i}) \supset D_{n + 1, \{i, i'\}}$ である．したがって Divisors,
Lemma \ref{divisors-lemma-difference-effective-Cartier-divisors} を適用すると，実際に
ある $X_{n + 1}$ 上の有効 Cartier 因子 $D_{n + 1, i}$ に対して
$b^{-1}(D_{n, i}) = D_{n + 1, i} +
\sum\nolimits_{i' \in I_n, i' \not = i} D_{n + 1, \{i, i'\}}$
となる．$D_{n + 1, \{i, i'\}}$ の近傍では，これらの因子 $D_{n + 1, i}$ は
Lemma \ref{obsolete-lemma-two-divisors} のプライム付き因子の役割を果たす．特に
Lemma \ref{obsolete-lemma-two-divisors} の (6) により，$i \not = i'$, $i, i' \in I_n$ ならば
$D_{n + 1, i} \cap D_{n + 1, i'} = \emptyset$ と結論できる．これはすでに因子
$D_{n + 1, i}$ たちの三重交叉が零であることを含意する．

\medskip\noindent
ここで $X$ の準コンパクト性を用いると，不変量
\begin{equation}
\label{obsolete-equation-invariant}
\epsilon(X, \{D_i\}_{i \in I}) =
\max\{\epsilon_Z(D_i, D_{i'}) \mid
Z \subset X,
\dim_\delta(Z) = d - 1,
\{i, i'\} \in P(I)\}
\end{equation}
は有限であると分かる．実際，各 $D_i$ は有限個の既約成分しかもたない．ある $n$ に
対して不変量 $\epsilon(X_n, \{D_{n, i}\}_{i \in I_n})$ は零であると主張する．
そうでなければ Lemma \ref{obsolete-lemma-two-divisors} により，正の整数の真に減少する列
$$
\epsilon(X, \{D_i\}_{i \in I})
=
\epsilon(X_0, \{D_{0, i}\}_{i \in I_0})
>
\epsilon(X_1, \{D_{1, i}\}_{i \in I_1})
>
\ldots
$$
が得られ，矛盾する．不変量 $\epsilon(X_n, \{D_{n, i}\}_{i \in I_n})$ が零となる
$n$ を取る．これは，整閉部分スキーム $Z \subset X_n$ と添字の対
$i, i' \in I_n$ で $\epsilon_Z(D_{n, i}, D_{n, i'}) > 0$ を満たすものがないことを
意味する．言い換えると，すべての対 $\{i, i'\} \in P(I_n)$ に対して，望みどおり
$\dim_\delta(D_{n, i}, D_{n, i'}) \leq d - 2$ である．
% Frozen obsolete.tex lines 3735--3900.
\medskip\noindent
次に，スキーム $X$ が準コンパクトであるとは仮定しない一般の場合を考える．証明の
最初の部分の考え方には，$X$ のある固定した点を支配する中心をもつ吹き上げが無限列を
なすおそれがあるという問題がある．これを避けるため，各段階で余次元
$\geq 3$ の適当な閉部分集合を除く．すなわち，次の形の射の列を帰納的に構成する：
$$
\xymatrix{
X = X_0 \\
U_0 \ar[u]^{j_0} & X_1 \ar[l]_{b_0} \\
 & U_1 \ar[u]^{j_1} & X_2 \ar[l]_{b_1} \\
 & & U_2 \ar[u]^{j_2} & X_3 \ar[l]_{b_2}
}
$$
各射 $j_n : U_n \to X_n$ は開埋入である．各射
$b_n : X_{n + 1} \to U_n$ は整スキームの固有双有理射である．準コンパクトな場合と
同様，$X_n$ 上に有効 Cartier 因子 $\{D_{n, i}\}_{i \in I_n}$ がある．各段階で，
任意の三重交叉 $D_{n, i} \cap D_{n, j} \cap D_{n, k}$ は空である．さらに各
$n \geq 0$ に対して $I_{n + 1} = I_n \amalg P(I_n)$ とする．ここで $P(I_n)$ は
$I_n$ の元の対の集合を表す．最後に
$$
b_n^{-1}(D_{n, i}|_{U_n}) = D_{n + 1, i} +
\sum\nolimits_{i' \in I_n, i' \not = i} D_{n + 1, \{i, i'\}}
$$
となるようにする．

\medskip\noindent
帰納法は $X_0 = X$, $I_0 = I$, $D_{0, i} = D_i$ と置いて始める．

\medskip\noindent
$(X_n, \{D_{n, i}\})$ が与えられたとき，開部分スキーム $U_n$ を次のように
構成する．各対 $\{i, i'\} \in P(I_n)$ に対して閉部分スキーム
$D_{n, i} \cap D_{n, i'}$ を考える．これには，$\delta$-次元 $d - 2$ の「良い」
既約成分と，$\delta$-次元 $d - 1$ の「悪い」既約成分がある．
$$
\text{Bad}(i, i')
=
\bigcup\nolimits_{W \subset D_{n, i} \cap D_{n, i'}
\text{ irred.\ comp. with }\dim_\delta(W) = d - 1} W
$$
と置き，同様に
$$
\text{Good}(i, i')
=
\bigcup\nolimits_{W \subset D_{n, i} \cap D_{n, i'}
\text{ irred.\ comp. with }\dim_\delta(W) = d - 2} W.
$$
と置く．すると
$D_{n, i} \cap D_{n, i'} = \text{Bad}(i, i') \cup \text{Good}(i, i')$ であり，さらに
$\dim_\delta(\text{Bad}(i, i') \cap \text{Good}(i, i')) \leq d - 3$ である．
$U_n$ を次のように選ぶ：
$$
U_n
=
X_n
\setminus
\bigcup\nolimits_{\{i, i'\} \in P(I_n)}
\text{Bad}(i, i') \cap \text{Good}(i, i').
$$
因子 $D_{n, i}$ の三重交叉に関する条件から，この合併は実際には直和であると分かる．
さらにスキームとして
$$
D_{n, i}|_{U_n} \cap D_{n, i'}|_{U_n}
=
Z_{n, i, i'} \amalg G_{n, i, i'}
$$
である．ここで $Z_{n, i, i'}$ は次元 $d - 1$ の $\delta$-純次元であり，
$G_{n, i, i'}$ は次元 $d - 2$ の $\delta$-純次元である（したがって位相的には，
$Z_{n, i, i'}$ は悪い成分の合併から良い成分との交叉を除いたものである）．最後に
$$
Z_n =
\bigcup\nolimits_{\{i, i'\} \in P(I_n)} Z_{n, i, i'} =
\coprod\nolimits_{\{i, i'\} \in P(I_n)} Z_{n, i, i'},
$$
と置き，$b_n : X_{n + 1} \to X_n$ を $Z_n$ に沿った吹き上げとする．Lemma
\ref{obsolete-lemma-two-divisors} は，各軌跡
$D_{n, i}|_{U_n} \cap D_{n, i'}|_{U_n}$ の周りで局所的に射
$b_n : X_{n + 1} \to X_n$ に適用できることに注意する．したがって証明の最初の部分と
まったく同様に，$\{i, i'\} \in P(I_n)$ に対する有効 Cartier 因子
$D_{n + 1, \{i, i'\}}$ と，$i \in I_n$ に対する有効 Cartier 因子
$D_{n + 1, i}$ で
$b_n^{-1}(D_{n, i}|_{U_n}) = D_{n + 1, i} +
\sum\nolimits_{i' \in I_n, i' \not = i} D_{n + 1, \{i, i'\}}$
を満たすものを得る．各 $n$ に対し，合成
$j_0 \circ \ldots \circ j_{n - 1} \circ b_{n - 1}$ として得られる射
$\pi_n : X_n \to X$ を記す．

\medskip\noindent
{\bf 主張：} 任意の準コンパクト開集合 $V \subset X$ に対し，十分大きいすべての
$n$ について，写像
$$
\pi_n^{-1}(V) \leftarrow \pi_{n + 1}^{-1}(V) \leftarrow \ldots
$$
はすべて同型である．実際，写像
$\pi_n^{-1}(V) \leftarrow \pi_{n + 1}^{-1}(V)$ が同型でなければ，ある
$\{i, i'\} \in P(I_n)$ に対して
$Z_{n, i, i'} \cap \pi_n^{-1}(V) \not = \emptyset$ である．したがって，
$\dim_\delta(W) = d - 1$ を満たす既約成分
$W \subset D_{n, i} \cap D_{n, i'}$ が存在する．特に
$\epsilon_W(D_{n, i}, D_{n, i'}) > 0$ である．Lemma
\ref{obsolete-lemma-two-divisors} を繰り返し適用すると
$$
\epsilon_W(D_{n, i}, D_{n, i'})
<
\epsilon(V, \{D_i|_V\}) - n
$$
を得る．ここで $\epsilon(V, \{D_i|_V\})$ は
(\ref{obsolete-equation-invariant}) のとおりである．$V$ は準コンパクトなので
$\epsilon(V, \{D_i|_V\}) < \infty$ であり，
$n > \epsilon(V, \{D_i|_V\})$ と取れば結果を得る．

\medskip\noindent
構成により $X_n \setminus U_n$ は
$\dim_\delta(X_n \setminus U_n) \leq d - 3$ を満たすことに注意する．その $X$ に
おける像を $T_n = \pi_n(X_n \setminus U_n)$ とする．上の射の図式をたどり，各
$b_n$ が閉写像であることを用いると，各 $n$ に対して
$T_0 \cup \ldots \cup T_n$ は $X$ の閉部分集合であると分かる．構成により，任意の
$t \in T_n$ は $\delta(t) \leq d - 3$ を満たす．したがって
$\overline{T_n} \subset X$ は $\dim_\delta(T_n) \leq d - 3$ を満たす閉部分集合である．
上の主張により，任意の準コンパクト開集合 $V \subset X$ に対して
$T_n \cap V \not = \emptyset$ となる $n$ は有限個しかない．したがって
$\{\overline{T_n}\}_{n \geq 0}$ は閉部分集合の局所有限な族であり，
$U = X \setminus \bigcup \overline{T_n}$ と置ける．これが補題の $U$ となる．

\medskip\noindent
$U$ の構成により $U_n \cap \pi_n^{-1}(U) = \pi_n^{-1}(U)$ であることに注意する．
したがって射
$$
b_n : \pi_{n + 1}^{-1}(U) \longrightarrow \pi_n^{-1}(U)
$$
はすべて固有である．さらに上の主張により，$X$ の各準コンパクト開集合上で最終的に
同型となる．ゆえに
$$
U' = \lim_n \pi_n^{-1}(U).
$$
と定義できる．誘導される射 $b : U' \to U$ は固有である．これは $U$ 上局所的な
性質であり，各コンパクト開集合上で極限が安定するからである．同様に，構成から得られる
包含 $I_n \to I_{n + 1}$ を用いて $J = \bigcup_{n \geq 0} I_n$ と置く．$j \in J$ に
対し，$j$ が $i \in I_{n_0}$ に対応するような $n_0$ を選び，
$D'_j = \lim_{n \geq n_0} D_{n, i}$ と定義する．これも $X$ 上局所的には射が安定する
ので意味をもつ．補題の他の主張は，$X$ が準コンパクトな場合と同様に確認される．
\end{proof}
% Frozen obsolete.tex lines 3901--4024.
\section{因子の交叉の可換性}
\label{obsolete-section-commutativity}

\noindent
この節の結果は元来，Chow Homology, Section
\ref{chow-section-commutativity} の補題に別証明を与え，Chow Homology, Lemma
\ref{chow-lemma-gysin-commutes-gysin} の弱い版を与えるために用いられていた．

\begin{lemma}
\label{obsolete-lemma-improved-additivity}
$(S, \delta)$ は Chow Homology, Situation \ref{chow-situation-setup} のとおりとする．
$X$ は $S$ 上局所有限型とする．$\{i_j : D_j \to X \}_{j \in J}$ を $X$ 上の
有効 Cartier 因子の局所有限な族とする．$n_j > 0$, $j\in J$ とする．
$D = \sum_{j \in J} n_j D_j$ と置き，包含射を $i : D \to X$ と記す．
$\alpha \in Z_{k + 1}(X)$ とする．このとき
$$
p : \coprod\nolimits_{j \in J} D_j \longrightarrow D
$$
は固有であり，$\CH_k(D)$ において
$$
i^*\alpha = p_*\left(\sum n_j i_j^*\alpha\right)
$$
が成り立つ．
\end{lemma}

\begin{proof}
有理同値の無限和に関する微妙な点のため，この補題の証明は予想より少し長くなる．
準コンパクトな場合，族 $D_j$ は有限で，結果はまったく容易であり，Chow Homology,
Lemma \ref{chow-lemma-compute-c1} と Divisors, Lemma
\ref{divisors-lemma-c1-additive} および定義から直接従う．

\medskip\noindent
族 $\{D_j\}_{j \in J}$ は局所有限なので，射 $p$ は固有である．
$W_a \subset X$ は $\delta$-次元 $k + 1$ の整閉部分スキームとして，
$\alpha = \sum_{a \in A} m_a [W_a]$ と書く．閉埋入を
$i_a : W_a \to X$ と記す．$\{W_a\}_{a \in A}$ が $X$ 上局所有限となるような
すべての $a \in A$ に対して $m_a \not = 0$ と仮定する．

\medskip\noindent
Chow Homology, Definition \ref{chow-definition-gysin-homomorphism} により，類
$i^*\alpha$ は，ある $\beta_a \in Z_k(W_a \cap D)$ に対するサイクル
$\sum m_a\beta_a$ の類であることに注意する．すなわち，
$W_a \not \subset D$ ならば $\beta_a = [D \cap W_a]_k$ であり，
$W_a \subset D$ ならば $\beta_a$ は
$c_1(\mathcal{O}_X(D)) \cap [W_a]$ を表すサイクルである．

\medskip\noindent
各 $a \in A$ に対して $J = J_{a, 1} \amalg J_{a, 2} \amalg J_{a, 3}$ と書く．ここで
\begin{enumerate}
\item $j \in J_{a, 1}$ であることと $W_a \cap D_j = \emptyset$ は同値である；
\item $j \in J_{a, 2}$ であることと
$W_a \not = W_a \cap D_1 \not = \emptyset$ は同値である；
\item $j \in J_{a, 3}$ であることと $W_a \subset D_j$ は同値である．
\end{enumerate}
族 $\{D_j\}$ は局所有限なので，$J_{a, 3}$ は有限集合である．各 $a \in A$ と
$j \in J$ に対して，サイクル $\beta_{a, j} \in Z_k(W_a \cap D_j)$ を次のように選ぶ：
\begin{enumerate}
\item $j \in J_{a, 1}$ ならば $\beta_{a, j} = 0$ と置く；
\item $j \in J_{a, 2}$ ならば $\beta_{a, j} = [D_j \cap W_a]_k$ と置く；
\item $j \in J_{a, 3}$ ならば，
$c_1(i_a^*\mathcal{O}_X(D_j)) \cap [W_j]$ を表す
$\beta_{a, j} \in Z_k(W_a)$ を選ぶ．
\end{enumerate}
$\CH_k(W_a \cap D)$ において
$$
\beta_a \sim_{rat}
\sum\nolimits_{j \in J} n_j \beta_{a, j}
$$
であると主張する．

\medskip\noindent
場合 I：$W_a \not \subset D$．この場合 $J_{a, 3} = \emptyset$ である．したがって
サイクルとして
$[D \cap W_a]_k = \sum n_j [D_j \cap W_a]_k$ を示せば十分である．これは Lemma
\ref{obsolete-lemma-sum-divisors-associated-Weil} である．

\medskip\noindent
場合 II：$W_a \subset D$．この場合，$\beta_a$ は
$c_1(i_a^*\mathcal{O}_X(D)) \cap [W_a]$ を表すサイクルである．
$D_{a, s} = \sum_{j \in J_{a, s}} n_jD_j$ として
$D = D_{a, 1} + D_{a, 2} + D_{a, 3}$ と書く．Divisors, Lemma
\ref{divisors-lemma-c1-additive} により
\begin{eqnarray*}
c_1(i_a^*\mathcal{O}_X(D)) \cap [W_a] & = &
c_1(i_a^*\mathcal{O}_X(D_{a, 1})) \cap [W_a] +
c_1(i_a^*\mathcal{O}_X(D_{a, 2})) \cap [W_a] \\
& &
 + c_1(i_a^*\mathcal{O}_X(D_{a, 3})) \cap [W_a].
\end{eqnarray*}
和の第1項は明らかに零である．$J_{a, 3}$ は有限なので，最後の項は
$\sum\nolimits_{j \in J_{a, 3}} n_jc_1(i_a^*\mathcal{L}_j) \cap [W_a]$
と一致する．Divisors, Lemma \ref{divisors-lemma-c1-additive} を参照せよ．これは
$\sum_{j \in J_{a, 3}} n_j \beta_{a, j}$ によって表される．最後に，場合 I により
中間の項はサイクル
$\sum\nolimits_{j \in J_{a, 2}} n_j[D_j \cap W_a]_k =
\sum_{j \in J_{a, 2}} n_j\beta_{a, j}$
によって表される．これでこの場合の主張を得る．

\medskip\noindent
これで補題の証明を終えられる．すなわち，$\beta_a$ の選択により
$i^*D \sim_{rat} \sum m_a\beta_a$ である．各 $a$ に対して
$\beta_a \sim_{rat} \sum_j \beta_{a, j}$ であり，この有理同値は
$D \cap W_a$ 上で成り立つ．閉部分スキーム $D \cap W_a$ の族は $D$ 上局所有限なので，
$D$ 上でも
$\sum m_a \beta_a \sim_{rat} \sum_{a, j} m_a\beta_{a, j}$
である（Chow Homology, Remark
\ref{chow-remark-infinite-sums-rational-equivalences} を参照せよ）！ さて，
$\sum_a m_a\beta_{a, j}$ は（$D_j$ 上のサイクルとみなすと）$i_j^*\alpha$ を表し，
したがって $\sum_{a, j} m_a\beta_{a, j}$ は
$p_* \sum_j i_j^*\alpha$ を表すことが明らかである．これで結論を得る．
\end{proof}
% Frozen obsolete.tex lines 4025--4104.
\begin{lemma}
\label{obsolete-lemma-commutativity-effective-Cartier-proper-intersection}
$(S, \delta)$ は Chow Homology, Situation \ref{chow-situation-setup} のとおりとする．
$X$ は $S$ 上局所有限型とする．$X$ は整で
$\dim_\delta(X) = n$ と仮定する．$D$, $D'$ を $X$ 上の有効 Cartier 因子とする．
$\dim_\delta(D \cap D') = n - 2$ と仮定する．$i : D \to X$，それぞれ
$i' : D' \to X$ を対応する閉埋入とする．このとき
\begin{enumerate}
\item サイクル $\alpha \in Z_{n - 2}(D \cap D')$ で，$D$ への押し出しが
$i^*[D']_{n - 1} \in \CH_{n - 2}(D)$ を表し，$D'$ への押し出しが
$(i')^*[D]_{n - 1} \in \CH_{n - 2}(D')$ を表すものが存在する；
\item $\CH_{n - 2}(X)$ において
$$
D \cdot [D']_{n - 1}
=
D' \cdot [D]_{n - 1}
$$
が成り立つ．
\end{enumerate}
\end{lemma}

\begin{proof}
(2) は (1) の自明な帰結である．$Z_a$ は $D$ の既約成分，$[Z_b]$ は $D'$ の
既約成分として，$[D]_{n - 1} = \sum n_a[Z_a]$ および
$[D']_{n - 1} = \sum m_b[Z_b]$ と書く．Chow Homology, Definition
\ref{chow-definition-gysin-homomorphism} によれば，
$i^*D' = \sum m_b i^*[Z_b]$ および
$(i')^*D = \sum n_a(i')^*[Z_a]$ である．仮定により，どの既約成分 $Z_b$ も
$D$ に含まれず，したがって定義から
$i^*[Z_b] = [Z_b\cap D]_{n - 2}$ である．同様に
$(i')^*[Z_a] = [Z_a \cap D']_{n - 2}$ である．したがって，実際に
$D \cap D'$ 上に台をもつサイクルの等式
$$
\sum n_a[Z_a \cap D']_{n - 2} = \sum m_b[Z_b \cap D]_{n - 2}
$$
を証明しようとしている．$W \subset X$ を
$\dim_\delta(W) = n - 2$ を満たす整閉部分スキームとする．その生成点を
$\xi \in W$ とする．$R = \mathcal{O}_{X, \xi}$ と置く．これは Noether 局所整域である．
$\dim(R) = 2$ であることに注意する．$D$，それぞれ $D'$ のイデアルを定める元を
$f \in R$，それぞれ $f' \in R$ とする．仮定により
$\dim(R/(f, f')) = 0$ である．$(f')$ 上の極小素イデアルを
$\mathfrak q'_1, \ldots, \mathfrak q'_t \subset R$ とし，$(f)$ 上の極小素イデアルを
$\mathfrak q_1, \ldots, \mathfrak q_s \subset R$ とする．上の等式は
$$
\sum_{i = 1, \ldots, s}
\text{length}_{R_{\mathfrak q_i}}(R_{\mathfrak q_i}/(f))
\text{ord}_{R/\mathfrak q_i}(f')
=
\sum_{j = 1, \ldots, t}
\text{length}_{R_{\mathfrak q'_j}}(R_{\mathfrak q'_j}/(f'))
\text{ord}_{R/\mathfrak q'_j}(f).
$$
に帰着する．Chow Homology, Lemma
\ref{chow-lemma-length-multiplication} を $M = R/(f)$ として適用すると，この等式の
左辺は
$$
\text{length}_R(R/(f, f'))
-
\text{length}_R(\Ker(f' : R/(f) \to R/(f)))
$$
に等しい．ここで，$x \bmod (f) \mapsto f'x \bmod (ff')$ により
$\Ker(f' : R/(f) \to R/(f))$ は $((f) \cap (f'))/(ff')$ と標準的に同型であることに
注意する．したがって左辺は
$$
\text{length}_R(R/(f, f'))
-
\text{length}_R((f) \cap (f')/(ff'))
$$
である．これは $f$ と $f'$ に関して対称なので，結論を得る．
\end{proof}
% Frozen obsolete.tex lines 4105--4193.
\begin{lemma}
\label{obsolete-lemma-commutativity-effective-Cartier-proper-intersection-infinite}
$(S, \delta)$ は Chow Homology, Situation \ref{chow-situation-setup} のとおりとする．
$X$ は $S$ 上局所有限型とする．$X$ は整で
$\dim_\delta(X) = n$ と仮定する．$\{D_j\}_{j \in J}$ を $X$ 上の有効 Cartier 因子の
局所有限な族とする．$n_j, m_j \geq 0$ を非負整数の族とする．
$D = \sum n_j D_j$ および $D' = \sum m_j D_j$ と置く．任意の
$j \not = j'$ に対して $\dim_\delta(D_j \cap D_{j'}) = n - 2$ と仮定する．このとき
$\CH_{n - 2}(X)$ において
$D \cdot [D']_{n - 1} = D' \cdot [D]_{n - 1}$ である．
\end{lemma}

\begin{proof}
和が有限の場合，たとえば $X$ が準コンパクトな場合には，この補題は Lemmas
\ref{obsolete-lemma-sum-divisors-associated-Weil} および
\ref{obsolete-lemma-commutativity-effective-Cartier-proper-intersection} の自明な帰結である．
したがって，読者には証明を飛ばすことを勧める．

\medskip\noindent
一般の場合の証明は次のとおりである．閉埋入を $i_j : D_j \to X$ とする．
$p : \coprod D_j \to X$ を射 $i_j$ たちの余積とする．$\{Z_a\}_{a \in A}$ を
$\bigcup D_j$ の既約成分の族とする．各 $j$ に対して
$$
[D_j]_{n - 1} = \sum d_{j, a}[Z_a].
$$
と書く．Lemma \ref{obsolete-lemma-sum-divisors-associated-Weil} により
$$
[D]_{n - 1} = \sum n_j d_{j, a} [Z_a],
\quad
[D']_{n - 1} = \sum m_j d_{j, a} [Z_a].
$$
である．Lemma \ref{obsolete-lemma-improved-additivity} により
$$
D \cdot [D']_{n - 1} = p_*\left(\sum n_j i_j^*[D']_{n - 1} \right),
\quad
D' \cdot [D]_{n - 1} = p_*\left(\sum m_{j'} i_{j'}^*[D]_{n - 1} \right).
$$
である．Gysin 準同型の定義と同様に（Chow Homology, Definition
\ref{chow-definition-gysin-homomorphism} を参照せよ），$i_j^*[Z_a]$ を表す
$D_j \cap Z_a$ 上のサイクル $\beta_{a, j}$ を選ぶ（実際，$Z_a$ が $D_j$ に
含まれないならば $\beta_{a, j} = [D_j \cap Z_a]_{n - 2}$ であり，この場合は
選択の余地がないことに注意する）．$p$ は各 $D_j$ に制限すると閉埋入なので，
$\beta_{a, j}$ を $X$ 上のサイクルとみなせるし，以下そうする．上で得た
$[D]_{n - 1}$ と $[D']_{n - 1}$ の公式を代入すると
$$
D \cdot [D']_{n - 1} =
\sum\nolimits_{j, j', a} n_j m_{j'} d_{j', a} \beta_{a, j},
\quad
D' \cdot [D]_{n - 1} =
\sum\nolimits_{j, j', a} m_{j'} n_j d_{j, a} \beta_{a, j'}.
$$
となる．さらに，同じ規約のもとで
$$
D_j \cdot [D_{j'}]_{n - 1} = \sum d_{j', a} \beta_{a, j}.
$$
でもある．この言葉で，Lemma
\ref{obsolete-lemma-commutativity-effective-Cartier-proper-intersection}（その証明も参照せよ）は，
$j \not = j'$ に対してサイクル
$\sum d_{j', a} \beta_{a, j}$ と $\sum d_{j, a} \beta_{a, j'}$ がサイクルとして
等しいと述べている！ したがって
\begin{eqnarray*}
D \cdot [D']_{n - 1}
& = &
\sum\nolimits_{j, j', a} n_j m_{j'} d_{j', a} \beta_{a, j} \\
& = &
\sum\nolimits_{j \not = j'} n_j m_{j'}
\left(\sum\nolimits_a d_{j', a} \beta_{a, j}\right) +
\sum\nolimits_{j, a} n_j m_j d_{j, a} \beta_{a, j} \\
& = &
\sum\nolimits_{j \not = j'} n_j m_{j'}
\left(\sum\nolimits_a d_{j, a} \beta_{a, j'}\right) +
\sum\nolimits_{j, a} n_j m_j d_{j, a} \beta_{a, j} \\
& = &
\sum\nolimits_{j, j', a} m_{j'} n_j d_{j, a} \beta_{a, j'} \\
& = &
D' \cdot [D]_{n - 1}
\end{eqnarray*}
となり，結論を得る．
\end{proof}
% Frozen obsolete.tex lines 4194--4311.
\begin{lemma}
\label{obsolete-lemma-commutativity-effective-Cartier}
$(S, \delta)$ は Chow Homology, Situation \ref{chow-situation-setup} のとおりとする．
$X$ は $S$ 上局所有限型とする．$X$ は整で
$\dim_\delta(X) = n$ と仮定する．$D$, $D'$ を $X$ 上の有効 Cartier 因子とする．
このとき $\CH_{n - 2}(X)$ において
$$
D \cdot [D']_{n - 1} = D' \cdot [D]_{n - 1}
$$
が成り立つ．
\end{lemma}

\begin{proof}[Lemma \ref{obsolete-lemma-commutativity-effective-Cartier} の第1証明]
まず $X$ が準コンパクトな場合に証明する．この場合，$X$ と有効 Cartier 因子の
二元集合 $\{D, D'\}$ に Lemma \ref{obsolete-lemma-blowing-up-intersections} を適用する．
すると固有射 $b : X' \to X$ と，二つずつ余次元 $\geq 2$ で交わる有効 Cartier
因子の有限族 $D'_j \subset X'$ が得られ，
$b^{-1}(D) = \sum n_j D'_j$ および $b^{-1}(D') = \sum m_j D'_j$ となる．
$Z_{n - 1}(X)$ において
$b_*[b^{-1}(D)]_{n - 1} = [D]_{n - 1}$ であり，$D'$ に対しても同様であることに
注意する．Lemma \ref{obsolete-lemma-push-pull-effective-Cartier} を参照せよ．したがって
Chow Homology, Lemma \ref{chow-lemma-pushforward-cap-c1} により，
$\CH_{n - 2}(X)$ において
$$
D \cdot [D']_{n - 1} = b_*\left(b^{-1}(D) \cdot [b^{-1}(D')]_{n - 1}\right)
$$
であり，他方の項についても同様である．したがって，補題は Lemma
\ref{obsolete-lemma-commutativity-effective-Cartier-proper-intersection-infinite} の
$\CH_{n - 2}(X')$ における等式
$b^{-1}(D) \cdot [b^{-1}(D')]_{n - 1} = b^{-1}(D') \cdot [b^{-1}(D)]_{n - 1}$
から従う．

\medskip\noindent
上の証明で参照した各補題は，一般の場合（$X$ を準コンパクトと仮定しない場合）にも
成り立つことに注意する．一般の場合の唯一の小さな変更は，Lemma
\ref{obsolete-lemma-blowing-up-intersections} を適用して得られる射 $b : U' \to U$ の終域が，
補集合の余次元が $\geq 3$ である開集合 $U \subset X$ となることである．したがって
Chow Homology, Lemma \ref{chow-lemma-restrict-to-open} により
$\CH_{n - 2}(U) = \CH_{n - 2}(X)$ であり，$X$ を $U$ で置き換えた後は，証明の残りは
変更なく進む．
\end{proof}

\begin{proof}[Lemma \ref{obsolete-lemma-commutativity-effective-Cartier} の第2証明]
$\mathcal{I} = \mathcal{O}_X(-D)$ および
$\mathcal{I}' = \mathcal{O}_X(-D')$ を，$D$ および $D'$ の可逆イデアル層とする．
$\mathcal{I}_{D'} = \mathcal{I} \otimes_{\mathcal{O}_X} \mathcal{O}_{D'}$
および
$\mathcal{I}'_D = \mathcal{I}' \otimes_{\mathcal{O}_X} \mathcal{O}_D$
と記す．包含写像 $\mathcal{I} \to \mathcal{O}_X$ を $D'$ に制限して写像
$$
\varphi : \mathcal{I}_{D'} \longrightarrow \mathcal{O}_{D'}
$$
を得られ，同様に
$$
\psi : \mathcal{I}'_D \longrightarrow \mathcal{O}_D
$$
を得る．明らかに
$$
\Coker(\varphi)
\cong
\mathcal{O}_{D \cap D'}
\cong
\Coker(\psi)
$$
かつ
$$
\Ker(\varphi)
\cong
\frac{\mathcal{I} \cap \mathcal{I}'}{\mathcal{I}\mathcal{I}'}
\cong
\Ker(\psi).
$$
である．したがって $K_0(\textit{Coh}_{\leq n - 1}(X))$ において
$$
\gamma =
[\mathcal{I}_{D'}] - [\mathcal{O}_{D'}]
=
[\mathcal{I}'_D] - [\mathcal{O}_D]
$$
である．一方，明らかに
$$
[\mathcal{I}'_D]_{n - 1} = [D]_{n - 1}, \quad
[\mathcal{I}_{D'}]_{n - 1} = [D']_{n - 1}.
$$
であり，また
$$
\mathcal{O}_X(D') \otimes \mathcal{I}'_D = \mathcal{O}_D, \quad
\mathcal{O}_X(D) \otimes \mathcal{I}_{D'} = \mathcal{O}_{D'}.
$$
である．Chow Homology, Lemma
\ref{chow-lemma-coherent-sheaf-cap-c1} を二度適用すると，これは元 $\gamma$ が
$B_{n - 2}(X)$ の元であり，
$c_1(\mathcal{O}_X(D')) \cap [D]_{n - 1}$ と
$c_1(\mathcal{O}_X(D)) \cap [D']_{n - 1}$ の両方へ写ることを意味する．これで結論を得る
（写像 $B_{n - 2}(X) \to \CH_{n - 2}(X)$ が well-defined であることがこの証明の
要点である）．
\end{proof}
% Frozen obsolete.tex lines 4312--4409.
\section{正則固有モデル上の双対化加群}
\label{obsolete-section-dualizing}

\noindent
Semistable Reduction, Situation \ref{models-situation-regular-model} では，Duality for
Schemes, Remark \ref{duality-remark-relative-dualizing-complex} で導入された
$f : X \to \Spec(R)$ の相対双対化複体を
$\omega_{X/R}^\bullet = f^!\mathcal{O}_{\Spec(R)}$ とした．Semistable Reduction,
Lemma \ref{models-lemma-gorenstein} により $f$ は相対次元 $1$ の Gorenstein 射なので，
Duality for Schemes, Lemmas
\ref{duality-lemma-affine-flat-Noetherian-gorenstein},
\ref{duality-lemma-CM-shriek}, および
\ref{duality-lemma-gorenstein-CM-morphism} を用いると，ある可逆
$\mathcal{O}_X$-加群 $\omega_X$ に対して
$$
\omega_{X/R}^\bullet = \omega_X[1]
$$
であることが分かる．この可逆加群はしばしば {it $R$ 上の $X$ の相対双対化加群}
と呼ばれる．$R$ は次元 $1$ の正則環（したがって Gorenstein 環）なので，
$\omega_R^\bullet = R[1]$ は $R$ の正規化された双対化複体である．したがって
$\omega_X = H^{-2}(f^!\omega_R^\bullet)$ であり，$\omega_X$ は単に相対双対化加群で
あるだけでなく双対化加群でもあることが分かる．Duality for Schemes, Example
\ref{duality-example-proper-over-local} を参照せよ．ゆえに Duality for Schemes, Lemma
\ref{duality-lemma-dualizing-module-proper-over-A} により，$\omega_X$ は函手
$$
\textit{Coh}(\mathcal{O}_X) \to \textit{Sets},\quad
\mathcal{F} \mapsto \Hom_R(H^1(X, \mathcal{F}), R)
$$
を表現する．これは Semistable Reduction, Situation
\ref{models-situation-regular-model} における相対双対化加群の別の定義を与える．
$\omega_X$ の形成は任意の基底変換と可換する（与えられた相対次元をもつ任意の
固有 Gorenstein 射について）．これは Duality for Schemes, Remark
\ref{duality-remark-relative-dualizing-complex} で論じた相対双対化複体に関する対応する
事実から従い，その事実は Duality for Schemes, Lemma
\ref{duality-lemma-proper-flat-base-change} に遡る．したがって $\omega_X$ は，
Algebraic Curves, Lemma \ref{curves-lemma-duality-dim-1-CM} で論じた $K$ 上の $C$ の
双対化加群 $\omega_C$ へ引き戻される．Algebraic Curves, Lemma
\ref{curves-lemma-duality-dim-1} により，$\omega_C$ は $\Omega_{C/K}$ と同型であることに
注意する．同様に，$\omega_X|_{X_k}$ は $k$ 上の $X_k$ の双対化加群
$\omega_{X_k}$ である．

\begin{lemma}
\label{obsolete-lemma-dualizing-components}
Semistable Reduction, Situation \ref{models-situation-regular-model} において，
$k$ 上の $C_i$ の双対化加群は
$$
\omega_{C_i} = \omega_X(C_i)|_{C_i}
$$
である．ここで $\omega_X$ は上のとおりである．
\end{lemma}

\begin{proof}
$t : C_i \to X$ を閉埋入とする．$t$ は有効 Cartier 因子の包含なので，Duality for
Schemes, Lemmas \ref{duality-lemma-twisted-inverse-image-closed} および
\ref{duality-lemma-sheaf-with-exact-support-effective-Cartier} から，任意の可逆
$\mathcal{O}_X$-加群 $\mathcal{L}$ に対して
$t^!(\mathcal{L}) = \mathcal{L}(C_i)|_{C_i}$ であると結論できる．可換図式
$$
\xymatrix{
C_i \ar[r]_t \ar[d]_g & X \ar[d]^f \\
\Spec(k) \ar[r]^s & \Spec(R)
}
$$
を考える．$C_i$ は Gorenstein 曲線であり（Semistable Reduction, Lemma
\ref{models-lemma-gorenstein}），その可逆双対化加群 $\omega_{C_i}$ は性質
$\omega_{C_i}[0] = g^!\mathcal{O}_{\Spec(k)}$ によって特徴づけられることに注意する．
Algebraic Curves, Lemma \ref{curves-lemma-duality-dim-1} とその証明，および Algebraic
Curves, Lemmas \ref{curves-lemma-duality-dim-1-CM} と
\ref{curves-lemma-rr} を参照せよ．一方，$s^!(R[1]) = k$ なので
$$
\omega_{C_i}[0] =
g^! s^!(R[1]) = t^!f^!(R[1]) = t^!\omega_X
$$
である．以上を組み合わせると補題の主張を得る．
\end{proof}
% Frozen obsolete.tex lines 4410--4544.
\section{重複して分割された参照}
\label{obsolete-section-duplicates}

\noindent
この節には，重複した結果，および以前は複数の内容を同時に述べていたが現在では
構成部分に分割された参照を集める．

\begin{lemma}
\label{obsolete-lemma-characterize-epi-mono}
$X$ を位相空間とする．
\begin{enumerate}
\item 前層の圏におけるエピ射（それぞれモノ射）は，ちょうど前層の全射
（それぞれ単射）である．
\item 層の圏におけるエピ射（それぞれモノ射）は，ちょうど層の全射
（それぞれ単射）であり，またすべての茎上で全射（それぞれ単射）である写像に
ちょうど一致する．
\item 集合の前層の全射（それぞれ単射）の層化は全射（それぞれ単射）である．
\end{enumerate}
\end{lemma}

\begin{proof}
Sheaves, Lemmas \ref{sheaves-lemma-characterize-epi-mono} および
\ref{sheaves-lemma-sheafify-epi-mono} を組み合わせればよい．
\end{proof}

\begin{lemma}
\label{obsolete-lemma-directed-colimit-finite-type}
$X$ をスキームとし，$X$ は準コンパクトかつ準分離的であると仮定する．
$\mathcal{F}$ を準連接 $\mathcal{O}_X$-加群とする．このとき $\mathcal{F}$ は，
その有限型準連接部分加群たちの有向余極限である．
\end{lemma}

\begin{proof}
これは Properties, Lemma
\ref{properties-lemma-quasi-coherent-colimit-finite-type} の重複である．
\end{proof}

\begin{lemma}
\label{obsolete-lemma-points-monomorphism}
$S$ をスキームとする．$X$ を $S$ 上の代数空間とする．写像
$\{\Spec(k) \to X \text{ monomorphism}\} \to |X|$ は単射である．
\end{lemma}

\begin{proof}
これは Properties of Spaces, Lemma
\ref{spaces-properties-lemma-points-monomorphism} の重複である．
\end{proof}

\begin{theorem}
\label{obsolete-theorem-equivalence-sheaves-point}
$K$ を体として $S = \Spec(K)$ とする．$\overline{s}$ を $S$ の幾何学的点とする．
$G = \text{Gal}_{\kappa(s)}$ を絶対 Galois 群とする．このとき圏同値
$\Sh(S_\etale) \to G\textit{-Sets}$,
$\mathcal{F} \mapsto \mathcal{F}_{\overline{s}}$ がある．
\end{theorem}

\begin{proof}
これは \'Etale Cohomology, Theorem
\ref{etale-cohomology-theorem-equivalence-sheaves-point} の重複である．
\end{proof}

\begin{remark}
\label{obsolete-remark-tangent-spaces}
重複したタグをたどってここに到達したはずである．実際の内容については
Formal Deformation Theory, Section
\ref{formal-defos-section-tangent-spaces} を参照されたい．
\end{remark}

\begin{lemma}
\label{obsolete-lemma-locally-ringed-space-direct-summand-free}
$X$ を局所環付き空間とする．有限自由 $\mathcal{O}_X$-加群の直和因子は有限局所自由である．
\end{lemma}

\begin{proof}
これは Modules, Lemma
\ref{modules-lemma-direct-summand-of-locally-free-is-locally-free} の重複である．
\end{proof}

\begin{lemma}
\label{obsolete-lemma-characterize-injective}
$R$ を環とし，$E$ を $R$-加群とする．次は同値である：
\begin{enumerate}
\item $E$ は入射的 $R$-加群である；
\item イデアル $I \subset R$ と加群準同型 $\varphi : I \to E$ が与えられたとき，
$\varphi$ を延長する $R$-加群準同型 $R \to E$ が存在する．
\end{enumerate}
\end{lemma}

\begin{proof}
これは Baer の判定法である．Injectives, Lemma
\ref{injectives-lemma-criterion-baer} を参照せよ．
\end{proof}

\begin{lemma}
\label{obsolete-lemma-periodic-length}
$R$ を局所環とする．
\begin{enumerate}
\item $(M, N, \varphi, \psi)$ が $2$-周期複体で，$M$, $N$ の長さが有限ならば，
$e_R(M, N, \varphi, \psi) = \text{length}_R(M) - \text{length}_R(N)$ である．
\item $(M, \varphi, \psi)$ が $(2, 1)$-周期複体で，$M$ の長さが有限ならば，
$e_R(M, \varphi, \psi) = 0$ である．
\item $2$-周期複体の短完全列
$$
0 \to (M_1, N_1, \varphi_1, \psi_1)
\to (M_2, N_2, \varphi_2, \psi_2)
\to (M_3, N_3, \varphi_3, \psi_3)
\to 0
$$
があるとする．三つのうち二つが有限長のコホモロジー加群をもつならば，三つ目も
そうであり，
$$
e_R(M_2, N_2, \varphi_2, \psi_2) =
e_R(M_1, N_1, \varphi_1, \psi_1) +
e_R(M_3, N_3, \varphi_3, \psi_3).
$$
が成り立つ．
\end{enumerate}
\end{lemma}

\begin{proof}
Chow Homology, Lemmas
\ref{chow-lemma-additivity-periodic-length} および
\ref{chow-lemma-finite-periodic-length} から従う．
\end{proof}
% Frozen obsolete.tex lines 4545--4650.

\begin{lemma}
\label{obsolete-lemma-extensions-of-rings}
$A$ を環とし，$I$ を $A$-加群とする．
\begin{enumerate}
\item $I$ が平方零イデアルであるような環の拡大
$0 \to I \to A' \to A \to 0$ の集合は，
$\Ext^1_A(\NL_{A/\mathbf{Z}}, I)$ と標準的に全単射である．
\item 環準同型 $A \to B$，$B$-加群 $N$，$A$-加群準同型 $c : I \to N$，および
平方零の核をもつ環の拡大が与えられているとする：
\begin{enumerate}
\item[(a)] $\alpha \in \Ext^1_A(\NL_{A/\mathbf{Z}}, I)$ に対応する
$0 \to I \to A' \to A \to 0$；
\item[(b)] $\beta \in \Ext^1_B(\NL_{B/\mathbf{Z}}, N)$ に対応する
$0 \to N \to B' \to B \to 0$．
\end{enumerate}
このとき，Deformation Theory, Equation (\ref{defos-equation-to-solve}) に組み込まれる
写像 $A' \to B'$ が存在するための必要十分条件は，$\beta$ と $\alpha$ が
$\Ext^1_A(\NL_{A/\mathbf{Z}}, N)$ の同じ元へ写ることである．
\end{enumerate}
\end{lemma}

\begin{proof}
Deformation Theory, Lemmas
\ref{defos-lemma-extensions-of-algebras} および
\ref{defos-lemma-extensions-of-algebras-functorial} から従う．
\end{proof}

\begin{lemma}
\label{obsolete-lemma-extensions-of-ringed-spaces}
$(S, \mathcal{O}_S)$ を環付き空間とし，$\mathcal{J}$ を
$\mathcal{O}_S$-加群とする．
\begin{enumerate}
\item $\mathcal{J}$ が平方零イデアルであるような環の層の拡大
$0 \to \mathcal{J} \to \mathcal{O}_{S'} \to \mathcal{O}_S \to 0$ の集合は，
$\Ext^1_{\mathcal{O}_S}(\NL_{S/\mathbf{Z}}, \mathcal{J})$ と標準的に全単射である．
\item 環付き空間の射
$f : (X, \mathcal{O}_X) \to (S, \mathcal{O}_S)$，$\mathcal{O}_X$-加群
$\mathcal{G}$，$f$-写像 $c : \mathcal{J} \to \mathcal{G}$，および平方零の核をもつ
環の層の拡大が与えられているとする：
\begin{enumerate}
\item[(a)]
$\alpha \in \Ext^1_{\mathcal{O}_S}(\NL_{S/\mathbf{Z}}, \mathcal{J})$ に対応する
$0 \to \mathcal{J} \to \mathcal{O}_{S'} \to \mathcal{O}_S \to 0$；
\item[(b)]
$\beta \in \Ext^1_{\mathcal{O}_X}(\NL_{X/\mathbf{Z}}, \mathcal{G})$ に対応する
$0 \to \mathcal{G} \to \mathcal{O}_{X'} \to \mathcal{O}_X \to 0$．
\end{enumerate}
このとき，Deformation Theory, Equation
(\ref{defos-equation-to-solve-ringed-spaces}) に組み込まれる射 $X' \to S'$ が
存在するための必要十分条件は，$\beta$ と $\alpha$ が
$\Ext^1_{\mathcal{O}_X}(Lf^*\NL_{S/\mathbf{Z}}, \mathcal{G})$ の同じ元へ写ることである．
\end{enumerate}
\end{lemma}

\begin{proof}
Deformation Theory, Lemmas
\ref{defos-lemma-extensions-of-relative-ringed-spaces} および
\ref{defos-lemma-extensions-of-relative-ringed-spaces-functorial} から従う．
\end{proof}

\begin{lemma}
\label{obsolete-lemma-extensions-of-ringed-topoi}
$(\Sh(\mathcal{B}), \mathcal{O}_\mathcal{B})$ を環付きトポスとし，
$\mathcal{J}$ を $\mathcal{O}_\mathcal{B}$-加群とする．
\begin{enumerate}
\item $\mathcal{J}$ が平方零イデアルであるような環の層の拡大
$0 \to \mathcal{J} \to \mathcal{O}_{\mathcal{B}'} \to
\mathcal{O}_\mathcal{B} \to 0$ の集合は，
$\Ext^1_{\mathcal{O}_\mathcal{B}}(
\NL_{\mathcal{O}_\mathcal{B}/\mathbf{Z}}, \mathcal{J})$ と標準的に全単射である．
\item 環付きトポスの射
$f : (\Sh(\mathcal{C}), \mathcal{O}) \to
(\Sh(\mathcal{B}), \mathcal{O}_\mathcal{B})$，$\mathcal{O}$-加群
$\mathcal{G}$，$f^{-1}\mathcal{O}_\mathcal{B}$-加群準同型
$c : f^{-1}\mathcal{J} \to \mathcal{G}$，および平方零の核をもつ環の層の拡大が
与えられているとする：
\begin{enumerate}
\item[(a)]
$\alpha \in \Ext^1_{\mathcal{O}_\mathcal{B}}(
\NL_{\mathcal{O}_\mathcal{B}/\mathbf{Z}}, \mathcal{J})$ に対応する
$0 \to \mathcal{J} \to \mathcal{O}_{\mathcal{B}'} \to
\mathcal{O}_\mathcal{B} \to 0$；
\item[(b)] $\beta \in \Ext^1_\mathcal{O}(\NL_{\mathcal{O}/\mathbf{Z}}, \mathcal{G})$
に対応する $0 \to \mathcal{G} \to \mathcal{O}' \to \mathcal{O} \to 0$．
\end{enumerate}
このとき，Deformation Theory, Equation
(\ref{defos-equation-to-solve-ringed-topoi}) に組み込まれる射
$(\Sh(\mathcal{C}), \mathcal{O}') \to
(\Sh(\mathcal{B}, \mathcal{O}_{\mathcal{B}'})$ が存在するための必要十分条件は，
$\beta$ と $\alpha$ が
$\Ext^1_\mathcal{O}(
Lf^*\NL_{\mathcal{O}_\mathcal{B}/\mathbf{Z}}, \mathcal{G})$ の同じ元へ写ることである．
\end{enumerate}
\end{lemma}

\begin{proof}
Deformation Theory, Lemmas
\ref{defos-lemma-extensions-of-relative-ringed-topoi} および
\ref{defos-lemma-extensions-of-relative-ringed-topoi-functorial} から従う．
\end{proof}
% Frozen obsolete.tex lines 4651--4729.
\begin{remark}
\label{obsolete-remark-examples-formal-defos}
このタグは以前，変形問題のいくつかの例を記述する節を指していた．現在では各例が
それぞれ独自の節をもつ．Deformation Problems, Sections
\ref{examples-defos-section-finite-projective-modules},
\ref{examples-defos-section-representations},
\ref{examples-defos-section-continuous-representations}, および
\ref{examples-defos-section-graded-algebras} を参照せよ．
\end{remark}

\begin{lemma}
\label{obsolete-lemma-examples-have-RS}
Deformation Problems, Examples
\ref{examples-defos-example-finite-projective-modules},
\ref{examples-defos-example-representations},
\ref{examples-defos-example-continuous-representations}, および
\ref{examples-defos-example-graded-algebras} は Rim--Schlessinger 条件 (RS) を満たす．
\end{lemma}

\begin{proof}
Deformation Problems, Lemmas
\ref{examples-defos-lemma-finite-projective-modules-RS},
\ref{examples-defos-lemma-representations-RS},
\ref{examples-defos-lemma-continuous-representations-RS}, および
\ref{examples-defos-lemma-graded-algebras-RS} から従う．
\end{proof}

\begin{lemma}
\label{obsolete-lemma-tangent-and-inf}
次の標準的な $k$-ベクトル空間の同定がある：
\begin{enumerate}
\item Deformation Problems, Example
\ref{examples-defos-example-finite-projective-modules} において，
$x_0 = (k, V)$ ならば，
$T_{x_0}\mathcal{F} = (0)$ および
$\text{Inf}_{x_0}(\mathcal{F}) = \text{End}_k(V)$ は有限次元である．
\item Deformation Problems, Example
\ref{examples-defos-example-representations} において，
$x_0 = (k, V, \rho_0)$ ならば，
$T_{x_0}\mathcal{F} = \Ext^1_{k[\Gamma]}(V, V) = H^1(\Gamma, \text{End}_k(V))$
および
$\text{Inf}_{x_0}(\mathcal{F}) = H^0(\Gamma, \text{End}_k(V))$
は，$\Gamma$ が有限生成ならば有限次元である．
\item Deformation Problems, Example
\ref{examples-defos-example-continuous-representations} において，
$x_0 = (k, V, \rho_0)$ ならば，
$T_{x_0}\mathcal{F} = H^1_{cont}(\Gamma, \text{End}_k(V))$
および
$\text{Inf}_{x_0}(\mathcal{F}) = H^0_{cont}(\Gamma, \text{End}_k(V))$
は，$\Gamma$ が位相的有限生成ならば有限次元である．
\item Deformation Problems, Example
\ref{examples-defos-example-graded-algebras} において，
$x_0 = (k, P)$ ならば，
$T_{x_0}\mathcal{F}$ および
$\text{Inf}_{x_0}(\mathcal{F}) = \text{Der}_k(P, P)$ は，$P$ が $k$ 上有限生成ならば
有限次元である．
\end{enumerate}
\end{lemma}

\begin{proof}
Deformation Problems, Lemmas
\ref{examples-defos-lemma-finite-projective-modules-TI},
\ref{examples-defos-lemma-representations-TI},
\ref{examples-defos-lemma-continuous-representations-TI}, および
\ref{examples-defos-lemma-graded-algebras-TI} から従う．
\end{proof}
